
% !TEX root = ../../../../Esperimenti/.tex/MEF.tex
% LTeX: language=it

\begin{tikzpicture}%
    %\begingroup Impostazioni
        \infopageKeysI
        \pgfplotsset{
            /pgfplots/group/plotGroup/.append style={
                group size=2 by 1,
            },
            plotRow/.append style={
                width=\plotLastWidth cm
            },
        }
    %\endgroup

    %\begingroup Righe
        \begin{groupplot}[
            group style={
                group name=Row1,
                plotGroup,
            },
            plotRow,
            row1Keys,
            plotLegendNE,
            %
            xticklabels={},
            % xmin=1.7158,xmax=5.4719,
            ymin=0,ymax=1.4,
        ]
            \nextgroupplot[%
                xmin=1.7064,xmax=5.4903,
                % xmin=1.9138,xmax=5.4719,
                % ymin=0,ymax=1.4,
            ] % This file was created by matlab2tikz.
%
\definecolor{mycolor1}{rgb}{0.00000,0.44706,0.69804}%
%
\addplot[ybar interval, fill=mycolor1, fill opacity=0.35, area legend, draw=none] table[row sep=crcr, x=Lower, y=Count] {%
Lower	Upper	Count\\
2.0755	2.3244	0.001822\\
2.3244	2.5732	0.020363\\
2.5732	2.822	0.15208\\
2.822	3.0708	0.56545\\
3.0708	3.3196	1.1158\\
3.3196	3.5684	1.0444\\
3.5684	3.8172	0.54091\\
3.8172	4.0661	0.26687\\
4.0661	4.3149	0.16966\\
4.3149	4.5637	0.10053\\
4.5637	4.8125	0.035904\\
4.8125	5.0613	0.00493\\
5.0613	5.3101	0.00032152\\
5.3101	5.3101	0.00032152\\
};
\addlegendentry{Istogramma esatto}

\addplot [color=mycolor1, only marks, every error bar/.append style={opacity=0.45}, mark=*, mark size=0pt, draw=none, forget plot]
 plot [error bars/.cd, y dir=both, y explicit, error bar style={line width=1pt, color=mycolor1}, error mark options={mark=none,mark size=0pt}]
 table[row sep=crcr, y error plus index=2, y error minus index=3]{%
2.2	0.001822	0.00085785	0.00085785\\
2.4488	0.020363	0.0044061	0.0044061\\
2.6976	0.15208	0.01936	0.01936\\
2.9464	0.56545	0.041558	0.041558\\
3.1952	1.1158	0.030793	0.030793\\
3.444	1.0444	0.03402	0.03402\\
3.6928	0.54091	0.032723	0.032723\\
3.9417	0.26687	0.01457	0.01457\\
4.1905	0.16966	0.007587	0.007587\\
4.4393	0.10053	0.0068977	0.0068977\\
4.6881	0.035904	0.0049146	0.0049146\\
4.9369	0.00493	0.0015536	0.0015536\\
5.1857	0.00032152	0.0003646	0.00032152\\
};
\addplot [color=mycolor1]
  table[row sep=crcr]{%
2.0755	0.00078348\\
2.1953	0.0020183\\
2.2731	0.0038435\\
2.3281	0.0061442\\
2.3734	0.0090882\\
2.409	0.012377\\
2.4414	0.016382\\
2.4706	0.021054\\
2.4965	0.026266\\
2.5191	0.031812\\
2.5418	0.038444\\
2.5612	0.045125\\
2.5806	0.052856\\
2.6001	0.061766\\
2.6195	0.071995\\
2.6357	0.081631\\
2.6519	0.092371\\
2.6681	0.1043\\
2.6843	0.11752\\
2.7004	0.13211\\
2.7166	0.14816\\
2.7328	0.16575\\
2.749	0.18495\\
2.7652	0.20585\\
2.7814	0.2285\\
2.7976	0.25294\\
2.8138	0.27922\\
2.83	0.30735\\
2.8461	0.33733\\
2.8656	0.37572\\
2.885	0.41665\\
2.9044	0.46\\
2.9271	0.51338\\
2.953	0.57758\\
2.9821	0.65293\\
3.0339	0.79078\\
3.0696	0.88422\\
3.0955	0.94916\\
3.1149	0.99522\\
3.1343	1.0384\\
3.1505	1.0717\\
3.1667	1.1022\\
3.1796	1.1245\\
3.1926	1.1447\\
3.2055	1.1627\\
3.2153	1.1747\\
3.225	1.1854\\
3.2347	1.1947\\
3.2444	1.2026\\
3.2541	1.2091\\
3.2638	1.2141\\
3.2735	1.2177\\
3.2833	1.2198\\
3.293	1.2204\\
3.3027	1.2196\\
3.3124	1.2174\\
3.3221	1.2137\\
3.3318	1.2087\\
3.3415	1.2023\\
3.3513	1.1945\\
3.361	1.1854\\
3.3707	1.1751\\
3.3836	1.1595\\
3.3966	1.1418\\
3.4095	1.1222\\
3.4225	1.101\\
3.4387	1.0721\\
3.4549	1.0412\\
3.4743	1.0016\\
3.497	0.95281\\
3.5261	0.88739\\
3.6103	0.69652\\
3.6362	0.64117\\
3.6588	0.59516\\
3.6783	0.55782\\
3.6977	0.52257\\
3.7171	0.48955\\
3.7366	0.45881\\
3.756	0.43039\\
3.7722	0.40845\\
3.7884	0.38807\\
3.8045	0.36922\\
3.8207	0.35182\\
3.8369	0.33582\\
3.8531	0.32115\\
3.8693	0.30771\\
3.8887	0.29311\\
3.9082	0.28001\\
3.9276	0.26826\\
3.947	0.2577\\
3.9697	0.2467\\
3.9923	0.23688\\
4.0182	0.22684\\
4.0474	0.21669\\
4.083	0.20541\\
4.1413	0.1883\\
4.2287	0.16258\\
4.287	0.14425\\
4.3582	0.12058\\
4.4586	0.087205\\
4.5072	0.072271\\
4.5492	0.060459\\
4.5881	0.050658\\
4.627	0.041988\\
4.6658	0.034454\\
4.7047	0.028012\\
4.7468	0.022171\\
4.7888	0.017388\\
4.8374	0.013\\
4.8892	0.0094253\\
4.9475	0.0064806\\
5.0155	0.0041198\\
5.0997	0.0022983\\
5.2098	0.0010322\\
5.3101	0.00047951\\
};
\addlegendentry{Adatt. BLN esatto}


\addplot[area legend, draw=none, fill=mycolor1, fill opacity=0.15, forget plot]
table[row sep=crcr] {%
x	y\\
2.0755	0.0010457\\
2.0788	0.0010706\\
2.082	0.0010961\\
2.0853	0.0011223\\
2.0885	0.0011492\\
2.0917	0.0011767\\
2.095	0.0012049\\
2.0982	0.0012338\\
2.1014	0.0012634\\
2.1047	0.0012938\\
2.1079	0.001325\\
2.1112	0.0013569\\
2.1144	0.0013897\\
2.1176	0.0014233\\
2.1209	0.0014578\\
2.1241	0.0014932\\
2.1274	0.0015295\\
2.1306	0.0015667\\
2.1338	0.0016049\\
2.1371	0.001644\\
2.1403	0.0016842\\
2.1435	0.0017255\\
2.1468	0.0017678\\
2.15	0.0018113\\
2.1533	0.0018559\\
2.1565	0.0019017\\
2.1597	0.0019487\\
2.163	0.0019969\\
2.1662	0.0020465\\
2.1694	0.0020973\\
2.1727	0.0021495\\
2.1759	0.0022032\\
2.1792	0.0022582\\
2.1824	0.0023148\\
2.1856	0.0023729\\
2.1889	0.0024326\\
2.1921	0.0024939\\
2.1953	0.0025568\\
2.1986	0.0026215\\
2.2018	0.002688\\
2.2051	0.0027563\\
2.2083	0.0028265\\
2.2115	0.0028986\\
2.2148	0.0029727\\
2.218	0.0030489\\
2.2212	0.0031272\\
2.2245	0.0032077\\
2.2277	0.0032904\\
2.231	0.0033754\\
2.2342	0.0034628\\
2.2374	0.0035527\\
2.2407	0.0036451\\
2.2439	0.0037401\\
2.2472	0.0038378\\
2.2504	0.0039382\\
2.2536	0.0040415\\
2.2569	0.0041477\\
2.2601	0.0042569\\
2.2633	0.0043693\\
2.2666	0.0044848\\
2.2698	0.0046036\\
2.2731	0.0047259\\
2.2763	0.0048516\\
2.2795	0.0049809\\
2.2828	0.0051139\\
2.286	0.0052507\\
2.2892	0.0053914\\
2.2925	0.0055362\\
2.2957	0.0056852\\
2.299	0.0058384\\
2.3022	0.005996\\
2.3054	0.0061582\\
2.3087	0.006325\\
2.3119	0.0064966\\
2.3151	0.0066732\\
2.3184	0.0068548\\
2.3216	0.0070417\\
2.3249	0.0072339\\
2.3281	0.0074317\\
2.3313	0.0076351\\
2.3346	0.0078445\\
2.3378	0.0080598\\
2.341	0.0082813\\
2.3443	0.0085092\\
2.3475	0.0087437\\
2.3508	0.0089848\\
2.354	0.0092329\\
2.3572	0.0094881\\
2.3605	0.0097506\\
2.3637	0.010021\\
2.367	0.010298\\
2.3702	0.010584\\
2.3734	0.010878\\
2.3767	0.01118\\
2.3799	0.011491\\
2.3831	0.01181\\
2.3864	0.012139\\
2.3896	0.012477\\
2.3929	0.012824\\
2.3961	0.013182\\
2.3993	0.013549\\
2.4026	0.013927\\
2.4058	0.014315\\
2.409	0.014714\\
2.4123	0.015125\\
2.4155	0.015546\\
2.4188	0.01598\\
2.422	0.016426\\
2.4252	0.016884\\
2.4285	0.017354\\
2.4317	0.017838\\
2.4349	0.018335\\
2.4382	0.018846\\
2.4414	0.01937\\
2.4447	0.019909\\
2.4479	0.020463\\
2.4511	0.021032\\
2.4544	0.021616\\
2.4576	0.022216\\
2.4608	0.022832\\
2.4641	0.023465\\
2.4673	0.024114\\
2.4706	0.024781\\
2.4738	0.025466\\
2.477	0.026169\\
2.4803	0.026891\\
2.4835	0.027632\\
2.4868	0.028392\\
2.49	0.029172\\
2.4932	0.029973\\
2.4965	0.030794\\
2.4997	0.031637\\
2.5029	0.032502\\
2.5062	0.033389\\
2.5094	0.034298\\
2.5127	0.035231\\
2.5159	0.036188\\
2.5191	0.037169\\
2.5224	0.038175\\
2.5256	0.039206\\
2.5288	0.040263\\
2.5321	0.041347\\
2.5353	0.042457\\
2.5386	0.043595\\
2.5418	0.044761\\
2.545	0.045956\\
2.5483	0.04718\\
2.5515	0.048434\\
2.5547	0.049719\\
2.558	0.051034\\
2.5612	0.052381\\
2.5645	0.05376\\
2.5677	0.055173\\
2.5709	0.056618\\
2.5742	0.058098\\
2.5774	0.059613\\
2.5806	0.061163\\
2.5839	0.062749\\
2.5871	0.064372\\
2.5904	0.066032\\
2.5936	0.06773\\
2.5968	0.069467\\
2.6001	0.071243\\
2.6033	0.07306\\
2.6066	0.074917\\
2.6098	0.076816\\
2.613	0.078756\\
2.6163	0.08074\\
2.6195	0.082767\\
2.6227	0.084838\\
2.626	0.086955\\
2.6292	0.089117\\
2.6325	0.091325\\
2.6357	0.09358\\
2.6389	0.095883\\
2.6422	0.098235\\
2.6454	0.10064\\
2.6486	0.10309\\
2.6519	0.10559\\
2.6551	0.10814\\
2.6584	0.11075\\
2.6616	0.1134\\
2.6648	0.11611\\
2.6681	0.11888\\
2.6713	0.1217\\
2.6745	0.12457\\
2.6778	0.12751\\
2.681	0.13049\\
2.6843	0.13354\\
2.6875	0.13665\\
2.6907	0.13981\\
2.694	0.14303\\
2.6972	0.14632\\
2.7004	0.14966\\
2.7037	0.15307\\
2.7069	0.15654\\
2.7102	0.16007\\
2.7134	0.16366\\
2.7166	0.16732\\
2.7199	0.17105\\
2.7231	0.17484\\
2.7263	0.17869\\
2.7296	0.18261\\
2.7328	0.1866\\
2.7361	0.19066\\
2.7393	0.19478\\
2.7425	0.19898\\
2.7458	0.20324\\
2.749	0.20757\\
2.7523	0.21197\\
2.7555	0.21644\\
2.7587	0.22098\\
2.762	0.2256\\
2.7652	0.23028\\
2.7684	0.23504\\
2.7717	0.23987\\
2.7749	0.24477\\
2.7782	0.24975\\
2.7814	0.2548\\
2.7846	0.25992\\
2.7879	0.26511\\
2.7911	0.27038\\
2.7943	0.27572\\
2.7976	0.28114\\
2.8008	0.28663\\
2.8041	0.29219\\
2.8073	0.29783\\
2.8105	0.30354\\
2.8138	0.30933\\
2.817	0.31519\\
2.8202	0.32113\\
2.8235	0.32714\\
2.8267	0.33322\\
2.83	0.33938\\
2.8332	0.34561\\
2.8364	0.35191\\
2.8397	0.35828\\
2.8429	0.36473\\
2.8461	0.37125\\
2.8494	0.37784\\
2.8526	0.3845\\
2.8559	0.39124\\
2.8591	0.39804\\
2.8623	0.40491\\
2.8656	0.41186\\
2.8688	0.41887\\
2.8721	0.42594\\
2.8753	0.43309\\
2.8785	0.4403\\
2.8818	0.44757\\
2.885	0.45491\\
2.8882	0.46232\\
2.8915	0.46978\\
2.8947	0.47731\\
2.898	0.4849\\
2.9012	0.49254\\
2.9044	0.50025\\
2.9077	0.50801\\
2.9109	0.51583\\
2.9141	0.5237\\
2.9174	0.53162\\
2.9206	0.5396\\
2.9239	0.54763\\
2.9271	0.55571\\
2.9303	0.56383\\
2.9336	0.572\\
2.9368	0.58021\\
2.94	0.58847\\
2.9433	0.59677\\
2.9465	0.6051\\
2.9498	0.61348\\
2.953	0.62189\\
2.9562	0.63033\\
2.9595	0.63881\\
2.9627	0.64731\\
2.9659	0.65584\\
2.9692	0.6644\\
2.9724	0.67299\\
2.9757	0.68159\\
2.9789	0.69022\\
2.9821	0.69886\\
2.9854	0.70752\\
2.9886	0.71619\\
2.9919	0.72487\\
2.9951	0.73356\\
2.9983	0.74225\\
3.0016	0.75095\\
3.0048	0.75965\\
3.008	0.76835\\
3.0113	0.77705\\
3.0145	0.78574\\
3.0178	0.79442\\
3.021	0.80309\\
3.0242	0.81174\\
3.0275	0.82038\\
3.0307	0.829\\
3.0339	0.8376\\
3.0372	0.84618\\
3.0404	0.85472\\
3.0437	0.86324\\
3.0469	0.87173\\
3.0501	0.88018\\
3.0534	0.88859\\
3.0566	0.89696\\
3.0598	0.90529\\
3.0631	0.91357\\
3.0663	0.9218\\
3.0696	0.92998\\
3.0728	0.93811\\
3.076	0.94618\\
3.0793	0.95419\\
3.0825	0.96214\\
3.0857	0.97002\\
3.089	0.97783\\
3.0922	0.98557\\
3.0955	0.99324\\
3.0987	1.0008\\
3.1019	1.0083\\
3.1052	1.0158\\
3.1084	1.0231\\
3.1117	1.0304\\
3.1149	1.0376\\
3.1181	1.0446\\
3.1214	1.0516\\
3.1246	1.0585\\
3.1278	1.0653\\
3.1311	1.072\\
3.1343	1.0786\\
3.1376	1.085\\
3.1408	1.0914\\
3.144	1.0977\\
3.1473	1.1038\\
3.1505	1.1098\\
3.1537	1.1157\\
3.157	1.1215\\
3.1602	1.1272\\
3.1635	1.1327\\
3.1667	1.1382\\
3.1699	1.1434\\
3.1732	1.1486\\
3.1764	1.1536\\
3.1796	1.1585\\
3.1829	1.1633\\
3.1861	1.1679\\
3.1894	1.1724\\
3.1926	1.1768\\
3.1958	1.181\\
3.1991	1.1851\\
3.2023	1.189\\
3.2055	1.1928\\
3.2088	1.1964\\
3.212	1.1999\\
3.2153	1.2033\\
3.2185	1.2065\\
3.2217	1.2096\\
3.225	1.2125\\
3.2282	1.2153\\
3.2315	1.2179\\
3.2347	1.2204\\
3.2379	1.2228\\
3.2412	1.225\\
3.2444	1.227\\
3.2476	1.2289\\
3.2509	1.2307\\
3.2541	1.2323\\
3.2574	1.2338\\
3.2606	1.2351\\
3.2638	1.2363\\
3.2671	1.2374\\
3.2703	1.2383\\
3.2735	1.2391\\
3.2768	1.2398\\
3.28	1.2403\\
3.2833	1.2406\\
3.2865	1.2409\\
3.2897	1.241\\
3.293	1.241\\
3.2962	1.2408\\
3.2994	1.2405\\
3.3027	1.2401\\
3.3059	1.2396\\
3.3092	1.2389\\
3.3124	1.2381\\
3.3156	1.2371\\
3.3189	1.2361\\
3.3221	1.2349\\
3.3253	1.2336\\
3.3286	1.2321\\
3.3318	1.2305\\
3.3351	1.2288\\
3.3383	1.2269\\
3.3415	1.2249\\
3.3448	1.2228\\
3.348	1.2206\\
3.3513	1.2182\\
3.3545	1.2157\\
3.3577	1.213\\
3.361	1.2102\\
3.3642	1.2073\\
3.3674	1.2043\\
3.3707	1.2011\\
3.3739	1.1978\\
3.3772	1.1944\\
3.3804	1.1909\\
3.3836	1.1872\\
3.3869	1.1834\\
3.3901	1.1795\\
3.3933	1.1754\\
3.3966	1.1712\\
3.3998	1.1669\\
3.4031	1.1625\\
3.4063	1.158\\
3.4095	1.1534\\
3.4128	1.1486\\
3.416	1.1438\\
3.4192	1.1388\\
3.4225	1.1337\\
3.4257	1.1285\\
3.429	1.1232\\
3.4322	1.1178\\
3.4354	1.1124\\
3.4387	1.1068\\
3.4419	1.1011\\
3.4451	1.0953\\
3.4484	1.0894\\
3.4516	1.0835\\
3.4549	1.0775\\
3.4581	1.0713\\
3.4613	1.0651\\
3.4646	1.0589\\
3.4678	1.0525\\
3.471	1.0461\\
3.4743	1.0396\\
3.4775	1.033\\
3.4808	1.0264\\
3.484	1.0197\\
3.4872	1.0129\\
3.4905	1.0061\\
3.4937	0.99924\\
3.497	0.99232\\
3.5002	0.98535\\
3.5034	0.97832\\
3.5067	0.97126\\
3.5099	0.96415\\
3.5131	0.957\\
3.5164	0.94981\\
3.5196	0.94259\\
3.5229	0.93534\\
3.5261	0.92805\\
3.5293	0.92074\\
3.5326	0.9134\\
3.5358	0.90604\\
3.539	0.89866\\
3.5423	0.89126\\
3.5455	0.88385\\
3.5488	0.87642\\
3.552	0.86899\\
3.5552	0.86154\\
3.5585	0.85409\\
3.5617	0.84664\\
3.5649	0.83918\\
3.5682	0.83172\\
3.5714	0.82427\\
3.5747	0.81682\\
3.5779	0.80938\\
3.5811	0.80195\\
3.5844	0.79453\\
3.5876	0.78713\\
3.5908	0.77973\\
3.5941	0.77236\\
3.5973	0.76501\\
3.6006	0.75767\\
3.6038	0.75036\\
3.607	0.74308\\
3.6103	0.73582\\
3.6135	0.72859\\
3.6168	0.72139\\
3.62	0.71422\\
3.6232	0.70709\\
3.6265	0.69999\\
3.6297	0.69293\\
3.6329	0.6859\\
3.6362	0.67891\\
3.6394	0.67197\\
3.6427	0.66507\\
3.6459	0.65821\\
3.6491	0.65139\\
3.6524	0.64462\\
3.6556	0.6379\\
3.6588	0.63123\\
3.6621	0.6246\\
3.6653	0.61803\\
3.6686	0.6115\\
3.6718	0.60503\\
3.675	0.59862\\
3.6783	0.59226\\
3.6815	0.58595\\
3.6847	0.5797\\
3.688	0.5735\\
3.6912	0.56737\\
3.6945	0.56129\\
3.6977	0.55527\\
3.7009	0.54931\\
3.7042	0.54341\\
3.7074	0.53757\\
3.7106	0.53179\\
3.7139	0.52607\\
3.7171	0.52042\\
3.7204	0.51483\\
3.7236	0.5093\\
3.7268	0.50383\\
3.7301	0.49843\\
3.7333	0.49309\\
3.7366	0.48781\\
3.7398	0.4826\\
3.743	0.47745\\
3.7463	0.47236\\
3.7495	0.46734\\
3.7527	0.46239\\
3.756	0.4575\\
3.7592	0.45267\\
3.7625	0.44791\\
3.7657	0.44321\\
3.7689	0.43857\\
3.7722	0.434\\
3.7754	0.42949\\
3.7786	0.42505\\
3.7819	0.42067\\
3.7851	0.41635\\
3.7884	0.4121\\
3.7916	0.40791\\
3.7948	0.40378\\
3.7981	0.39971\\
3.8013	0.3957\\
3.8045	0.39176\\
3.8078	0.38787\\
3.811	0.38405\\
3.8143	0.38029\\
3.8175	0.37658\\
3.8207	0.37294\\
3.824	0.36935\\
3.8272	0.36582\\
3.8304	0.36235\\
3.8337	0.35893\\
3.8369	0.35557\\
3.8402	0.35227\\
3.8434	0.34902\\
3.8466	0.34583\\
3.8499	0.34269\\
3.8531	0.3396\\
3.8564	0.33657\\
3.8596	0.33359\\
3.8628	0.33066\\
3.8661	0.32778\\
3.8693	0.32495\\
3.8725	0.32217\\
3.8758	0.31944\\
3.879	0.31676\\
3.8823	0.31413\\
3.8855	0.31154\\
3.8887	0.30899\\
3.892	0.3065\\
3.8952	0.30404\\
3.8984	0.30163\\
3.9017	0.29927\\
3.9049	0.29694\\
3.9082	0.29466\\
3.9114	0.29242\\
3.9146	0.29021\\
3.9179	0.28805\\
3.9211	0.28593\\
3.9243	0.28384\\
3.9276	0.28179\\
3.9308	0.27978\\
3.9341	0.2778\\
3.9373	0.27585\\
3.9405	0.27394\\
3.9438	0.27207\\
3.947	0.27022\\
3.9502	0.26841\\
3.9535	0.26663\\
3.9567	0.26488\\
3.96	0.26316\\
3.9632	0.26147\\
3.9664	0.2598\\
3.9697	0.25817\\
3.9729	0.25656\\
3.9762	0.25497\\
3.9794	0.25342\\
3.9826	0.25188\\
3.9859	0.25037\\
3.9891	0.24889\\
3.9923	0.24742\\
3.9956	0.24598\\
3.9988	0.24456\\
4.0021	0.24316\\
4.0053	0.24178\\
4.0085	0.24042\\
4.0118	0.23908\\
4.015	0.23776\\
4.0182	0.23645\\
4.0215	0.23516\\
4.0247	0.23389\\
4.028	0.23263\\
4.0312	0.23139\\
4.0344	0.23017\\
4.0377	0.22895\\
4.0409	0.22776\\
4.0441	0.22657\\
4.0474	0.2254\\
4.0506	0.22424\\
4.0539	0.22309\\
4.0571	0.22195\\
4.0603	0.22083\\
4.0636	0.21971\\
4.0668	0.21861\\
4.07	0.21751\\
4.0733	0.21642\\
4.0765	0.21534\\
4.0798	0.21427\\
4.083	0.21321\\
4.0862	0.21215\\
4.0895	0.21111\\
4.0927	0.21006\\
4.0959	0.20903\\
4.0992	0.208\\
4.1024	0.20698\\
4.1057	0.20596\\
4.1089	0.20495\\
4.1121	0.20394\\
4.1154	0.20294\\
4.1186	0.20194\\
4.1219	0.20094\\
4.1251	0.19995\\
4.1283	0.19896\\
4.1316	0.19798\\
4.1348	0.197\\
4.138	0.19602\\
4.1413	0.19504\\
4.1445	0.19407\\
4.1478	0.19309\\
4.151	0.19212\\
4.1542	0.19115\\
4.1575	0.19019\\
4.1607	0.18922\\
4.1639	0.18825\\
4.1672	0.18729\\
4.1704	0.18633\\
4.1737	0.18536\\
4.1769	0.1844\\
4.1801	0.18344\\
4.1834	0.18247\\
4.1866	0.18151\\
4.1898	0.18054\\
4.1931	0.17958\\
4.1963	0.17861\\
4.1996	0.17765\\
4.2028	0.17668\\
4.206	0.17571\\
4.2093	0.17474\\
4.2125	0.17377\\
4.2157	0.1728\\
4.219	0.17182\\
4.2222	0.17084\\
4.2255	0.16987\\
4.2287	0.16888\\
4.2319	0.1679\\
4.2352	0.16691\\
4.2384	0.16593\\
4.2417	0.16493\\
4.2449	0.16394\\
4.2481	0.16294\\
4.2514	0.16194\\
4.2546	0.16094\\
4.2578	0.15993\\
4.2611	0.15892\\
4.2643	0.15791\\
4.2676	0.1569\\
4.2708	0.15588\\
4.274	0.15485\\
4.2773	0.15383\\
4.2805	0.1528\\
4.2837	0.15176\\
4.287	0.15073\\
4.2902	0.14969\\
4.2935	0.14865\\
4.2967	0.1476\\
4.2999	0.14655\\
4.3032	0.14549\\
4.3064	0.14444\\
4.3096	0.14338\\
4.3129	0.14231\\
4.3161	0.14125\\
4.3194	0.14018\\
4.3226	0.1391\\
4.3258	0.13803\\
4.3291	0.13695\\
4.3323	0.13587\\
4.3355	0.13478\\
4.3388	0.1337\\
4.342	0.13261\\
4.3453	0.13152\\
4.3485	0.13042\\
4.3517	0.12933\\
4.355	0.12823\\
4.3582	0.12713\\
4.3615	0.12603\\
4.3647	0.12493\\
4.3679	0.12382\\
4.3712	0.12272\\
4.3744	0.12161\\
4.3776	0.1205\\
4.3809	0.1194\\
4.3841	0.11829\\
4.3874	0.11718\\
4.3906	0.11607\\
4.3938	0.11496\\
4.3971	0.11385\\
4.4003	0.11275\\
4.4035	0.11164\\
4.4068	0.11053\\
4.41	0.10943\\
4.4133	0.10832\\
4.4165	0.10722\\
4.4197	0.10612\\
4.423	0.10502\\
4.4262	0.10392\\
4.4294	0.10282\\
4.4327	0.10173\\
4.4359	0.10064\\
4.4392	0.099552\\
4.4424	0.098467\\
4.4456	0.097385\\
4.4489	0.096306\\
4.4521	0.095231\\
4.4553	0.094159\\
4.4586	0.093091\\
4.4618	0.092027\\
4.4651	0.090967\\
4.4683	0.089912\\
4.4715	0.088861\\
4.4748	0.087814\\
4.478	0.086773\\
4.4813	0.085736\\
4.4845	0.084705\\
4.4877	0.083678\\
4.491	0.082657\\
4.4942	0.081642\\
4.4974	0.080633\\
4.5007	0.079629\\
4.5039	0.078631\\
4.5072	0.077639\\
4.5104	0.076653\\
4.5136	0.075674\\
4.5169	0.074701\\
4.5201	0.073735\\
4.5233	0.072775\\
4.5266	0.071822\\
4.5298	0.070876\\
4.5331	0.069937\\
4.5363	0.069004\\
4.5395	0.068079\\
4.5428	0.067161\\
4.546	0.06625\\
4.5492	0.065347\\
4.5525	0.064451\\
4.5557	0.063562\\
4.559	0.062681\\
4.5622	0.061807\\
4.5654	0.060941\\
4.5687	0.060083\\
4.5719	0.059233\\
4.5751	0.05839\\
4.5784	0.057555\\
4.5816	0.056727\\
4.5849	0.055908\\
4.5881	0.055096\\
4.5913	0.054293\\
4.5946	0.053497\\
4.5978	0.052709\\
4.6011	0.051929\\
4.6043	0.051158\\
4.6075	0.050394\\
4.6108	0.049638\\
4.614	0.04889\\
4.6172	0.04815\\
4.6205	0.047418\\
4.6237	0.046694\\
4.627	0.045978\\
4.6302	0.04527\\
4.6334	0.04457\\
4.6367	0.043878\\
4.6399	0.043194\\
4.6431	0.042517\\
4.6464	0.041849\\
4.6496	0.041188\\
4.6529	0.040536\\
4.6561	0.039891\\
4.6593	0.039254\\
4.6626	0.038625\\
4.6658	0.038003\\
4.669	0.037389\\
4.6723	0.036783\\
4.6755	0.036184\\
4.6788	0.035593\\
4.682	0.03501\\
4.6852	0.034434\\
4.6885	0.033865\\
4.6917	0.033304\\
4.6949	0.03275\\
4.6982	0.032204\\
4.7014	0.031665\\
4.7047	0.031133\\
4.7079	0.030608\\
4.7111	0.030091\\
4.7144	0.02958\\
4.7176	0.029077\\
4.7209	0.02858\\
4.7241	0.028091\\
4.7273	0.027608\\
4.7306	0.027132\\
4.7338	0.026663\\
4.737	0.0262\\
4.7403	0.025745\\
4.7435	0.025295\\
4.7468	0.024853\\
4.75	0.024416\\
4.7532	0.023986\\
4.7565	0.023563\\
4.7597	0.023146\\
4.7629	0.022735\\
4.7662	0.02233\\
4.7694	0.021931\\
4.7727	0.021538\\
4.7759	0.021152\\
4.7791	0.020771\\
4.7824	0.020396\\
4.7856	0.020026\\
4.7888	0.019663\\
4.7921	0.019305\\
4.7953	0.018953\\
4.7986	0.018606\\
4.8018	0.018265\\
4.805	0.017929\\
4.8083	0.017598\\
4.8115	0.017273\\
4.8147	0.016953\\
4.818	0.016638\\
4.8212	0.016328\\
4.8245	0.016023\\
4.8277	0.015724\\
4.8309	0.015429\\
4.8342	0.015139\\
4.8374	0.014853\\
4.8406	0.014573\\
4.8439	0.014297\\
4.8471	0.014025\\
4.8504	0.013758\\
4.8536	0.013496\\
4.8568	0.013238\\
4.8601	0.012985\\
4.8633	0.012735\\
4.8666	0.01249\\
4.8698	0.012249\\
4.873	0.012013\\
4.8763	0.01178\\
4.8795	0.011551\\
4.8827	0.011327\\
4.886	0.011106\\
4.8892	0.010889\\
4.8925	0.010676\\
4.8957	0.010466\\
4.8989	0.01026\\
4.9022	0.010058\\
4.9054	0.0098594\\
4.9086	0.0096643\\
4.9119	0.0094727\\
4.9151	0.0092844\\
4.9184	0.0090996\\
4.9216	0.008918\\
4.9248	0.0087397\\
4.9281	0.0085646\\
4.9313	0.0083927\\
4.9345	0.0082238\\
4.9378	0.0080581\\
4.941	0.0078953\\
4.9443	0.0077355\\
4.9475	0.0075787\\
4.9507	0.0074247\\
4.954	0.0072735\\
4.9572	0.0071252\\
4.9604	0.0069795\\
4.9637	0.0068366\\
4.9669	0.0066963\\
4.9702	0.0065587\\
4.9734	0.0064236\\
4.9766	0.006291\\
4.9799	0.006161\\
4.9831	0.0060333\\
4.9864	0.0059081\\
4.9896	0.0057852\\
4.9928	0.0056647\\
4.9961	0.0055465\\
4.9993	0.0054305\\
5.0025	0.0053167\\
5.0058	0.0052051\\
5.009	0.0050956\\
5.0123	0.0049883\\
5.0155	0.004883\\
5.0187	0.0047797\\
5.022	0.0046784\\
5.0252	0.0045791\\
5.0284	0.0044818\\
5.0317	0.0043863\\
5.0349	0.0042927\\
5.0382	0.0042009\\
5.0414	0.0041109\\
5.0446	0.0040227\\
5.0479	0.0039362\\
5.0511	0.0038514\\
5.0543	0.0037683\\
5.0576	0.0036869\\
5.0608	0.003607\\
5.0641	0.0035288\\
5.0673	0.0034521\\
5.0705	0.003377\\
5.0738	0.0033033\\
5.077	0.0032312\\
5.0802	0.0031604\\
5.0835	0.0030912\\
5.0867	0.0030233\\
5.09	0.0029568\\
5.0932	0.0028916\\
5.0964	0.0028278\\
5.0997	0.0027652\\
5.1029	0.002704\\
5.1062	0.0026439\\
5.1094	0.0025852\\
5.1126	0.0025276\\
5.1159	0.0024712\\
5.1191	0.002416\\
5.1223	0.0023619\\
5.1256	0.002309\\
5.1288	0.0022571\\
5.1321	0.0022063\\
5.1353	0.0021566\\
5.1385	0.0021079\\
5.1418	0.0020602\\
5.145	0.0020136\\
5.1482	0.0019679\\
5.1515	0.0019232\\
5.1547	0.0018794\\
5.158	0.0018365\\
5.1612	0.0017946\\
5.1644	0.0017535\\
5.1677	0.0017133\\
5.1709	0.001674\\
5.1741	0.0016355\\
5.1774	0.0015979\\
5.1806	0.001561\\
5.1839	0.0015249\\
5.1871	0.0014896\\
5.1903	0.0014551\\
5.1936	0.0014213\\
5.1968	0.0013883\\
5.2	0.0013559\\
5.2033	0.0013243\\
5.2065	0.0012933\\
5.2098	0.0012631\\
5.213	0.0012335\\
5.2162	0.0012045\\
5.2195	0.0011762\\
5.2227	0.0011485\\
5.226	0.0011214\\
5.2292	0.0010949\\
5.2324	0.001069\\
5.2357	0.0010436\\
5.2389	0.0010188\\
5.2421	0.00099462\\
5.2454	0.00097093\\
5.2486	0.00094777\\
5.2519	0.00092513\\
5.2551	0.00090299\\
5.2583	0.00088135\\
5.2616	0.0008602\\
5.2648	0.00083953\\
5.268	0.00081932\\
5.2713	0.00079956\\
5.2745	0.00078026\\
5.2778	0.00076139\\
5.281	0.00074295\\
5.2842	0.00072493\\
5.2875	0.00070732\\
5.2907	0.00069011\\
5.2939	0.00067329\\
5.2972	0.00065687\\
5.3004	0.00064081\\
5.3037	0.00062513\\
5.3069	0.00060981\\
5.3101	0.00059484\\
5.3101	0.00036417\\
5.3069	0.00037376\\
5.3037	0.00038358\\
5.3004	0.00039365\\
5.2972	0.00040397\\
5.2939	0.00041455\\
5.2907	0.00042538\\
5.2875	0.00043648\\
5.2842	0.00044786\\
5.281	0.00045951\\
5.2778	0.00047145\\
5.2745	0.00048369\\
5.2713	0.00049622\\
5.268	0.00050906\\
5.2648	0.00052221\\
5.2616	0.00053569\\
5.2583	0.00054949\\
5.2551	0.00056363\\
5.2519	0.00057811\\
5.2486	0.00059294\\
5.2454	0.00060813\\
5.2421	0.00062369\\
5.2389	0.00063962\\
5.2357	0.00065593\\
5.2324	0.00067264\\
5.2292	0.00068975\\
5.226	0.00070727\\
5.2227	0.00072521\\
5.2195	0.00074358\\
5.2162	0.00076239\\
5.213	0.00078164\\
5.2098	0.00080136\\
5.2065	0.00082154\\
5.2033	0.0008422\\
5.2	0.00086335\\
5.1968	0.000885\\
5.1936	0.00090717\\
5.1903	0.00092985\\
5.1871	0.00095307\\
5.1839	0.00097684\\
5.1806	0.0010012\\
5.1774	0.0010261\\
5.1741	0.0010515\\
5.1709	0.0010776\\
5.1677	0.0011043\\
5.1644	0.0011316\\
5.1612	0.0011595\\
5.158	0.0011881\\
5.1547	0.0012174\\
5.1515	0.0012473\\
5.1482	0.001278\\
5.145	0.0013093\\
5.1418	0.0013413\\
5.1385	0.0013741\\
5.1353	0.0014077\\
5.1321	0.001442\\
5.1288	0.0014771\\
5.1256	0.001513\\
5.1223	0.0015498\\
5.1191	0.0015873\\
5.1159	0.0016258\\
5.1126	0.0016651\\
5.1094	0.0017053\\
5.1062	0.0017464\\
5.1029	0.0017884\\
5.0997	0.0018314\\
5.0964	0.0018754\\
5.0932	0.0019203\\
5.09	0.0019663\\
5.0867	0.0020133\\
5.0835	0.0020613\\
5.0802	0.0021105\\
5.077	0.0021607\\
5.0738	0.002212\\
5.0705	0.0022645\\
5.0673	0.0023182\\
5.0641	0.002373\\
5.0608	0.0024291\\
5.0576	0.0024864\\
5.0543	0.002545\\
5.0511	0.0026049\\
5.0479	0.0026661\\
5.0446	0.0027286\\
5.0414	0.0027925\\
5.0382	0.0028578\\
5.0349	0.0029245\\
5.0317	0.0029927\\
5.0284	0.0030624\\
5.0252	0.0031336\\
5.022	0.0032063\\
5.0187	0.0032806\\
5.0155	0.0033566\\
5.0123	0.0034341\\
5.009	0.0035133\\
5.0058	0.0035943\\
5.0025	0.0036769\\
4.9993	0.0037614\\
4.9961	0.0038476\\
4.9928	0.0039357\\
4.9896	0.0040256\\
4.9864	0.0041175\\
4.9831	0.0042113\\
4.9799	0.0043071\\
4.9766	0.0044049\\
4.9734	0.0045048\\
4.9702	0.0046067\\
4.9669	0.0047108\\
4.9637	0.0048171\\
4.9604	0.0049256\\
4.9572	0.0050363\\
4.954	0.0051494\\
4.9507	0.0052648\\
4.9475	0.0053825\\
4.9443	0.0055027\\
4.941	0.0056254\\
4.9378	0.0057506\\
4.9345	0.0058783\\
4.9313	0.0060086\\
4.9281	0.0061416\\
4.9248	0.0062773\\
4.9216	0.0064157\\
4.9184	0.006557\\
4.9151	0.006701\\
4.9119	0.006848\\
4.9086	0.0069979\\
4.9054	0.0071507\\
4.9022	0.0073067\\
4.8989	0.0074657\\
4.8957	0.0076278\\
4.8925	0.0077932\\
4.8892	0.0079618\\
4.886	0.0081336\\
4.8827	0.0083089\\
4.8795	0.0084876\\
4.8763	0.0086697\\
4.873	0.0088553\\
4.8698	0.0090446\\
4.8666	0.0092374\\
4.8633	0.009434\\
4.8601	0.0096343\\
4.8568	0.0098384\\
4.8536	0.010046\\
4.8504	0.010258\\
4.8471	0.010474\\
4.8439	0.010694\\
4.8406	0.010918\\
4.8374	0.011146\\
4.8342	0.011379\\
4.8309	0.011616\\
4.8277	0.011857\\
4.8245	0.012102\\
4.8212	0.012352\\
4.818	0.012607\\
4.8147	0.012866\\
4.8115	0.01313\\
4.8083	0.013398\\
4.805	0.013672\\
4.8018	0.01395\\
4.7986	0.014233\\
4.7953	0.014521\\
4.7921	0.014815\\
4.7888	0.015113\\
4.7856	0.015416\\
4.7824	0.015725\\
4.7791	0.01604\\
4.7759	0.016359\\
4.7727	0.016684\\
4.7694	0.017015\\
4.7662	0.017351\\
4.7629	0.017693\\
4.7597	0.01804\\
4.7565	0.018394\\
4.7532	0.018753\\
4.75	0.019118\\
4.7468	0.019489\\
4.7435	0.019867\\
4.7403	0.02025\\
4.737	0.02064\\
4.7338	0.021036\\
4.7306	0.021438\\
4.7273	0.021846\\
4.7241	0.022261\\
4.7209	0.022683\\
4.7176	0.023111\\
4.7144	0.023546\\
4.7111	0.023987\\
4.7079	0.024436\\
4.7047	0.024891\\
4.7014	0.025353\\
4.6982	0.025822\\
4.6949	0.026298\\
4.6917	0.026781\\
4.6885	0.027271\\
4.6852	0.027768\\
4.682	0.028273\\
4.6788	0.028784\\
4.6755	0.029304\\
4.6723	0.02983\\
4.669	0.030364\\
4.6658	0.030905\\
4.6626	0.031454\\
4.6593	0.032011\\
4.6561	0.032575\\
4.6529	0.033146\\
4.6496	0.033726\\
4.6464	0.034312\\
4.6431	0.034907\\
4.6399	0.03551\\
4.6367	0.03612\\
4.6334	0.036738\\
4.6302	0.037364\\
4.627	0.037998\\
4.6237	0.038639\\
4.6205	0.039289\\
4.6172	0.039946\\
4.614	0.040612\\
4.6108	0.041285\\
4.6075	0.041966\\
4.6043	0.042655\\
4.6011	0.043352\\
4.5978	0.044057\\
4.5946	0.04477\\
4.5913	0.045491\\
4.5881	0.04622\\
4.5849	0.046956\\
4.5816	0.047701\\
4.5784	0.048453\\
4.5751	0.049213\\
4.5719	0.049981\\
4.5687	0.050756\\
4.5654	0.05154\\
4.5622	0.052331\\
4.559	0.05313\\
4.5557	0.053936\\
4.5525	0.05475\\
4.5492	0.055571\\
4.546	0.0564\\
4.5428	0.057236\\
4.5395	0.05808\\
4.5363	0.05893\\
4.5331	0.059788\\
4.5298	0.060654\\
4.5266	0.061526\\
4.5233	0.062405\\
4.5201	0.063291\\
4.5169	0.064184\\
4.5136	0.065084\\
4.5104	0.06599\\
4.5072	0.066903\\
4.5039	0.067823\\
4.5007	0.068748\\
4.4974	0.069681\\
4.4942	0.070619\\
4.491	0.071563\\
4.4877	0.072514\\
4.4845	0.07347\\
4.4813	0.074432\\
4.478	0.0754\\
4.4748	0.076373\\
4.4715	0.077352\\
4.4683	0.078336\\
4.4651	0.079325\\
4.4618	0.08032\\
4.4586	0.081319\\
4.4553	0.082323\\
4.4521	0.083332\\
4.4489	0.084346\\
4.4456	0.085364\\
4.4424	0.086386\\
4.4392	0.087413\\
4.4359	0.088443\\
4.4327	0.089478\\
4.4294	0.090516\\
4.4262	0.091558\\
4.423	0.092604\\
4.4197	0.093653\\
4.4165	0.094705\\
4.4133	0.09576\\
4.41	0.096819\\
4.4068	0.09788\\
4.4035	0.098943\\
4.4003	0.10001\\
4.3971	0.10108\\
4.3938	0.10215\\
4.3906	0.10322\\
4.3874	0.1043\\
4.3841	0.10537\\
4.3809	0.10645\\
4.3776	0.10753\\
4.3744	0.10861\\
4.3712	0.1097\\
4.3679	0.11078\\
4.3647	0.11186\\
4.3615	0.11295\\
4.3582	0.11403\\
4.355	0.11512\\
4.3517	0.11621\\
4.3485	0.11729\\
4.3453	0.11838\\
4.342	0.11947\\
4.3388	0.12055\\
4.3355	0.12164\\
4.3323	0.12272\\
4.3291	0.12381\\
4.3258	0.12489\\
4.3226	0.12597\\
4.3194	0.12705\\
4.3161	0.12813\\
4.3129	0.12921\\
4.3096	0.13029\\
4.3064	0.13136\\
4.3032	0.13244\\
4.2999	0.13351\\
4.2967	0.13457\\
4.2935	0.13564\\
4.2902	0.1367\\
4.287	0.13777\\
4.2837	0.13882\\
4.2805	0.13988\\
4.2773	0.14093\\
4.274	0.14198\\
4.2708	0.14303\\
4.2676	0.14407\\
4.2643	0.14511\\
4.2611	0.14614\\
4.2578	0.14717\\
4.2546	0.1482\\
4.2514	0.14923\\
4.2481	0.15024\\
4.2449	0.15126\\
4.2417	0.15227\\
4.2384	0.15328\\
4.2352	0.15428\\
4.2319	0.15528\\
4.2287	0.15628\\
4.2255	0.15727\\
4.2222	0.15825\\
4.219	0.15923\\
4.2157	0.16021\\
4.2125	0.16118\\
4.2093	0.16215\\
4.206	0.16311\\
4.2028	0.16407\\
4.1996	0.16503\\
4.1963	0.16598\\
4.1931	0.16692\\
4.1898	0.16786\\
4.1866	0.1688\\
4.1834	0.16973\\
4.1801	0.17066\\
4.1769	0.17159\\
4.1737	0.17251\\
4.1704	0.17343\\
4.1672	0.17434\\
4.1639	0.17525\\
4.1607	0.17616\\
4.1575	0.17707\\
4.1542	0.17797\\
4.151	0.17887\\
4.1478	0.17976\\
4.1445	0.18066\\
4.1413	0.18155\\
4.138	0.18244\\
4.1348	0.18333\\
4.1316	0.18422\\
4.1283	0.18511\\
4.1251	0.186\\
4.1219	0.18688\\
4.1186	0.18777\\
4.1154	0.18866\\
4.1121	0.18955\\
4.1089	0.19043\\
4.1057	0.19132\\
4.1024	0.19222\\
4.0992	0.19311\\
4.0959	0.19401\\
4.0927	0.1949\\
4.0895	0.19581\\
4.0862	0.19671\\
4.083	0.19762\\
4.0798	0.19853\\
4.0765	0.19945\\
4.0733	0.20037\\
4.07	0.2013\\
4.0668	0.20223\\
4.0636	0.20317\\
4.0603	0.20412\\
4.0571	0.20507\\
4.0539	0.20603\\
4.0506	0.207\\
4.0474	0.20798\\
4.0441	0.20897\\
4.0409	0.20996\\
4.0377	0.21097\\
4.0344	0.21198\\
4.0312	0.21301\\
4.028	0.21404\\
4.0247	0.21509\\
4.0215	0.21616\\
4.0182	0.21723\\
4.015	0.21832\\
4.0118	0.21942\\
4.0085	0.22053\\
4.0053	0.22166\\
4.0021	0.22281\\
3.9988	0.22397\\
3.9956	0.22515\\
3.9923	0.22634\\
3.9891	0.22756\\
3.9859	0.22879\\
3.9826	0.23004\\
3.9794	0.2313\\
3.9762	0.23259\\
3.9729	0.2339\\
3.9697	0.23523\\
3.9664	0.23658\\
3.9632	0.23796\\
3.96	0.23936\\
3.9567	0.24078\\
3.9535	0.24222\\
3.9502	0.24369\\
3.947	0.24519\\
3.9438	0.24671\\
3.9405	0.24825\\
3.9373	0.24983\\
3.9341	0.25143\\
3.9308	0.25307\\
3.9276	0.25473\\
3.9243	0.25642\\
3.9211	0.25814\\
3.9179	0.2599\\
3.9146	0.26168\\
3.9114	0.2635\\
3.9082	0.26536\\
3.9049	0.26724\\
3.9017	0.26917\\
3.8984	0.27112\\
3.8952	0.27312\\
3.892	0.27515\\
3.8887	0.27722\\
3.8855	0.27933\\
3.8823	0.28148\\
3.879	0.28366\\
3.8758	0.28589\\
3.8725	0.28816\\
3.8693	0.29047\\
3.8661	0.29283\\
3.8628	0.29523\\
3.8596	0.29767\\
3.8564	0.30016\\
3.8531	0.30269\\
3.8499	0.30527\\
3.8466	0.3079\\
3.8434	0.31058\\
3.8402	0.3133\\
3.8369	0.31608\\
3.8337	0.3189\\
3.8304	0.32177\\
3.8272	0.3247\\
3.824	0.32768\\
3.8207	0.33071\\
3.8175	0.33379\\
3.8143	0.33693\\
3.811	0.34012\\
3.8078	0.34337\\
3.8045	0.34667\\
3.8013	0.35003\\
3.7981	0.35345\\
3.7948	0.35692\\
3.7916	0.36045\\
3.7884	0.36404\\
3.7851	0.36769\\
3.7819	0.3714\\
3.7786	0.37517\\
3.7754	0.379\\
3.7722	0.38289\\
3.7689	0.38685\\
3.7657	0.39086\\
3.7625	0.39493\\
3.7592	0.39907\\
3.756	0.40327\\
3.7527	0.40754\\
3.7495	0.41187\\
3.7463	0.41626\\
3.743	0.42071\\
3.7398	0.42523\\
3.7366	0.42982\\
3.7333	0.43447\\
3.7301	0.43918\\
3.7268	0.44396\\
3.7236	0.4488\\
3.7204	0.45371\\
3.7171	0.45868\\
3.7139	0.46372\\
3.7106	0.46882\\
3.7074	0.47399\\
3.7042	0.47922\\
3.7009	0.48452\\
3.6977	0.48988\\
3.6945	0.4953\\
3.6912	0.50079\\
3.688	0.50634\\
3.6847	0.51196\\
3.6815	0.51764\\
3.6783	0.52338\\
3.675	0.52918\\
3.6718	0.53504\\
3.6686	0.54097\\
3.6653	0.54695\\
3.6621	0.553\\
3.6588	0.5591\\
3.6556	0.56527\\
3.6524	0.57149\\
3.6491	0.57776\\
3.6459	0.5841\\
3.6427	0.59049\\
3.6394	0.59693\\
3.6362	0.60343\\
3.6329	0.60998\\
3.6297	0.61658\\
3.6265	0.62324\\
3.6232	0.62994\\
3.62	0.63669\\
3.6168	0.64349\\
3.6135	0.65033\\
3.6103	0.65722\\
3.607	0.66415\\
3.6038	0.67112\\
3.6006	0.67813\\
3.5973	0.68519\\
3.5941	0.69228\\
3.5908	0.6994\\
3.5876	0.70656\\
3.5844	0.71376\\
3.5811	0.72098\\
3.5779	0.72823\\
3.5747	0.73551\\
3.5714	0.74282\\
3.5682	0.75015\\
3.5649	0.75751\\
3.5617	0.76488\\
3.5585	0.77227\\
3.5552	0.77968\\
3.552	0.7871\\
3.5488	0.79454\\
3.5455	0.80198\\
3.5423	0.80943\\
3.539	0.81689\\
3.5358	0.82435\\
3.5326	0.83181\\
3.5293	0.83927\\
3.5261	0.84673\\
3.5229	0.85418\\
3.5196	0.86163\\
3.5164	0.86906\\
3.5131	0.87648\\
3.5099	0.88389\\
3.5067	0.89127\\
3.5034	0.89864\\
3.5002	0.90598\\
3.497	0.9133\\
3.4937	0.92059\\
3.4905	0.92784\\
3.4872	0.93507\\
3.484	0.94225\\
3.4808	0.9494\\
3.4775	0.95651\\
3.4743	0.96357\\
3.471	0.97059\\
3.4678	0.97755\\
3.4646	0.98446\\
3.4613	0.99132\\
3.4581	0.99812\\
3.4549	1.0049\\
3.4516	1.0115\\
3.4484	1.0181\\
3.4451	1.0247\\
3.4419	1.0311\\
3.4387	1.0375\\
3.4354	1.0438\\
3.4322	1.05\\
3.429	1.0562\\
3.4257	1.0622\\
3.4225	1.0682\\
3.4192	1.0741\\
3.416	1.0799\\
3.4128	1.0855\\
3.4095	1.0911\\
3.4063	1.0966\\
3.4031	1.102\\
3.3998	1.1072\\
3.3966	1.1123\\
3.3933	1.1174\\
3.3901	1.1223\\
3.3869	1.1271\\
3.3836	1.1317\\
3.3804	1.1362\\
3.3772	1.1406\\
3.3739	1.1449\\
3.3707	1.149\\
3.3674	1.153\\
3.3642	1.1569\\
3.361	1.1606\\
3.3577	1.1641\\
3.3545	1.1675\\
3.3513	1.1708\\
3.348	1.1739\\
3.3448	1.1768\\
3.3415	1.1796\\
3.3383	1.1822\\
3.3351	1.1846\\
3.3318	1.1869\\
3.3286	1.189\\
3.3253	1.1909\\
3.3221	1.1926\\
3.3189	1.1942\\
3.3156	1.1955\\
3.3124	1.1967\\
3.3092	1.1977\\
3.3059	1.1985\\
3.3027	1.1992\\
3.2994	1.1996\\
3.2962	1.1998\\
3.293	1.1999\\
3.2897	1.1998\\
3.2865	1.1994\\
3.2833	1.1989\\
3.28	1.1982\\
3.2768	1.1973\\
3.2735	1.1962\\
3.2703	1.1949\\
3.2671	1.1935\\
3.2638	1.1918\\
3.2606	1.19\\
3.2574	1.188\\
3.2541	1.1858\\
3.2509	1.1834\\
3.2476	1.1809\\
3.2444	1.1781\\
3.2412	1.1752\\
3.2379	1.1722\\
3.2347	1.1689\\
3.2315	1.1655\\
3.2282	1.162\\
3.225	1.1583\\
3.2217	1.1544\\
3.2185	1.1503\\
3.2153	1.1461\\
3.212	1.1418\\
3.2088	1.1373\\
3.2055	1.1327\\
3.2023	1.1279\\
3.1991	1.1229\\
3.1958	1.1179\\
3.1926	1.1126\\
3.1894	1.1073\\
3.1861	1.1018\\
3.1829	1.0962\\
3.1796	1.0905\\
3.1764	1.0846\\
3.1732	1.0786\\
3.1699	1.0725\\
3.1667	1.0663\\
3.1635	1.0599\\
3.1602	1.0535\\
3.157	1.0469\\
3.1537	1.0403\\
3.1505	1.0335\\
3.1473	1.0266\\
3.144	1.0196\\
3.1408	1.0126\\
3.1376	1.0054\\
3.1343	0.99815\\
3.1311	0.9908\\
3.1278	0.98338\\
3.1246	0.97587\\
3.1214	0.96828\\
3.1181	0.96062\\
3.1149	0.95288\\
3.1117	0.94507\\
3.1084	0.9372\\
3.1052	0.92926\\
3.1019	0.92126\\
3.0987	0.91319\\
3.0955	0.90508\\
3.0922	0.89691\\
3.089	0.88869\\
3.0857	0.88042\\
3.0825	0.87211\\
3.0793	0.86375\\
3.076	0.85536\\
3.0728	0.84693\\
3.0696	0.83847\\
3.0663	0.82997\\
3.0631	0.82145\\
3.0598	0.8129\\
3.0566	0.80433\\
3.0534	0.79575\\
3.0501	0.78714\\
3.0469	0.77852\\
3.0437	0.76989\\
3.0404	0.76125\\
3.0372	0.7526\\
3.0339	0.74395\\
3.0307	0.7353\\
3.0275	0.72665\\
3.0242	0.71801\\
3.021	0.70937\\
3.0178	0.70073\\
3.0145	0.69211\\
3.0113	0.68351\\
3.008	0.67491\\
3.0048	0.66634\\
3.0016	0.65779\\
2.9983	0.64925\\
2.9951	0.64075\\
2.9919	0.63227\\
2.9886	0.62381\\
2.9854	0.61539\\
2.9821	0.607\\
2.9789	0.59865\\
2.9757	0.59033\\
2.9724	0.58205\\
2.9692	0.57381\\
2.9659	0.56561\\
2.9627	0.55745\\
2.9595	0.54934\\
2.9562	0.54128\\
2.953	0.53326\\
2.9498	0.5253\\
2.9465	0.51739\\
2.9433	0.50952\\
2.94	0.50172\\
2.9368	0.49397\\
2.9336	0.48627\\
2.9303	0.47863\\
2.9271	0.47106\\
2.9239	0.46354\\
2.9206	0.45608\\
2.9174	0.44869\\
2.9141	0.44136\\
2.9109	0.43409\\
2.9077	0.42689\\
2.9044	0.41976\\
2.9012	0.41269\\
2.898	0.40569\\
2.8947	0.39876\\
2.8915	0.3919\\
2.8882	0.38511\\
2.885	0.37839\\
2.8818	0.37174\\
2.8785	0.36516\\
2.8753	0.35865\\
2.8721	0.35222\\
2.8688	0.34586\\
2.8656	0.33957\\
2.8623	0.33336\\
2.8591	0.32722\\
2.8559	0.32116\\
2.8526	0.31517\\
2.8494	0.30925\\
2.8461	0.30341\\
2.8429	0.29764\\
2.8397	0.29195\\
2.8364	0.28633\\
2.8332	0.28079\\
2.83	0.27533\\
2.8267	0.26993\\
2.8235	0.26462\\
2.8202	0.25937\\
2.817	0.25421\\
2.8138	0.24911\\
2.8105	0.24409\\
2.8073	0.23915\\
2.8041	0.23427\\
2.8008	0.22948\\
2.7976	0.22475\\
2.7943	0.2201\\
2.7911	0.21552\\
2.7879	0.21101\\
2.7846	0.20657\\
2.7814	0.2022\\
2.7782	0.19791\\
2.7749	0.19368\\
2.7717	0.18952\\
2.7684	0.18544\\
2.7652	0.18142\\
2.762	0.17747\\
2.7587	0.17359\\
2.7555	0.16977\\
2.7523	0.16602\\
2.749	0.16234\\
2.7458	0.15872\\
2.7425	0.15517\\
2.7393	0.15168\\
2.7361	0.14825\\
2.7328	0.14489\\
2.7296	0.14159\\
2.7263	0.13835\\
2.7231	0.13517\\
2.7199	0.13205\\
2.7166	0.12899\\
2.7134	0.12599\\
2.7102	0.12305\\
2.7069	0.12016\\
2.7037	0.11733\\
2.7004	0.11456\\
2.6972	0.11184\\
2.694	0.10917\\
2.6907	0.10656\\
2.6875	0.10401\\
2.6843	0.1015\\
2.681	0.099047\\
2.6778	0.096644\\
2.6745	0.09429\\
2.6713	0.091986\\
2.6681	0.089731\\
2.6648	0.087523\\
2.6616	0.085362\\
2.6584	0.083247\\
2.6551	0.081178\\
2.6519	0.079154\\
2.6486	0.077174\\
2.6454	0.075238\\
2.6422	0.073344\\
2.6389	0.071492\\
2.6357	0.069681\\
2.6325	0.067911\\
2.6292	0.066181\\
2.626	0.064489\\
2.6227	0.062837\\
2.6195	0.061222\\
2.6163	0.059644\\
2.613	0.058103\\
2.6098	0.056597\\
2.6066	0.055127\\
2.6033	0.053691\\
2.6001	0.052289\\
2.5968	0.05092\\
2.5936	0.049583\\
2.5904	0.048278\\
2.5871	0.047005\\
2.5839	0.045762\\
2.5806	0.044549\\
2.5774	0.043366\\
2.5742	0.042211\\
2.5709	0.041085\\
2.5677	0.039986\\
2.5645	0.038915\\
2.5612	0.03787\\
2.558	0.036851\\
2.5547	0.035857\\
2.5515	0.034888\\
2.5483	0.033944\\
2.545	0.033023\\
2.5418	0.032126\\
2.5386	0.031251\\
2.5353	0.030399\\
2.5321	0.029568\\
2.5288	0.028759\\
2.5256	0.027971\\
2.5224	0.027203\\
2.5191	0.026454\\
2.5159	0.025726\\
2.5127	0.025016\\
2.5094	0.024325\\
2.5062	0.023652\\
2.5029	0.022996\\
2.4997	0.022358\\
2.4965	0.021737\\
2.4932	0.021132\\
2.49	0.020543\\
2.4868	0.01997\\
2.4835	0.019412\\
2.4803	0.01887\\
2.477	0.018341\\
2.4738	0.017827\\
2.4706	0.017327\\
2.4673	0.01684\\
2.4641	0.016367\\
2.4608	0.015906\\
2.4576	0.015458\\
2.4544	0.015022\\
2.4511	0.014598\\
2.4479	0.014186\\
2.4447	0.013784\\
2.4414	0.013394\\
2.4382	0.013015\\
2.4349	0.012646\\
2.4317	0.012287\\
2.4285	0.011939\\
2.4252	0.0116\\
2.422	0.01127\\
2.4188	0.010949\\
2.4155	0.010638\\
2.4123	0.010335\\
2.409	0.01004\\
2.4058	0.009754\\
2.4026	0.0094758\\
2.3993	0.0092053\\
2.3961	0.0089425\\
2.3929	0.008687\\
2.3896	0.0084388\\
2.3864	0.0081975\\
2.3831	0.0079631\\
2.3799	0.0077352\\
2.3767	0.0075139\\
2.3734	0.0072987\\
2.3702	0.0070897\\
2.367	0.0068866\\
2.3637	0.0066893\\
2.3605	0.0064975\\
2.3572	0.0063113\\
2.354	0.0061303\\
2.3508	0.0059545\\
2.3475	0.0057837\\
2.3443	0.0056178\\
2.341	0.0054566\\
2.3378	0.0053\\
2.3346	0.0051479\\
2.3313	0.0050001\\
2.3281	0.0048566\\
2.3249	0.0047172\\
2.3216	0.0045818\\
2.3184	0.0044503\\
2.3151	0.0043226\\
2.3119	0.0041985\\
2.3087	0.004078\\
2.3054	0.003961\\
2.3022	0.0038473\\
2.299	0.0037369\\
2.2957	0.0036297\\
2.2925	0.0035256\\
2.2892	0.0034245\\
2.286	0.0033263\\
2.2828	0.0032309\\
2.2795	0.0031383\\
2.2763	0.0030484\\
2.2731	0.0029611\\
2.2698	0.0028763\\
2.2666	0.002794\\
2.2633	0.002714\\
2.2601	0.0026363\\
2.2569	0.0025609\\
2.2536	0.0024877\\
2.2504	0.0024166\\
2.2472	0.0023476\\
2.2439	0.0022806\\
2.2407	0.0022155\\
2.2374	0.0021523\\
2.2342	0.002091\\
2.231	0.0020314\\
2.2277	0.0019735\\
2.2245	0.0019174\\
2.2212	0.0018628\\
2.218	0.0018099\\
2.2148	0.0017584\\
2.2115	0.0017085\\
2.2083	0.00166\\
2.2051	0.001613\\
2.2018	0.0015673\\
2.1986	0.0015229\\
2.1953	0.0014798\\
2.1921	0.001438\\
2.1889	0.0013974\\
2.1856	0.0013579\\
2.1824	0.0013196\\
2.1792	0.0012824\\
2.1759	0.0012463\\
2.1727	0.0012113\\
2.1694	0.0011772\\
2.1662	0.0011442\\
2.163	0.0011121\\
2.1597	0.0010809\\
2.1565	0.0010506\\
2.1533	0.0010212\\
2.15	0.00099267\\
2.1468	0.00096495\\
2.1435	0.00093802\\
2.1403	0.00091188\\
2.1371	0.00088649\\
2.1338	0.00086183\\
2.1306	0.00083788\\
2.1274	0.00081462\\
2.1241	0.00079203\\
2.1209	0.00077009\\
2.1176	0.00074878\\
2.1144	0.00072808\\
2.1112	0.00070797\\
2.1079	0.00068844\\
2.1047	0.00066947\\
2.1014	0.00065103\\
2.0982	0.00063313\\
2.095	0.00061574\\
2.0917	0.00059884\\
2.0885	0.00058242\\
2.0853	0.00056647\\
2.082	0.00055097\\
2.0788	0.00053592\\
2.0755	0.00052128\\
}--cycle;
\addplot[ybar interval, fill=black, fill opacity=0.35, area legend, draw=none] table[row sep=crcr, x=Lower, y=Count] {%
Lower	Upper	Count\\
2.0755	2.3244	0.075022\\
2.3244	2.5732	0.13933\\
2.5732	2.822	0.35368\\
2.822	3.0708	0.76094\\
3.0708	3.3196	0.94314\\
3.3196	3.5684	0.75022\\
3.5684	3.8172	0.52516\\
3.8172	4.0661	0.2465\\
4.0661	4.3149	0.10717\\
4.3149	4.5637	0.053587\\
4.5637	4.8125	0.04287\\
4.8125	5.0613	0\\
5.0613	5.3101	0.021435\\
5.3101	5.3101	0.021435\\
};
\addlegendentry{Istogramma reale}

\addplot [color=black]
  table[row sep=crcr]{%
2.0755	0.016864\\
2.1144	0.021976\\
2.15	0.027791\\
2.1824	0.034175\\
2.2115	0.040944\\
2.2407	0.048804\\
2.2666	0.056806\\
2.2925	0.065853\\
2.3184	0.076034\\
2.3443	0.087435\\
2.3702	0.10014\\
2.3929	0.11239\\
2.4155	0.12576\\
2.4382	0.14028\\
2.4608	0.15599\\
2.4835	0.17292\\
2.5062	0.19111\\
2.5288	0.21056\\
2.5515	0.23127\\
2.5774	0.25648\\
2.6033	0.28329\\
2.6292	0.31165\\
2.6584	0.34529\\
2.6875	0.38062\\
2.7199	0.42161\\
2.7555	0.4684\\
2.8041	0.5341\\
2.8947	0.65712\\
2.9303	0.70335\\
2.9595	0.73941\\
2.9854	0.76973\\
3.008	0.79468\\
3.0307	0.81793\\
3.0501	0.83636\\
3.0696	0.8533\\
3.089	0.86863\\
3.1052	0.88011\\
3.1214	0.89036\\
3.1376	0.89934\\
3.1537	0.907\\
3.1699	0.91332\\
3.1861	0.91826\\
3.1991	0.92122\\
3.212	0.92327\\
3.225	0.92442\\
3.2379	0.92466\\
3.2509	0.92399\\
3.2638	0.92242\\
3.2768	0.91996\\
3.2897	0.91661\\
3.3059	0.91119\\
3.3221	0.90444\\
3.3383	0.89637\\
3.3545	0.88705\\
3.3707	0.8765\\
3.3869	0.86479\\
3.4063	0.84927\\
3.4257	0.83225\\
3.4451	0.81384\\
3.4678	0.79078\\
3.4905	0.76619\\
3.5164	0.73648\\
3.5455	0.70139\\
3.5811	0.65673\\
3.6362	0.58564\\
3.7042	0.49801\\
3.743	0.44963\\
3.7754	0.4109\\
3.8045	0.37756\\
3.8337	0.34584\\
3.8596	0.31913\\
3.8855	0.29389\\
3.9114	0.27016\\
3.9373	0.24797\\
3.9632	0.22732\\
3.9891	0.20819\\
4.015	0.19056\\
4.0409	0.17437\\
4.0668	0.15957\\
4.0927	0.14609\\
4.1186	0.13385\\
4.1445	0.12277\\
4.1704	0.11277\\
4.1996	0.10269\\
4.2287	0.093758\\
4.2611	0.085015\\
4.2935	0.077372\\
4.3291	0.070062\\
4.3679	0.063196\\
4.41	0.056832\\
4.4586	0.050596\\
4.5136	0.044623\\
4.5784	0.03868\\
4.6593	0.032391\\
4.7565	0.025978\\
4.8666	0.019813\\
4.9799	0.014524\\
5.0964	0.010153\\
5.2162	0.006727\\
5.3101	0.0047195\\
};
\addlegendentry{Adatt. BLN reale}

            % \input{../../../.tex/20/KS/SizeDistributionFittingsfig2.tex}
    
            \nextgroupplot[%
                xmin=2.0755,xmax=5.4719,
                % xmin=1.9138,xmax=5.4719,
                % ymin=0,ymax=1.4,
            ] % This file was created by matlab2tikz.
%
\definecolor{mycolor1}{rgb}{0.00000,0.44706,0.69804}%
%
\addplot[ybar interval, fill=mycolor1, fill opacity=0.35, area legend, draw=none] table[row sep=crcr, x=Lower, y=Count] {%
Lower	Upper	Count\\
2.0755	2.3244	0\\
2.3244	2.5732	0.0069664\\
2.5732	2.822	0.083275\\
2.822	3.0708	0.48518\\
3.0708	3.3196	1.1996\\
3.3196	3.5684	1.1326\\
3.5684	3.8172	0.52698\\
3.8172	4.0661	0.254\\
4.0661	4.3149	0.17255\\
4.3149	4.5637	0.12829\\
4.5637	4.8125	0.02958\\
4.8125	5.0613	0\\
5.0613	5.3101	0\\
5.3101	5.3101	0\\
};
\addlegendentry{Istogramma esatto}

\addplot [color=mycolor1, only marks, every error bar/.append style={opacity=0.45}, mark=*, mark size=0pt, draw=none, forget plot]
 plot [error bars/.cd, y dir=both, y explicit, error bar style={line width=1pt, color=mycolor1}, error mark options={mark=none,mark size=0pt}]
 table[row sep=crcr, y error plus index=2, y error minus index=3]{%
2.2	0	0	0\\
2.4488	0.0069664	0.0018479	0.0018479\\
2.6976	0.083275	0.0057258	0.0057258\\
2.9464	0.48518	0.010548	0.010548\\
3.1952	1.1996	0.016387	0.016387\\
3.444	1.1326	0.016065	0.016065\\
3.6928	0.52698	0.011557	0.011557\\
3.9417	0.254	0.0081973	0.0081973\\
4.1905	0.17255	0.0061315	0.0061315\\
4.4393	0.12829	0.0043745	0.0043745\\
4.6881	0.02958	0.0020527	0.0020527\\
4.9369	0	0	0\\
5.1857	0	0	0\\
};
\addplot [color=mycolor1]
  table[row sep=crcr]{%
2.0755	0.00030664\\
2.2633	0.0012405\\
2.3572	0.0026999\\
2.4188	0.0046988\\
2.4641	0.0072103\\
2.4997	0.010179\\
2.5321	0.013979\\
2.558	0.018033\\
2.5806	0.022524\\
2.6033	0.028098\\
2.6227	0.03391\\
2.6422	0.040845\\
2.6584	0.047614\\
2.6745	0.055402\\
2.6907	0.064333\\
2.7069	0.074536\\
2.7231	0.08615\\
2.7393	0.099316\\
2.7555	0.11418\\
2.7684	0.1274\\
2.7814	0.14187\\
2.7943	0.15767\\
2.8073	0.17487\\
2.8202	0.19355\\
2.8332	0.21375\\
2.8494	0.24125\\
2.8656	0.27131\\
2.8818	0.30401\\
2.898	0.33938\\
2.9141	0.37743\\
2.9303	0.41812\\
2.9465	0.46137\\
2.9659	0.51647\\
2.9854	0.57476\\
3.008	0.64624\\
3.0372	0.74228\\
3.1246	1.0342\\
3.144	1.0941\\
3.1602	1.141\\
3.1764	1.1847\\
3.1894	1.217\\
3.2023	1.2466\\
3.2153	1.2732\\
3.225	1.2912\\
3.2347	1.3073\\
3.2444	1.3216\\
3.2541	1.3339\\
3.2638	1.3441\\
3.2735	1.3523\\
3.28	1.3566\\
3.2865	1.36\\
3.293	1.3625\\
3.2994	1.364\\
3.3059	1.3645\\
3.3124	1.3641\\
3.3189	1.3628\\
3.3253	1.3606\\
3.3318	1.3575\\
3.3383	1.3534\\
3.3448	1.3484\\
3.3545	1.3394\\
3.3642	1.3284\\
3.3739	1.3155\\
3.3836	1.3009\\
3.3933	1.2846\\
3.4063	1.2604\\
3.4192	1.2335\\
3.4322	1.2043\\
3.4484	1.1648\\
3.4646	1.1225\\
3.484	1.0689\\
3.5099	0.99393\\
3.5973	0.73774\\
3.62	0.676\\
3.6394	0.6259\\
3.6588	0.57879\\
3.675	0.54199\\
3.6912	0.50754\\
3.7074	0.47549\\
3.7236	0.44587\\
3.7398	0.41866\\
3.756	0.39383\\
3.7722	0.3713\\
3.7884	0.351\\
3.8013	0.33628\\
3.8143	0.32286\\
3.8304	0.30779\\
3.8466	0.2945\\
3.8628	0.28285\\
3.879	0.27269\\
3.8952	0.26388\\
3.9114	0.25626\\
3.9276	0.24969\\
3.947	0.243\\
3.9664	0.2374\\
3.9891	0.23194\\
4.0182	0.22612\\
4.0603	0.21898\\
4.1219	0.20852\\
4.1575	0.20134\\
4.1898	0.19369\\
4.2222	0.18485\\
4.2546	0.1748\\
4.287	0.16361\\
4.3226	0.15018\\
4.3679	0.13186\\
4.4877	0.082703\\
4.5233	0.069574\\
4.5557	0.058714\\
4.5881	0.049\\
4.6205	0.040478\\
4.6529	0.03313\\
4.6852	0.026889\\
4.7209	0.021188\\
4.7565	0.016563\\
4.7953	0.012561\\
4.8406	0.0090185\\
4.8925	0.0061208\\
4.954	0.00383\\
5.0317	0.0021021\\
5.1385	0.00091533\\
5.3101	0.00023711\\
};
\addlegendentry{Adatt. BLN esatto}


\addplot[area legend, draw=none, fill=mycolor1, fill opacity=0.15, forget plot]
table[row sep=crcr] {%
x	y\\
2.0755	0.000429\\
2.0788	0.0004389\\
2.082	0.00044902\\
2.0853	0.00045937\\
2.0885	0.00046995\\
2.0917	0.00048077\\
2.095	0.00049183\\
2.0982	0.00050314\\
2.1014	0.00051471\\
2.1047	0.00052654\\
2.1079	0.00053863\\
2.1112	0.000551\\
2.1144	0.00056364\\
2.1176	0.00057658\\
2.1209	0.0005898\\
2.1241	0.00060333\\
2.1274	0.00061717\\
2.1306	0.00063132\\
2.1338	0.0006458\\
2.1371	0.0006606\\
2.1403	0.00067575\\
2.1435	0.00069125\\
2.1468	0.0007071\\
2.15	0.00072332\\
2.1533	0.00073991\\
2.1565	0.0007569\\
2.1597	0.00077427\\
2.163	0.00079205\\
2.1662	0.00081025\\
2.1694	0.00082888\\
2.1727	0.00084794\\
2.1759	0.00086746\\
2.1792	0.00088743\\
2.1824	0.00090788\\
2.1856	0.00092882\\
2.1889	0.00095026\\
2.1921	0.00097221\\
2.1953	0.0009947\\
2.1986	0.0010177\\
2.2018	0.0010413\\
2.2051	0.0010655\\
2.2083	0.0010902\\
2.2115	0.0011156\\
2.2148	0.0011415\\
2.218	0.0011682\\
2.2212	0.0011954\\
2.2245	0.0012234\\
2.2277	0.0012521\\
2.231	0.0012815\\
2.2342	0.0013116\\
2.2374	0.0013425\\
2.2407	0.0013741\\
2.2439	0.0014066\\
2.2472	0.0014399\\
2.2504	0.0014741\\
2.2536	0.0015091\\
2.2569	0.0015451\\
2.2601	0.001582\\
2.2633	0.0016199\\
2.2666	0.0016588\\
2.2698	0.0016987\\
2.2731	0.0017396\\
2.2763	0.0017817\\
2.2795	0.0018249\\
2.2828	0.0018692\\
2.286	0.0019147\\
2.2892	0.0019615\\
2.2925	0.0020096\\
2.2957	0.002059\\
2.299	0.0021097\\
2.3022	0.0021619\\
2.3054	0.0022155\\
2.3087	0.0022706\\
2.3119	0.0023273\\
2.3151	0.0023855\\
2.3184	0.0024454\\
2.3216	0.002507\\
2.3249	0.0025704\\
2.3281	0.0026356\\
2.3313	0.0027027\\
2.3346	0.0027717\\
2.3378	0.0028427\\
2.341	0.0029158\\
2.3443	0.002991\\
2.3475	0.0030685\\
2.3508	0.0031482\\
2.354	0.0032303\\
2.3572	0.0033149\\
2.3605	0.003402\\
2.3637	0.0034917\\
2.367	0.0035841\\
2.3702	0.0036793\\
2.3734	0.0037774\\
2.3767	0.0038785\\
2.3799	0.0039826\\
2.3831	0.00409\\
2.3864	0.0042006\\
2.3896	0.0043147\\
2.3929	0.0044323\\
2.3961	0.0045536\\
2.3993	0.0046786\\
2.4026	0.0048075\\
2.4058	0.0049405\\
2.409	0.0050777\\
2.4123	0.0052191\\
2.4155	0.0053651\\
2.4188	0.0055156\\
2.422	0.0056709\\
2.4252	0.0058312\\
2.4285	0.0059965\\
2.4317	0.0061671\\
2.4349	0.0063432\\
2.4382	0.0065249\\
2.4414	0.0067125\\
2.4447	0.006906\\
2.4479	0.0071058\\
2.4511	0.007312\\
2.4544	0.0075249\\
2.4576	0.0077446\\
2.4608	0.0079715\\
2.4641	0.0082057\\
2.4673	0.0084475\\
2.4706	0.0086971\\
2.4738	0.0089548\\
2.477	0.0092208\\
2.4803	0.0094956\\
2.4835	0.0097792\\
2.4868	0.010072\\
2.49	0.010374\\
2.4932	0.010687\\
2.4965	0.011009\\
2.4997	0.011342\\
2.5029	0.011685\\
2.5062	0.01204\\
2.5094	0.012406\\
2.5127	0.012785\\
2.5159	0.013175\\
2.5191	0.013578\\
2.5224	0.013994\\
2.5256	0.014424\\
2.5288	0.014867\\
2.5321	0.015325\\
2.5353	0.015798\\
2.5386	0.016285\\
2.5418	0.016789\\
2.545	0.017308\\
2.5483	0.017845\\
2.5515	0.018398\\
2.5547	0.018969\\
2.558	0.019558\\
2.5612	0.020166\\
2.5645	0.020793\\
2.5677	0.02144\\
2.5709	0.022107\\
2.5742	0.022795\\
2.5774	0.023505\\
2.5806	0.024237\\
2.5839	0.024992\\
2.5871	0.02577\\
2.5904	0.026572\\
2.5936	0.0274\\
2.5968	0.028252\\
2.6001	0.029131\\
2.6033	0.030037\\
2.6066	0.03097\\
2.6098	0.031931\\
2.613	0.032922\\
2.6163	0.033942\\
2.6195	0.034993\\
2.6227	0.036076\\
2.626	0.03719\\
2.6292	0.038338\\
2.6325	0.039519\\
2.6357	0.040735\\
2.6389	0.041987\\
2.6422	0.043275\\
2.6454	0.0446\\
2.6486	0.045963\\
2.6519	0.047366\\
2.6551	0.048808\\
2.6584	0.050291\\
2.6616	0.051816\\
2.6648	0.053384\\
2.6681	0.054996\\
2.6713	0.056652\\
2.6745	0.058355\\
2.6778	0.060104\\
2.681	0.0619\\
2.6843	0.063746\\
2.6875	0.065641\\
2.6907	0.067587\\
2.694	0.069585\\
2.6972	0.071636\\
2.7004	0.073741\\
2.7037	0.075901\\
2.7069	0.078118\\
2.7102	0.080392\\
2.7134	0.082724\\
2.7166	0.085115\\
2.7199	0.087567\\
2.7231	0.090081\\
2.7263	0.092658\\
2.7296	0.095299\\
2.7328	0.098005\\
2.7361	0.10078\\
2.7393	0.10362\\
2.7425	0.10652\\
2.7458	0.1095\\
2.749	0.11255\\
2.7523	0.11567\\
2.7555	0.11887\\
2.7587	0.12213\\
2.762	0.12548\\
2.7652	0.1289\\
2.7684	0.13239\\
2.7717	0.13597\\
2.7749	0.13963\\
2.7782	0.14336\\
2.7814	0.14718\\
2.7846	0.15108\\
2.7879	0.15507\\
2.7911	0.15914\\
2.7943	0.1633\\
2.7976	0.16754\\
2.8008	0.17188\\
2.8041	0.1763\\
2.8073	0.18081\\
2.8105	0.18542\\
2.8138	0.19011\\
2.817	0.1949\\
2.8202	0.19978\\
2.8235	0.20476\\
2.8267	0.20984\\
2.83	0.21501\\
2.8332	0.22027\\
2.8364	0.22564\\
2.8397	0.2311\\
2.8429	0.23667\\
2.8461	0.24233\\
2.8494	0.2481\\
2.8526	0.25396\\
2.8559	0.25993\\
2.8591	0.266\\
2.8623	0.27218\\
2.8656	0.27845\\
2.8688	0.28483\\
2.8721	0.29132\\
2.8753	0.29791\\
2.8785	0.3046\\
2.8818	0.3114\\
2.885	0.3183\\
2.8882	0.32531\\
2.8915	0.33242\\
2.8947	0.33964\\
2.898	0.34697\\
2.9012	0.35439\\
2.9044	0.36193\\
2.9077	0.36956\\
2.9109	0.37731\\
2.9141	0.38515\\
2.9174	0.3931\\
2.9206	0.40116\\
2.9239	0.40931\\
2.9271	0.41757\\
2.9303	0.42593\\
2.9336	0.43438\\
2.9368	0.44294\\
2.94	0.4516\\
2.9433	0.46036\\
2.9465	0.46921\\
2.9498	0.47816\\
2.953	0.4872\\
2.9562	0.49633\\
2.9595	0.50556\\
2.9627	0.51488\\
2.9659	0.52429\\
2.9692	0.53379\\
2.9724	0.54337\\
2.9757	0.55304\\
2.9789	0.56279\\
2.9821	0.57262\\
2.9854	0.58253\\
2.9886	0.59252\\
2.9919	0.60259\\
2.9951	0.61273\\
2.9983	0.62293\\
3.0016	0.63321\\
3.0048	0.64356\\
3.008	0.65397\\
3.0113	0.66444\\
3.0145	0.67497\\
3.0178	0.68555\\
3.021	0.69619\\
3.0242	0.70688\\
3.0275	0.71762\\
3.0307	0.72841\\
3.0339	0.73923\\
3.0372	0.7501\\
3.0404	0.761\\
3.0437	0.77193\\
3.0469	0.78289\\
3.0501	0.79388\\
3.0534	0.80489\\
3.0566	0.81592\\
3.0598	0.82696\\
3.0631	0.83801\\
3.0663	0.84907\\
3.0696	0.86012\\
3.0728	0.87118\\
3.076	0.88223\\
3.0793	0.89328\\
3.0825	0.9043\\
3.0857	0.91531\\
3.089	0.92629\\
3.0922	0.93725\\
3.0955	0.94817\\
3.0987	0.95906\\
3.1019	0.9699\\
3.1052	0.9807\\
3.1084	0.99145\\
3.1117	1.0021\\
3.1149	1.0128\\
3.1181	1.0233\\
3.1214	1.0338\\
3.1246	1.0442\\
3.1278	1.0546\\
3.1311	1.0648\\
3.1343	1.075\\
3.1376	1.0851\\
3.1408	1.095\\
3.144	1.1049\\
3.1473	1.1146\\
3.1505	1.1243\\
3.1537	1.1338\\
3.157	1.1432\\
3.1602	1.1525\\
3.1635	1.1616\\
3.1667	1.1706\\
3.1699	1.1795\\
3.1732	1.1882\\
3.1764	1.1967\\
3.1796	1.2051\\
3.1829	1.2134\\
3.1861	1.2215\\
3.1894	1.2294\\
3.1926	1.2371\\
3.1958	1.2447\\
3.1991	1.2521\\
3.2023	1.2593\\
3.2055	1.2663\\
3.2088	1.2731\\
3.212	1.2798\\
3.2153	1.2862\\
3.2185	1.2924\\
3.2217	1.2985\\
3.225	1.3043\\
3.2282	1.3099\\
3.2315	1.3153\\
3.2347	1.3205\\
3.2379	1.3255\\
3.2412	1.3303\\
3.2444	1.3348\\
3.2476	1.3391\\
3.2509	1.3432\\
3.2541	1.347\\
3.2574	1.3507\\
3.2606	1.3541\\
3.2638	1.3572\\
3.2671	1.3602\\
3.2703	1.3629\\
3.2735	1.3653\\
3.2768	1.3675\\
3.28	1.3695\\
3.2833	1.3713\\
3.2865	1.3728\\
3.2897	1.3741\\
3.293	1.3751\\
3.2962	1.3759\\
3.2994	1.3765\\
3.3027	1.3768\\
3.3059	1.3769\\
3.3092	1.3767\\
3.3124	1.3763\\
3.3156	1.3757\\
3.3189	1.3748\\
3.3221	1.3738\\
3.3253	1.3724\\
3.3286	1.3709\\
3.3318	1.3691\\
3.3351	1.3671\\
3.3383	1.3649\\
3.3415	1.3624\\
3.3448	1.3598\\
3.348	1.3569\\
3.3513	1.3538\\
3.3545	1.3505\\
3.3577	1.347\\
3.361	1.3433\\
3.3642	1.3393\\
3.3674	1.3352\\
3.3707	1.3309\\
3.3739	1.3264\\
3.3772	1.3217\\
3.3804	1.3168\\
3.3836	1.3117\\
3.3869	1.3064\\
3.3901	1.301\\
3.3933	1.2954\\
3.3966	1.2896\\
3.3998	1.2836\\
3.4031	1.2775\\
3.4063	1.2712\\
3.4095	1.2648\\
3.4128	1.2582\\
3.416	1.2514\\
3.4192	1.2445\\
3.4225	1.2375\\
3.4257	1.2303\\
3.429	1.223\\
3.4322	1.2156\\
3.4354	1.208\\
3.4387	1.2003\\
3.4419	1.1925\\
3.4451	1.1846\\
3.4484	1.1765\\
3.4516	1.1684\\
3.4549	1.1601\\
3.4581	1.1517\\
3.4613	1.1433\\
3.4646	1.1347\\
3.4678	1.1261\\
3.471	1.1173\\
3.4743	1.1085\\
3.4775	1.0996\\
3.4808	1.0907\\
3.484	1.0816\\
3.4872	1.0725\\
3.4905	1.0634\\
3.4937	1.0541\\
3.497	1.0448\\
3.5002	1.0355\\
3.5034	1.0261\\
3.5067	1.0167\\
3.5099	1.0073\\
3.5131	0.99778\\
3.5164	0.98826\\
3.5196	0.97871\\
3.5229	0.96913\\
3.5261	0.95954\\
3.5293	0.94993\\
3.5326	0.94031\\
3.5358	0.93068\\
3.539	0.92104\\
3.5423	0.91141\\
3.5455	0.90177\\
3.5488	0.89214\\
3.552	0.88252\\
3.5552	0.87292\\
3.5585	0.86333\\
3.5617	0.85375\\
3.5649	0.8442\\
3.5682	0.83468\\
3.5714	0.82518\\
3.5747	0.81572\\
3.5779	0.80628\\
3.5811	0.79689\\
3.5844	0.78754\\
3.5876	0.77823\\
3.5908	0.76896\\
3.5941	0.75975\\
3.5973	0.75058\\
3.6006	0.74147\\
3.6038	0.73242\\
3.607	0.72342\\
3.6103	0.71449\\
3.6135	0.70562\\
3.6168	0.69681\\
3.62	0.68807\\
3.6232	0.6794\\
3.6265	0.6708\\
3.6297	0.66228\\
3.6329	0.65383\\
3.6362	0.64546\\
3.6394	0.63716\\
3.6427	0.62895\\
3.6459	0.62082\\
3.6491	0.61278\\
3.6524	0.60481\\
3.6556	0.59694\\
3.6588	0.58915\\
3.6621	0.58146\\
3.6653	0.57385\\
3.6686	0.56634\\
3.6718	0.55891\\
3.675	0.55159\\
3.6783	0.54435\\
3.6815	0.53722\\
3.6847	0.53018\\
3.688	0.52323\\
3.6912	0.51639\\
3.6945	0.50964\\
3.6977	0.50299\\
3.7009	0.49645\\
3.7042	0.49\\
3.7074	0.48366\\
3.7106	0.47741\\
3.7139	0.47127\\
3.7171	0.46523\\
3.7204	0.45929\\
3.7236	0.45346\\
3.7268	0.44772\\
3.7301	0.44209\\
3.7333	0.43657\\
3.7366	0.43114\\
3.7398	0.42582\\
3.743	0.42059\\
3.7463	0.41547\\
3.7495	0.41045\\
3.7527	0.40553\\
3.756	0.40071\\
3.7592	0.39599\\
3.7625	0.39137\\
3.7657	0.38685\\
3.7689	0.38242\\
3.7722	0.37809\\
3.7754	0.37385\\
3.7786	0.36971\\
3.7819	0.36566\\
3.7851	0.3617\\
3.7884	0.35783\\
3.7916	0.35405\\
3.7948	0.35036\\
3.7981	0.34675\\
3.8013	0.34323\\
3.8045	0.33979\\
3.8078	0.33644\\
3.811	0.33317\\
3.8143	0.32997\\
3.8175	0.32686\\
3.8207	0.32382\\
3.824	0.32086\\
3.8272	0.31797\\
3.8304	0.31515\\
3.8337	0.31241\\
3.8369	0.30974\\
3.8402	0.30713\\
3.8434	0.30459\\
3.8466	0.30212\\
3.8499	0.29971\\
3.8531	0.29737\\
3.8564	0.29509\\
3.8596	0.29286\\
3.8628	0.2907\\
3.8661	0.28859\\
3.8693	0.28654\\
3.8725	0.28455\\
3.8758	0.28261\\
3.879	0.28072\\
3.8823	0.27889\\
3.8855	0.2771\\
3.8887	0.27536\\
3.892	0.27367\\
3.8952	0.27203\\
3.8984	0.27043\\
3.9017	0.26888\\
3.9049	0.26736\\
3.9082	0.26589\\
3.9114	0.26446\\
3.9146	0.26307\\
3.9179	0.26172\\
3.9211	0.2604\\
3.9243	0.25912\\
3.9276	0.25788\\
3.9308	0.25666\\
3.9341	0.25548\\
3.9373	0.25434\\
3.9405	0.25322\\
3.9438	0.25213\\
3.947	0.25107\\
3.9502	0.25004\\
3.9535	0.24903\\
3.9567	0.24805\\
3.96	0.24709\\
3.9632	0.24616\\
3.9664	0.24525\\
3.9697	0.24436\\
3.9729	0.24349\\
3.9762	0.24264\\
3.9794	0.24181\\
3.9826	0.24099\\
3.9859	0.2402\\
3.9891	0.23942\\
3.9923	0.23865\\
3.9956	0.23791\\
3.9988	0.23717\\
4.0021	0.23645\\
4.0053	0.23574\\
4.0085	0.23504\\
4.0118	0.23435\\
4.015	0.23367\\
4.0182	0.233\\
4.0215	0.23234\\
4.0247	0.23169\\
4.028	0.23105\\
4.0312	0.23041\\
4.0344	0.22978\\
4.0377	0.22915\\
4.0409	0.22853\\
4.0441	0.22791\\
4.0474	0.2273\\
4.0506	0.22669\\
4.0539	0.22608\\
4.0571	0.22548\\
4.0603	0.22487\\
4.0636	0.22427\\
4.0668	0.22367\\
4.07	0.22307\\
4.0733	0.22247\\
4.0765	0.22186\\
4.0798	0.22126\\
4.083	0.22065\\
4.0862	0.22005\\
4.0895	0.21944\\
4.0927	0.21883\\
4.0959	0.21821\\
4.0992	0.21759\\
4.1024	0.21697\\
4.1057	0.21634\\
4.1089	0.21571\\
4.1121	0.21508\\
4.1154	0.21444\\
4.1186	0.21379\\
4.1219	0.21314\\
4.1251	0.21249\\
4.1283	0.21182\\
4.1316	0.21115\\
4.1348	0.21048\\
4.138	0.20979\\
4.1413	0.2091\\
4.1445	0.2084\\
4.1478	0.2077\\
4.151	0.20699\\
4.1542	0.20626\\
4.1575	0.20553\\
4.1607	0.20479\\
4.1639	0.20405\\
4.1672	0.20329\\
4.1704	0.20253\\
4.1737	0.20175\\
4.1769	0.20097\\
4.1801	0.20017\\
4.1834	0.19937\\
4.1866	0.19855\\
4.1898	0.19773\\
4.1931	0.19689\\
4.1963	0.19605\\
4.1996	0.19519\\
4.2028	0.19433\\
4.206	0.19345\\
4.2093	0.19256\\
4.2125	0.19166\\
4.2157	0.19075\\
4.219	0.18983\\
4.2222	0.1889\\
4.2255	0.18796\\
4.2287	0.187\\
4.2319	0.18604\\
4.2352	0.18506\\
4.2384	0.18407\\
4.2417	0.18307\\
4.2449	0.18206\\
4.2481	0.18104\\
4.2514	0.18001\\
4.2546	0.17896\\
4.2578	0.17791\\
4.2611	0.17684\\
4.2643	0.17576\\
4.2676	0.17467\\
4.2708	0.17358\\
4.274	0.17247\\
4.2773	0.17135\\
4.2805	0.17021\\
4.2837	0.16907\\
4.287	0.16792\\
4.2902	0.16676\\
4.2935	0.16559\\
4.2967	0.1644\\
4.2999	0.16321\\
4.3032	0.16201\\
4.3064	0.1608\\
4.3096	0.15958\\
4.3129	0.15835\\
4.3161	0.15711\\
4.3194	0.15586\\
4.3226	0.1546\\
4.3258	0.15333\\
4.3291	0.15206\\
4.3323	0.15078\\
4.3355	0.14949\\
4.3388	0.14819\\
4.342	0.14688\\
4.3453	0.14557\\
4.3485	0.14425\\
4.3517	0.14292\\
4.355	0.14158\\
4.3582	0.14024\\
4.3615	0.1389\\
4.3647	0.13754\\
4.3679	0.13618\\
4.3712	0.13482\\
4.3744	0.13345\\
4.3776	0.13208\\
4.3809	0.1307\\
4.3841	0.12932\\
4.3874	0.12793\\
4.3906	0.12654\\
4.3938	0.12515\\
4.3971	0.12376\\
4.4003	0.12236\\
4.4035	0.12096\\
4.4068	0.11956\\
4.41	0.11815\\
4.4133	0.11675\\
4.4165	0.11535\\
4.4197	0.11394\\
4.423	0.11254\\
4.4262	0.11114\\
4.4294	0.10973\\
4.4327	0.10833\\
4.4359	0.10693\\
4.4392	0.10554\\
4.4424	0.10414\\
4.4456	0.10275\\
4.4489	0.10136\\
4.4521	0.099978\\
4.4553	0.098598\\
4.4586	0.097222\\
4.4618	0.095852\\
4.4651	0.094486\\
4.4683	0.093126\\
4.4715	0.091773\\
4.4748	0.090425\\
4.478	0.089085\\
4.4813	0.087751\\
4.4845	0.086425\\
4.4877	0.085107\\
4.491	0.083798\\
4.4942	0.082496\\
4.4974	0.081204\\
4.5007	0.07992\\
4.5039	0.078646\\
4.5072	0.077382\\
4.5104	0.076128\\
4.5136	0.074885\\
4.5169	0.073652\\
4.5201	0.07243\\
4.5233	0.071219\\
4.5266	0.07002\\
4.5298	0.068832\\
4.5331	0.067656\\
4.5363	0.066493\\
4.5395	0.065342\\
4.5428	0.064203\\
4.546	0.063078\\
4.5492	0.061965\\
4.5525	0.060866\\
4.5557	0.059781\\
4.559	0.058709\\
4.5622	0.057651\\
4.5654	0.056607\\
4.5687	0.055578\\
4.5719	0.054563\\
4.5751	0.053563\\
4.5784	0.052578\\
4.5816	0.051607\\
4.5849	0.050652\\
4.5881	0.049712\\
4.5913	0.048787\\
4.5946	0.047877\\
4.5978	0.046982\\
4.6011	0.046103\\
4.6043	0.045238\\
4.6075	0.044388\\
4.6108	0.043552\\
4.614	0.042731\\
4.6172	0.041923\\
4.6205	0.041129\\
4.6237	0.040349\\
4.627	0.039581\\
4.6302	0.038826\\
4.6334	0.038084\\
4.6367	0.037354\\
4.6399	0.036636\\
4.6431	0.03593\\
4.6464	0.035235\\
4.6496	0.034552\\
4.6529	0.03388\\
4.6561	0.033219\\
4.6593	0.032569\\
4.6626	0.03193\\
4.6658	0.031301\\
4.669	0.030682\\
4.6723	0.030074\\
4.6755	0.029476\\
4.6788	0.028888\\
4.682	0.02831\\
4.6852	0.027742\\
4.6885	0.027183\\
4.6917	0.026634\\
4.6949	0.026095\\
4.6982	0.025564\\
4.7014	0.025043\\
4.7047	0.024531\\
4.7079	0.024028\\
4.7111	0.023534\\
4.7144	0.023049\\
4.7176	0.022572\\
4.7209	0.022104\\
4.7241	0.021644\\
4.7273	0.021193\\
4.7306	0.02075\\
4.7338	0.020314\\
4.737	0.019887\\
4.7403	0.019468\\
4.7435	0.019057\\
4.7468	0.018653\\
4.75	0.018257\\
4.7532	0.017869\\
4.7565	0.017487\\
4.7597	0.017113\\
4.7629	0.016746\\
4.7662	0.016387\\
4.7694	0.016034\\
4.7727	0.015688\\
4.7759	0.015349\\
4.7791	0.015016\\
4.7824	0.01469\\
4.7856	0.01437\\
4.7888	0.014057\\
4.7921	0.01375\\
4.7953	0.013449\\
4.7986	0.013154\\
4.8018	0.012865\\
4.805	0.012582\\
4.8083	0.012304\\
4.8115	0.012032\\
4.8147	0.011766\\
4.818	0.011505\\
4.8212	0.01125\\
4.8245	0.011\\
4.8277	0.010755\\
4.8309	0.010515\\
4.8342	0.01028\\
4.8374	0.01005\\
4.8406	0.0098245\\
4.8439	0.009604\\
4.8471	0.0093882\\
4.8504	0.0091769\\
4.8536	0.00897\\
4.8568	0.0087676\\
4.8601	0.0085695\\
4.8633	0.0083756\\
4.8666	0.0081858\\
4.8698	0.0080001\\
4.873	0.0078185\\
4.8763	0.0076407\\
4.8795	0.0074668\\
4.8827	0.0072966\\
4.886	0.0071302\\
4.8892	0.0069673\\
4.8925	0.006808\\
4.8957	0.0066522\\
4.8989	0.0064999\\
4.9022	0.0063508\\
4.9054	0.006205\\
4.9086	0.0060625\\
4.9119	0.0059231\\
4.9151	0.0057868\\
4.9184	0.0056535\\
4.9216	0.0055231\\
4.9248	0.0053957\\
4.9281	0.0052711\\
4.9313	0.0051492\\
4.9345	0.0050301\\
4.9378	0.0049137\\
4.941	0.0047999\\
4.9443	0.0046886\\
4.9475	0.0045799\\
4.9507	0.0044735\\
4.954	0.0043696\\
4.9572	0.0042681\\
4.9604	0.0041688\\
4.9637	0.0040718\\
4.9669	0.0039769\\
4.9702	0.0038842\\
4.9734	0.0037937\\
4.9766	0.0037052\\
4.9799	0.0036187\\
4.9831	0.0035342\\
4.9864	0.0034516\\
4.9896	0.0033709\\
4.9928	0.003292\\
4.9961	0.0032149\\
4.9993	0.0031397\\
5.0025	0.0030661\\
5.0058	0.0029942\\
5.009	0.002924\\
5.0123	0.0028554\\
5.0155	0.0027884\\
5.0187	0.0027229\\
5.022	0.0026589\\
5.0252	0.0025964\\
5.0284	0.0025353\\
5.0317	0.0024757\\
5.0349	0.0024174\\
5.0382	0.0023605\\
5.0414	0.0023049\\
5.0446	0.0022506\\
5.0479	0.0021975\\
5.0511	0.0021457\\
5.0543	0.002095\\
5.0576	0.0020456\\
5.0608	0.0019973\\
5.0641	0.0019501\\
5.0673	0.001904\\
5.0705	0.001859\\
5.0738	0.001815\\
5.077	0.001772\\
5.0802	0.0017301\\
5.0835	0.0016891\\
5.0867	0.0016491\\
5.09	0.00161\\
5.0932	0.0015718\\
5.0964	0.0015345\\
5.0997	0.0014981\\
5.1029	0.0014626\\
5.1062	0.0014278\\
5.1094	0.0013939\\
5.1126	0.0013608\\
5.1159	0.0013284\\
5.1191	0.0012968\\
5.1223	0.0012659\\
5.1256	0.0012358\\
5.1288	0.0012063\\
5.1321	0.0011776\\
5.1353	0.0011495\\
5.1385	0.0011221\\
5.1418	0.0010953\\
5.145	0.0010692\\
5.1482	0.0010436\\
5.1515	0.0010187\\
5.1547	0.00099432\\
5.158	0.00097053\\
5.1612	0.0009473\\
5.1644	0.00092462\\
5.1677	0.00090247\\
5.1709	0.00088083\\
5.1741	0.00085971\\
5.1774	0.00083908\\
5.1806	0.00081893\\
5.1839	0.00079926\\
5.1871	0.00078005\\
5.1903	0.00076129\\
5.1936	0.00074297\\
5.1968	0.00072508\\
5.2	0.00070761\\
5.2033	0.00069056\\
5.2065	0.0006739\\
5.2098	0.00065764\\
5.213	0.00064176\\
5.2162	0.00062625\\
5.2195	0.00061111\\
5.2227	0.00059632\\
5.226	0.00058189\\
5.2292	0.00056779\\
5.2324	0.00055403\\
5.2357	0.00054059\\
5.2389	0.00052747\\
5.2421	0.00051466\\
5.2454	0.00050215\\
5.2486	0.00048993\\
5.2519	0.00047801\\
5.2551	0.00046636\\
5.2583	0.000455\\
5.2616	0.0004439\\
5.2648	0.00043306\\
5.268	0.00042248\\
5.2713	0.00041215\\
5.2745	0.00040207\\
5.2778	0.00039222\\
5.281	0.00038261\\
5.2842	0.00037323\\
5.2875	0.00036407\\
5.2907	0.00035513\\
5.2939	0.0003464\\
5.2972	0.00033787\\
5.3004	0.00032955\\
5.3037	0.00032143\\
5.3069	0.0003135\\
5.3101	0.00030577\\
5.3101	0.00016846\\
5.3069	0.00017311\\
5.3037	0.00017789\\
5.3004	0.0001828\\
5.2972	0.00018784\\
5.2939	0.00019302\\
5.2907	0.00019834\\
5.2875	0.00020381\\
5.2842	0.00020942\\
5.281	0.00021519\\
5.2778	0.00022112\\
5.2745	0.0002272\\
5.2713	0.00023346\\
5.268	0.00023988\\
5.2648	0.00024648\\
5.2616	0.00025326\\
5.2583	0.00026022\\
5.2551	0.00026738\\
5.2519	0.00027473\\
5.2486	0.00028228\\
5.2454	0.00029003\\
5.2421	0.000298\\
5.2389	0.00030619\\
5.2357	0.0003146\\
5.2324	0.00032324\\
5.2292	0.00033211\\
5.226	0.00034123\\
5.2227	0.0003506\\
5.2195	0.00036022\\
5.2162	0.00037011\\
5.213	0.00038027\\
5.2098	0.00039071\\
5.2065	0.00040143\\
5.2033	0.00041244\\
5.2	0.00042376\\
5.1968	0.00043538\\
5.1936	0.00044733\\
5.1903	0.0004596\\
5.1871	0.00047221\\
5.1839	0.00048516\\
5.1806	0.00049847\\
5.1774	0.00051214\\
5.1741	0.00052619\\
5.1709	0.00054062\\
5.1677	0.00055544\\
5.1644	0.00057068\\
5.1612	0.00058633\\
5.158	0.00060241\\
5.1547	0.00061893\\
5.1515	0.0006359\\
5.1482	0.00065334\\
5.145	0.00067125\\
5.1418	0.00068966\\
5.1385	0.00070857\\
5.1353	0.000728\\
5.1321	0.00074796\\
5.1288	0.00076847\\
5.1256	0.00078954\\
5.1223	0.00081119\\
5.1191	0.00083343\\
5.1159	0.00085628\\
5.1126	0.00087976\\
5.1094	0.00090388\\
5.1062	0.00092866\\
5.1029	0.00095411\\
5.0997	0.00098027\\
5.0964	0.0010071\\
5.0932	0.0010347\\
5.09	0.0010631\\
5.0867	0.0010922\\
5.0835	0.0011222\\
5.0802	0.0011529\\
5.077	0.0011845\\
5.0738	0.0012169\\
5.0705	0.0012503\\
5.0673	0.0012845\\
5.0641	0.0013197\\
5.0608	0.0013558\\
5.0576	0.0013929\\
5.0543	0.001431\\
5.0511	0.0014702\\
5.0479	0.0015104\\
5.0446	0.0015517\\
5.0414	0.0015942\\
5.0382	0.0016378\\
5.0349	0.0016825\\
5.0317	0.0017285\\
5.0284	0.0017757\\
5.0252	0.0018242\\
5.022	0.001874\\
5.0187	0.0019252\\
5.0155	0.0019777\\
5.0123	0.0020317\\
5.009	0.002087\\
5.0058	0.0021439\\
5.0025	0.0022023\\
4.9993	0.0022623\\
4.9961	0.0023239\\
4.9928	0.0023871\\
4.9896	0.002452\\
4.9864	0.0025186\\
4.9831	0.0025871\\
4.9799	0.0026573\\
4.9766	0.0027294\\
4.9734	0.0028034\\
4.9702	0.0028794\\
4.9669	0.0029573\\
4.9637	0.0030374\\
4.9604	0.0031195\\
4.9572	0.0032038\\
4.954	0.0032903\\
4.9507	0.0033791\\
4.9475	0.0034702\\
4.9443	0.0035637\\
4.941	0.0036596\\
4.9378	0.003758\\
4.9345	0.003859\\
4.9313	0.0039626\\
4.9281	0.0040689\\
4.9248	0.0041779\\
4.9216	0.0042897\\
4.9184	0.0044044\\
4.9151	0.004522\\
4.9119	0.0046427\\
4.9086	0.0047664\\
4.9054	0.0048932\\
4.9022	0.0050233\\
4.8989	0.0051567\\
4.8957	0.0052934\\
4.8925	0.0054336\\
4.8892	0.0055774\\
4.886	0.0057247\\
4.8827	0.0058757\\
4.8795	0.0060304\\
4.8763	0.006189\\
4.873	0.0063516\\
4.8698	0.0065181\\
4.8666	0.0066888\\
4.8633	0.0068636\\
4.8601	0.0070428\\
4.8568	0.0072263\\
4.8536	0.0074142\\
4.8504	0.0076067\\
4.8471	0.0078039\\
4.8439	0.0080058\\
4.8406	0.0082126\\
4.8374	0.0084243\\
4.8342	0.008641\\
4.8309	0.0088629\\
4.8277	0.0090901\\
4.8245	0.0093226\\
4.8212	0.0095605\\
4.818	0.0098041\\
4.8147	0.010053\\
4.8115	0.010308\\
4.8083	0.010569\\
4.805	0.010836\\
4.8018	0.011109\\
4.7986	0.011388\\
4.7953	0.011674\\
4.7921	0.011966\\
4.7888	0.012264\\
4.7856	0.01257\\
4.7824	0.012882\\
4.7791	0.013201\\
4.7759	0.013527\\
4.7727	0.01386\\
4.7694	0.014201\\
4.7662	0.014549\\
4.7629	0.014905\\
4.7597	0.015268\\
4.7565	0.015639\\
4.7532	0.016018\\
4.75	0.016405\\
4.7468	0.016801\\
4.7435	0.017204\\
4.7403	0.017616\\
4.737	0.018037\\
4.7338	0.018466\\
4.7306	0.018904\\
4.7273	0.019352\\
4.7241	0.019808\\
4.7209	0.020273\\
4.7176	0.020748\\
4.7144	0.021232\\
4.7111	0.021726\\
4.7079	0.022229\\
4.7047	0.022743\\
4.7014	0.023266\\
4.6982	0.023799\\
4.6949	0.024343\\
4.6917	0.024897\\
4.6885	0.025461\\
4.6852	0.026035\\
4.682	0.026621\\
4.6788	0.027217\\
4.6755	0.027824\\
4.6723	0.028441\\
4.669	0.02907\\
4.6658	0.029709\\
4.6626	0.03036\\
4.6593	0.031022\\
4.6561	0.031695\\
4.6529	0.032379\\
4.6496	0.033074\\
4.6464	0.033781\\
4.6431	0.034498\\
4.6399	0.035227\\
4.6367	0.035967\\
4.6334	0.036717\\
4.6302	0.037479\\
4.627	0.038251\\
4.6237	0.039034\\
4.6205	0.039827\\
4.6172	0.040631\\
4.614	0.041444\\
4.6108	0.042267\\
4.6075	0.0431\\
4.6043	0.043942\\
4.6011	0.044794\\
4.5978	0.045654\\
4.5946	0.046523\\
4.5913	0.047401\\
4.5881	0.048288\\
4.5849	0.049184\\
4.5816	0.050088\\
4.5784	0.051002\\
4.5751	0.051924\\
4.5719	0.052855\\
4.5687	0.053795\\
4.5654	0.054744\\
4.5622	0.055702\\
4.559	0.05667\\
4.5557	0.057646\\
4.5525	0.058632\\
4.5492	0.059628\\
4.546	0.060633\\
4.5428	0.061647\\
4.5395	0.062671\\
4.5363	0.063704\\
4.5331	0.064746\\
4.5298	0.065798\\
4.5266	0.066859\\
4.5233	0.06793\\
4.5201	0.06901\\
4.5169	0.070098\\
4.5136	0.071197\\
4.5104	0.072304\\
4.5072	0.07342\\
4.5039	0.074545\\
4.5007	0.075678\\
4.4974	0.07682\\
4.4942	0.077971\\
4.491	0.07913\\
4.4877	0.080298\\
4.4845	0.081474\\
4.4813	0.082657\\
4.478	0.083849\\
4.4748	0.085048\\
4.4715	0.086254\\
4.4683	0.087468\\
4.4651	0.08869\\
4.4618	0.089918\\
4.4586	0.091153\\
4.4553	0.092394\\
4.4521	0.093642\\
4.4489	0.094897\\
4.4456	0.096157\\
4.4424	0.097423\\
4.4392	0.098694\\
4.4359	0.099971\\
4.4327	0.10125\\
4.4294	0.10254\\
4.4262	0.10383\\
4.423	0.10513\\
4.4197	0.10643\\
4.4165	0.10773\\
4.4133	0.10904\\
4.41	0.11035\\
4.4068	0.11166\\
4.4035	0.11297\\
4.4003	0.11429\\
4.3971	0.11561\\
4.3938	0.11693\\
4.3906	0.11826\\
4.3874	0.11958\\
4.3841	0.1209\\
4.3809	0.12223\\
4.3776	0.12355\\
4.3744	0.12488\\
4.3712	0.1262\\
4.3679	0.12753\\
4.3647	0.12885\\
4.3615	0.13017\\
4.3582	0.13149\\
4.355	0.1328\\
4.3517	0.13412\\
4.3485	0.13543\\
4.3453	0.13674\\
4.342	0.13804\\
4.3388	0.13934\\
4.3355	0.14063\\
4.3323	0.14192\\
4.3291	0.14321\\
4.3258	0.14449\\
4.3226	0.14576\\
4.3194	0.14703\\
4.3161	0.14829\\
4.3129	0.14955\\
4.3096	0.1508\\
4.3064	0.15204\\
4.3032	0.15327\\
4.2999	0.15449\\
4.2967	0.15571\\
4.2935	0.15692\\
4.2902	0.15812\\
4.287	0.15931\\
4.2837	0.16049\\
4.2805	0.16166\\
4.2773	0.16282\\
4.274	0.16397\\
4.2708	0.16511\\
4.2676	0.16624\\
4.2643	0.16735\\
4.2611	0.16846\\
4.2578	0.16956\\
4.2546	0.17064\\
4.2514	0.17171\\
4.2481	0.17277\\
4.2449	0.17382\\
4.2417	0.17486\\
4.2384	0.17588\\
4.2352	0.17689\\
4.2319	0.17789\\
4.2287	0.17887\\
4.2255	0.17985\\
4.2222	0.1808\\
4.219	0.18175\\
4.2157	0.18268\\
4.2125	0.1836\\
4.2093	0.18451\\
4.206	0.1854\\
4.2028	0.18628\\
4.1996	0.18714\\
4.1963	0.18799\\
4.1931	0.18883\\
4.1898	0.18965\\
4.1866	0.19046\\
4.1834	0.19126\\
4.1801	0.19204\\
4.1769	0.19281\\
4.1737	0.19356\\
4.1704	0.19431\\
4.1672	0.19503\\
4.1639	0.19575\\
4.1607	0.19645\\
4.1575	0.19714\\
4.1542	0.19781\\
4.151	0.19847\\
4.1478	0.19912\\
4.1445	0.19976\\
4.1413	0.20039\\
4.138	0.201\\
4.1348	0.2016\\
4.1316	0.20219\\
4.1283	0.20277\\
4.1251	0.20334\\
4.1219	0.20389\\
4.1186	0.20444\\
4.1154	0.20498\\
4.1121	0.20551\\
4.1089	0.20603\\
4.1057	0.20654\\
4.1024	0.20704\\
4.0992	0.20753\\
4.0959	0.20802\\
4.0927	0.2085\\
4.0895	0.20898\\
4.0862	0.20945\\
4.083	0.20992\\
4.0798	0.21038\\
4.0765	0.21084\\
4.0733	0.21129\\
4.07	0.21174\\
4.0668	0.21219\\
4.0636	0.21264\\
4.0603	0.21309\\
4.0571	0.21354\\
4.0539	0.21399\\
4.0506	0.21444\\
4.0474	0.2149\\
4.0441	0.21536\\
4.0409	0.21582\\
4.0377	0.21629\\
4.0344	0.21676\\
4.0312	0.21724\\
4.028	0.21773\\
4.0247	0.21822\\
4.0215	0.21872\\
4.0182	0.21924\\
4.015	0.21976\\
4.0118	0.2203\\
4.0085	0.22085\\
4.0053	0.22141\\
4.0021	0.22199\\
3.9988	0.22258\\
3.9956	0.22319\\
3.9923	0.22381\\
3.9891	0.22446\\
3.9859	0.22512\\
3.9826	0.2258\\
3.9794	0.22651\\
3.9762	0.22723\\
3.9729	0.22798\\
3.9697	0.22876\\
3.9664	0.22955\\
3.9632	0.23038\\
3.96	0.23123\\
3.9567	0.23211\\
3.9535	0.23303\\
3.9502	0.23397\\
3.947	0.23494\\
3.9438	0.23595\\
3.9405	0.23699\\
3.9373	0.23806\\
3.9341	0.23917\\
3.9308	0.24032\\
3.9276	0.24151\\
3.9243	0.24273\\
3.9211	0.244\\
3.9179	0.24531\\
3.9146	0.24666\\
3.9114	0.24805\\
3.9082	0.24949\\
3.9049	0.25098\\
3.9017	0.25251\\
3.8984	0.25409\\
3.8952	0.25572\\
3.892	0.2574\\
3.8887	0.25914\\
3.8855	0.26092\\
3.8823	0.26276\\
3.879	0.26466\\
3.8758	0.26661\\
3.8725	0.26862\\
3.8693	0.27069\\
3.8661	0.27282\\
3.8628	0.275\\
3.8596	0.27725\\
3.8564	0.27956\\
3.8531	0.28194\\
3.8499	0.28438\\
3.8466	0.28688\\
3.8434	0.28945\\
3.8402	0.29209\\
3.8369	0.2948\\
3.8337	0.29758\\
3.8304	0.30042\\
3.8272	0.30334\\
3.824	0.30633\\
3.8207	0.30939\\
3.8175	0.31253\\
3.8143	0.31574\\
3.811	0.31902\\
3.8078	0.32238\\
3.8045	0.32582\\
3.8013	0.32933\\
3.7981	0.33292\\
3.7948	0.33659\\
3.7916	0.34034\\
3.7884	0.34417\\
3.7851	0.34808\\
3.7819	0.35207\\
3.7786	0.35614\\
3.7754	0.36029\\
3.7722	0.36452\\
3.7689	0.36884\\
3.7657	0.37324\\
3.7625	0.37772\\
3.7592	0.38229\\
3.756	0.38695\\
3.7527	0.39169\\
3.7495	0.39651\\
3.7463	0.40142\\
3.743	0.40642\\
3.7398	0.41151\\
3.7366	0.41668\\
3.7333	0.42195\\
3.7301	0.4273\\
3.7268	0.43274\\
3.7236	0.43828\\
3.7204	0.4439\\
3.7171	0.44962\\
3.7139	0.45543\\
3.7106	0.46133\\
3.7074	0.46732\\
3.7042	0.47341\\
3.7009	0.47959\\
3.6977	0.48586\\
3.6945	0.49222\\
3.6912	0.49868\\
3.688	0.50524\\
3.6847	0.51189\\
3.6815	0.51863\\
3.6783	0.52546\\
3.675	0.53239\\
3.6718	0.53941\\
3.6686	0.54652\\
3.6653	0.55373\\
3.6621	0.56103\\
3.6588	0.56842\\
3.6556	0.5759\\
3.6524	0.58347\\
3.6491	0.59113\\
3.6459	0.59888\\
3.6427	0.60671\\
3.6394	0.61464\\
3.6362	0.62264\\
3.6329	0.63074\\
3.6297	0.63891\\
3.6265	0.64717\\
3.6232	0.65551\\
3.62	0.66392\\
3.6168	0.67241\\
3.6135	0.68098\\
3.6103	0.68962\\
3.607	0.69834\\
3.6038	0.70712\\
3.6006	0.71597\\
3.5973	0.72489\\
3.5941	0.73387\\
3.5908	0.74291\\
3.5876	0.75202\\
3.5844	0.76117\\
3.5811	0.77039\\
3.5779	0.77965\\
3.5747	0.78896\\
3.5714	0.79832\\
3.5682	0.80772\\
3.5649	0.81716\\
3.5617	0.82664\\
3.5585	0.83616\\
3.5552	0.8457\\
3.552	0.85527\\
3.5488	0.86487\\
3.5455	0.87448\\
3.5423	0.88412\\
3.539	0.89376\\
3.5358	0.90342\\
3.5326	0.91309\\
3.5293	0.92275\\
3.5261	0.93242\\
3.5229	0.94208\\
3.5196	0.95173\\
3.5164	0.96137\\
3.5131	0.97099\\
3.5099	0.98059\\
3.5067	0.99017\\
3.5034	0.99971\\
3.5002	1.0092\\
3.497	1.0187\\
3.4937	1.0281\\
3.4905	1.0375\\
3.4872	1.0468\\
3.484	1.0561\\
3.4808	1.0653\\
3.4775	1.0745\\
3.4743	1.0835\\
3.471	1.0925\\
3.4678	1.1014\\
3.4646	1.1103\\
3.4613	1.119\\
3.4581	1.1277\\
3.4549	1.1362\\
3.4516	1.1447\\
3.4484	1.153\\
3.4451	1.1612\\
3.4419	1.1694\\
3.4387	1.1773\\
3.4354	1.1852\\
3.4322	1.1929\\
3.429	1.2005\\
3.4257	1.208\\
3.4225	1.2153\\
3.4192	1.2224\\
3.416	1.2294\\
3.4128	1.2363\\
3.4095	1.243\\
3.4063	1.2495\\
3.4031	1.2558\\
3.3998	1.262\\
3.3966	1.268\\
3.3933	1.2738\\
3.3901	1.2794\\
3.3869	1.2849\\
3.3836	1.2901\\
3.3804	1.2952\\
3.3772	1.3\\
3.3739	1.3047\\
3.3707	1.3091\\
3.3674	1.3134\\
3.3642	1.3174\\
3.361	1.3212\\
3.3577	1.3248\\
3.3545	1.3282\\
3.3513	1.3314\\
3.348	1.3344\\
3.3448	1.3371\\
3.3415	1.3396\\
3.3383	1.3419\\
3.3351	1.344\\
3.3318	1.3458\\
3.3286	1.3474\\
3.3253	1.3488\\
3.3221	1.3499\\
3.3189	1.3508\\
3.3156	1.3515\\
3.3124	1.352\\
3.3092	1.3522\\
3.3059	1.3522\\
3.3027	1.3519\\
3.2994	1.3515\\
3.2962	1.3508\\
3.293	1.3498\\
3.2897	1.3486\\
3.2865	1.3472\\
3.2833	1.3456\\
3.28	1.3437\\
3.2768	1.3417\\
3.2735	1.3393\\
3.2703	1.3368\\
3.2671	1.334\\
3.2638	1.331\\
3.2606	1.3278\\
3.2574	1.3243\\
3.2541	1.3207\\
3.2509	1.3168\\
3.2476	1.3127\\
3.2444	1.3084\\
3.2412	1.3038\\
3.2379	1.2991\\
3.2347	1.2942\\
3.2315	1.289\\
3.2282	1.2837\\
3.225	1.2781\\
3.2217	1.2723\\
3.2185	1.2664\\
3.2153	1.2603\\
3.212	1.2539\\
3.2088	1.2474\\
3.2055	1.2407\\
3.2023	1.2338\\
3.1991	1.2268\\
3.1958	1.2195\\
3.1926	1.2121\\
3.1894	1.2046\\
3.1861	1.1968\\
3.1829	1.1889\\
3.1796	1.1809\\
3.1764	1.1727\\
3.1732	1.1644\\
3.1699	1.1559\\
3.1667	1.1472\\
3.1635	1.1385\\
3.1602	1.1296\\
3.157	1.1205\\
3.1537	1.1114\\
3.1505	1.1021\\
3.1473	1.0927\\
3.144	1.0832\\
3.1408	1.0736\\
3.1376	1.0639\\
3.1343	1.0541\\
3.1311	1.0442\\
3.1278	1.0342\\
3.1246	1.0241\\
3.1214	1.0139\\
3.1181	1.0037\\
3.1149	0.99337\\
3.1117	0.98298\\
3.1084	0.97253\\
3.1052	0.96201\\
3.1019	0.95144\\
3.0987	0.94082\\
3.0955	0.93014\\
3.0922	0.91943\\
3.089	0.90867\\
3.0857	0.89787\\
3.0825	0.88705\\
3.0793	0.8762\\
3.076	0.86532\\
3.0728	0.85443\\
3.0696	0.84352\\
3.0663	0.83259\\
3.0631	0.82167\\
3.0598	0.81073\\
3.0566	0.7998\\
3.0534	0.78888\\
3.0501	0.77796\\
3.0469	0.76705\\
3.0437	0.75616\\
3.0404	0.7453\\
3.0372	0.73445\\
3.0339	0.72363\\
3.0307	0.71285\\
3.0275	0.7021\\
3.0242	0.69138\\
3.021	0.68071\\
3.0178	0.67009\\
3.0145	0.65951\\
3.0113	0.64898\\
3.008	0.63851\\
3.0048	0.6281\\
3.0016	0.61775\\
2.9983	0.60746\\
2.9951	0.59724\\
2.9919	0.58708\\
2.9886	0.577\\
2.9854	0.567\\
2.9821	0.55707\\
2.9789	0.54722\\
2.9757	0.53745\\
2.9724	0.52776\\
2.9692	0.51817\\
2.9659	0.50866\\
2.9627	0.49923\\
2.9595	0.48991\\
2.9562	0.48067\\
2.953	0.47153\\
2.9498	0.46249\\
2.9465	0.45354\\
2.9433	0.44469\\
2.94	0.43595\\
2.9368	0.4273\\
2.9336	0.41876\\
2.9303	0.41032\\
2.9271	0.40199\\
2.9239	0.39376\\
2.9206	0.38564\\
2.9174	0.37762\\
2.9141	0.36971\\
2.9109	0.36191\\
2.9077	0.35422\\
2.9044	0.34664\\
2.9012	0.33917\\
2.898	0.3318\\
2.8947	0.32455\\
2.8915	0.3174\\
2.8882	0.31037\\
2.885	0.30344\\
2.8818	0.29662\\
2.8785	0.28991\\
2.8753	0.28331\\
2.8721	0.27682\\
2.8688	0.27044\\
2.8656	0.26417\\
2.8623	0.258\\
2.8591	0.25194\\
2.8559	0.24598\\
2.8526	0.24014\\
2.8494	0.23439\\
2.8461	0.22875\\
2.8429	0.22322\\
2.8397	0.21779\\
2.8364	0.21246\\
2.8332	0.20723\\
2.83	0.2021\\
2.8267	0.19707\\
2.8235	0.19214\\
2.8202	0.18731\\
2.817	0.18257\\
2.8138	0.17793\\
2.8105	0.17339\\
2.8073	0.16893\\
2.8041	0.16457\\
2.8008	0.16031\\
2.7976	0.15613\\
2.7943	0.15204\\
2.7911	0.14804\\
2.7879	0.14412\\
2.7846	0.1403\\
2.7814	0.13655\\
2.7782	0.13289\\
2.7749	0.12931\\
2.7717	0.12582\\
2.7684	0.1224\\
2.7652	0.11906\\
2.762	0.1158\\
2.7587	0.11261\\
2.7555	0.1095\\
2.7523	0.10646\\
2.749	0.10349\\
2.7458	0.1006\\
2.7425	0.097772\\
2.7393	0.095016\\
2.7361	0.092326\\
2.7328	0.089703\\
2.7296	0.087145\\
2.7263	0.08465\\
2.7231	0.082218\\
2.7199	0.079848\\
2.7166	0.077537\\
2.7134	0.075286\\
2.7102	0.073092\\
2.7069	0.070955\\
2.7037	0.068873\\
2.7004	0.066845\\
2.6972	0.064871\\
2.694	0.062949\\
2.6907	0.061078\\
2.6875	0.059257\\
2.6843	0.057484\\
2.681	0.055759\\
2.6778	0.054081\\
2.6745	0.052449\\
2.6713	0.050861\\
2.6681	0.049317\\
2.6648	0.047815\\
2.6616	0.046355\\
2.6584	0.044936\\
2.6551	0.043556\\
2.6519	0.042215\\
2.6486	0.040912\\
2.6454	0.039646\\
2.6422	0.038416\\
2.6389	0.037221\\
2.6357	0.03606\\
2.6325	0.034933\\
2.6292	0.033838\\
2.626	0.032776\\
2.6227	0.031744\\
2.6195	0.030742\\
2.6163	0.02977\\
2.613	0.028826\\
2.6098	0.02791\\
2.6066	0.027022\\
2.6033	0.02616\\
2.6001	0.025324\\
2.5968	0.024513\\
2.5936	0.023726\\
2.5904	0.022963\\
2.5871	0.022223\\
2.5839	0.021506\\
2.5806	0.020811\\
2.5774	0.020137\\
2.5742	0.019483\\
2.5709	0.01885\\
2.5677	0.018237\\
2.5645	0.017642\\
2.5612	0.017066\\
2.558	0.016508\\
2.5547	0.015968\\
2.5515	0.015444\\
2.5483	0.014937\\
2.545	0.014446\\
2.5418	0.013971\\
2.5386	0.01351\\
2.5353	0.013065\\
2.5321	0.012633\\
2.5288	0.012216\\
2.5256	0.011812\\
2.5224	0.011421\\
2.5191	0.011042\\
2.5159	0.010676\\
2.5127	0.010322\\
2.5094	0.0099793\\
2.5062	0.0096478\\
2.5029	0.0093272\\
2.4997	0.0090171\\
2.4965	0.0087172\\
2.4932	0.0084272\\
2.49	0.0081468\\
2.4868	0.0078757\\
2.4835	0.0076136\\
2.4803	0.0073602\\
2.477	0.0071153\\
2.4738	0.0068785\\
2.4706	0.0066497\\
2.4673	0.0064286\\
2.4641	0.0062148\\
2.4608	0.0060083\\
2.4576	0.0058088\\
2.4544	0.005616\\
2.4511	0.0054297\\
2.4479	0.0052497\\
2.4447	0.0050759\\
2.4414	0.0049079\\
2.4382	0.0047457\\
2.4349	0.0045889\\
2.4317	0.0044376\\
2.4285	0.0042914\\
2.4252	0.0041502\\
2.422	0.0040138\\
2.4188	0.0038821\\
2.4155	0.0037549\\
2.4123	0.003632\\
2.409	0.0035134\\
2.4058	0.0033989\\
2.4026	0.0032883\\
2.3993	0.0031814\\
2.3961	0.0030783\\
2.3929	0.0029787\\
2.3896	0.0028825\\
2.3864	0.0027896\\
2.3831	0.0027\\
2.3799	0.0026134\\
2.3767	0.0025297\\
2.3734	0.002449\\
2.3702	0.002371\\
2.367	0.0022957\\
2.3637	0.002223\\
2.3605	0.0021527\\
2.3572	0.0020849\\
2.354	0.0020194\\
2.3508	0.0019562\\
2.3475	0.0018951\\
2.3443	0.001836\\
2.341	0.001779\\
2.3378	0.001724\\
2.3346	0.0016708\\
2.3313	0.0016194\\
2.3281	0.0015697\\
2.3249	0.0015217\\
2.3216	0.0014754\\
2.3184	0.0014306\\
2.3151	0.0013873\\
2.3119	0.0013455\\
2.3087	0.0013051\\
2.3054	0.001266\\
2.3022	0.0012282\\
2.299	0.0011917\\
2.2957	0.0011564\\
2.2925	0.0011222\\
2.2892	0.0010892\\
2.286	0.0010573\\
2.2828	0.0010264\\
2.2795	0.00099652\\
2.2763	0.00096762\\
2.2731	0.00093967\\
2.2698	0.00091261\\
2.2666	0.00088644\\
2.2633	0.0008611\\
2.2601	0.00083658\\
2.2569	0.00081285\\
2.2536	0.00078987\\
2.2504	0.00076762\\
2.2472	0.00074607\\
2.2439	0.0007252\\
2.2407	0.00070499\\
2.2374	0.00068541\\
2.2342	0.00066644\\
2.231	0.00064806\\
2.2277	0.00063025\\
2.2245	0.00061298\\
2.2212	0.00059625\\
2.218	0.00058003\\
2.2148	0.0005643\\
2.2115	0.00054904\\
2.2083	0.00053425\\
2.2051	0.0005199\\
2.2018	0.00050598\\
2.1986	0.00049247\\
2.1953	0.00047936\\
2.1921	0.00046664\\
2.1889	0.0004543\\
2.1856	0.00044231\\
2.1824	0.00043067\\
2.1792	0.00041938\\
2.1759	0.0004084\\
2.1727	0.00039774\\
2.1694	0.00038739\\
2.1662	0.00037733\\
2.163	0.00036756\\
2.1597	0.00035806\\
2.1565	0.00034883\\
2.1533	0.00033986\\
2.15	0.00033114\\
2.1468	0.00032266\\
2.1435	0.00031441\\
2.1403	0.00030639\\
2.1371	0.0002986\\
2.1338	0.00029101\\
2.1306	0.00028363\\
2.1274	0.00027645\\
2.1241	0.00026946\\
2.1209	0.00026266\\
2.1176	0.00025604\\
2.1144	0.0002496\\
2.1112	0.00024333\\
2.1079	0.00023723\\
2.1047	0.00023128\\
2.1014	0.00022549\\
2.0982	0.00021986\\
2.095	0.00021437\\
2.0917	0.00020902\\
2.0885	0.00020381\\
2.0853	0.00019874\\
2.082	0.00019379\\
2.0788	0.00018897\\
2.0755	0.00018428\\
}--cycle;
\addplot[ybar interval, fill=black, fill opacity=0.35, area legend, draw=none] table[row sep=crcr, x=Lower, y=Count] {%
Lower	Upper	Count\\
2.0755	2.3244	0.075022\\
2.3244	2.5732	0.13933\\
2.5732	2.822	0.35368\\
2.822	3.0708	0.76094\\
3.0708	3.3196	0.94314\\
3.3196	3.5684	0.75022\\
3.5684	3.8172	0.52516\\
3.8172	4.0661	0.2465\\
4.0661	4.3149	0.10717\\
4.3149	4.5637	0.053587\\
4.5637	4.8125	0.04287\\
4.8125	5.0613	0\\
5.0613	5.3101	0.021435\\
5.3101	5.3101	0.021435\\
};
\addlegendentry{Istogramma reale}

\addplot [color=black]
  table[row sep=crcr]{%
2.0755	0.016864\\
2.1144	0.021976\\
2.15	0.027791\\
2.1824	0.034175\\
2.2115	0.040944\\
2.2407	0.048804\\
2.2666	0.056806\\
2.2925	0.065853\\
2.3184	0.076034\\
2.3443	0.087435\\
2.3702	0.10014\\
2.3929	0.11239\\
2.4155	0.12576\\
2.4382	0.14028\\
2.4608	0.15599\\
2.4835	0.17292\\
2.5062	0.19111\\
2.5288	0.21056\\
2.5515	0.23127\\
2.5774	0.25648\\
2.6033	0.28329\\
2.6292	0.31165\\
2.6584	0.34529\\
2.6875	0.38062\\
2.7199	0.42161\\
2.7555	0.4684\\
2.8041	0.5341\\
2.8947	0.65712\\
2.9303	0.70335\\
2.9595	0.73941\\
2.9854	0.76973\\
3.008	0.79468\\
3.0307	0.81793\\
3.0501	0.83636\\
3.0696	0.8533\\
3.089	0.86863\\
3.1052	0.88011\\
3.1214	0.89036\\
3.1376	0.89934\\
3.1537	0.907\\
3.1699	0.91332\\
3.1861	0.91826\\
3.1991	0.92122\\
3.212	0.92327\\
3.225	0.92442\\
3.2379	0.92466\\
3.2509	0.92399\\
3.2638	0.92242\\
3.2768	0.91996\\
3.2897	0.91661\\
3.3059	0.91119\\
3.3221	0.90444\\
3.3383	0.89637\\
3.3545	0.88705\\
3.3707	0.8765\\
3.3869	0.86479\\
3.4063	0.84927\\
3.4257	0.83225\\
3.4451	0.81384\\
3.4678	0.79078\\
3.4905	0.76619\\
3.5164	0.73648\\
3.5455	0.70139\\
3.5811	0.65673\\
3.6362	0.58564\\
3.7042	0.49801\\
3.743	0.44963\\
3.7754	0.4109\\
3.8045	0.37756\\
3.8337	0.34584\\
3.8596	0.31913\\
3.8855	0.29389\\
3.9114	0.27016\\
3.9373	0.24797\\
3.9632	0.22732\\
3.9891	0.20819\\
4.015	0.19056\\
4.0409	0.17437\\
4.0668	0.15957\\
4.0927	0.14609\\
4.1186	0.13385\\
4.1445	0.12277\\
4.1704	0.11277\\
4.1996	0.10269\\
4.2287	0.093758\\
4.2611	0.085015\\
4.2935	0.077372\\
4.3291	0.070062\\
4.3679	0.063196\\
4.41	0.056832\\
4.4586	0.050596\\
4.5136	0.044623\\
4.5784	0.03868\\
4.6593	0.032391\\
4.7565	0.025978\\
4.8666	0.019813\\
4.9799	0.014524\\
5.0964	0.010153\\
5.2162	0.006727\\
5.3101	0.0047195\\
};
\addlegendentry{Adatt. BLN reale}

            % \input{../../../.tex/20/KS/SizeDistributionFittingsfig2.tex}
        \end{groupplot}
    
        \begin{groupplot}[
            group style={
                group name=Row2,
                plotGroup,
            },
            plotRow,
            row1Keys,
            plotLegendNE,
            %
            % xmin=1.7158,xmax=5.4719,
            ymin=0,ymax=1.4,
        ]
            \nextgroupplot[%
                at={($(Row1 c1r1.south west)-(0,.5*\yPlotSep em)-(0,\plotHeight cm)$)},
                xmin=1.7064,xmax=5.4903,
                % xmin=1.5262,xmax=5.4903,
                % ymin=0,ymax=1.4,
            ] % This file was created by matlab2tikz.
%
\definecolor{mycolor1}{rgb}{0.83529,0.36863,0.00000}%
%
\addplot[ybar interval, fill=mycolor1, fill opacity=0.35, area legend, draw=none] table[row sep=crcr, x=Lower, y=Count] {%
Lower	Upper	Count\\
1.7064	1.9836	9.6197e-05\\
1.9836	2.2608	0.00086577\\
2.2608	2.538	0.01289\\
2.538	2.8152	0.12775\\
2.8152	3.0925	0.60065\\
3.0925	3.3697	1.1215\\
3.3697	3.6469	0.88539\\
3.6469	3.9241	0.44664\\
3.9241	4.2013	0.22981\\
4.2013	4.4785	0.12429\\
4.4785	4.7557	0.04704\\
4.7557	5.0329	0.0097158\\
5.0329	5.3101	0.00076957\\
5.3101	5.3101	0.00076957\\
};
\addlegendentry{Istogramma appr.}

\addplot [color=mycolor1, only marks, every error bar/.append style={opacity=0.45}, mark=*, mark size=0pt, draw=none, forget plot]
 plot [error bars/.cd, y dir=both, y explicit, error bar style={line width=1pt, color=mycolor1}, error mark options={mark=none,mark size=0pt}]
 table[row sep=crcr, y error plus index=2, y error minus index=3]{%
1.845	9.6197e-05	0.00019087	9.6197e-05\\
2.1222	0.00086577	0.00061238	0.00061238\\
2.3994	0.01289	0.0029498	0.0029498\\
2.6766	0.12775	0.016101	0.016101\\
2.9538	0.60065	0.040269	0.040269\\
3.2311	1.1215	0.024597	0.024597\\
3.5083	0.88539	0.033302	0.033302\\
3.7855	0.44664	0.022533	0.022533\\
4.0627	0.22981	0.011394	0.011394\\
4.3399	0.12429	0.0073685	0.0073685\\
4.6171	0.04704	0.0048143	0.0048143\\
4.8943	0.0097158	0.0019656	0.0019656\\
5.1715	0.00076957	0.00052044	0.00052044\\
};
\addplot [color=mycolor1]
  table[row sep=crcr]{%
1.7064	4.0817e-05\\
2.0635	0.00068831\\
2.1934	0.0019195\\
2.2728	0.0036902\\
2.3305	0.0060213\\
2.3738	0.0087461\\
2.4098	0.011966\\
2.4423	0.015873\\
2.4712	0.020392\\
2.4964	0.025358\\
2.5217	0.031476\\
2.5433	0.037809\\
2.5649	0.045319\\
2.583	0.052604\\
2.601	0.06094\\
2.6191	0.070446\\
2.6371	0.081245\\
2.6551	0.093463\\
2.6732	0.10723\\
2.6876	0.11945\\
2.702	0.13281\\
2.7165	0.14737\\
2.7309	0.16319\\
2.7453	0.18033\\
2.7597	0.19885\\
2.7742	0.21877\\
2.7922	0.24572\\
2.8102	0.27497\\
2.8283	0.30656\\
2.8463	0.34046\\
2.8644	0.37663\\
2.8824	0.41498\\
2.904	0.46369\\
2.9257	0.51504\\
2.9509	0.57776\\
2.9834	0.66158\\
3.0628	0.86836\\
3.0844	0.92155\\
3.106	0.97181\\
3.1241	1.0109\\
3.1385	1.04\\
3.1529	1.0669\\
3.1674	1.0915\\
3.1818	1.1135\\
3.1926	1.1283\\
3.2034	1.1415\\
3.2143	1.1531\\
3.2251	1.163\\
3.2359	1.1712\\
3.2467	1.1777\\
3.2576	1.1825\\
3.2648	1.1847\\
3.272	1.1861\\
3.2792	1.1868\\
3.2864	1.1867\\
3.2936	1.1859\\
3.3008	1.1843\\
3.3117	1.1805\\
3.3225	1.175\\
3.3333	1.168\\
3.3441	1.1595\\
3.355	1.1495\\
3.3658	1.1381\\
3.3766	1.1254\\
3.391	1.1065\\
3.4055	1.0855\\
3.4199	1.0628\\
3.4379	1.0321\\
3.456	0.99934\\
3.4776	0.95787\\
3.5065	0.9001\\
3.5966	0.71731\\
3.6219	0.66913\\
3.6472	0.62337\\
3.6688	0.58637\\
3.6904	0.55158\\
3.7121	0.51908\\
3.7301	0.49378\\
3.7482	0.47007\\
3.7662	0.44794\\
3.7842	0.42732\\
3.8023	0.40814\\
3.8203	0.39032\\
3.8383	0.37378\\
3.86	0.35547\\
3.8816	0.33868\\
3.9033	0.32324\\
3.9285	0.3067\\
3.9538	0.2915\\
3.9826	0.27546\\
4.0151	0.25875\\
4.0512	0.24139\\
4.0981	0.22013\\
4.1558	0.19522\\
4.2243	0.16681\\
4.2893	0.14103\\
4.3434	0.12065\\
4.3939	0.10278\\
4.4408	0.087387\\
4.4841	0.074336\\
4.5237	0.063423\\
4.5634	0.053555\\
4.6031	0.044748\\
4.6428	0.036991\\
4.6825	0.030252\\
4.7257	0.023999\\
4.769	0.018804\\
4.8159	0.014243\\
4.8664	0.010406\\
4.9241	0.0071451\\
4.9891	0.0045902\\
5.0684	0.0026085\\
5.1694	0.0012293\\
5.3101	0.00040808\\
};
\addlegendentry{Adatt. BLN appr.}


\addplot[area legend, draw=none, fill=mycolor1, fill opacity=0.15, forget plot]
table[row sep=crcr] {%
x	y\\
1.7064	5.8891e-05\\
1.71	6.0656e-05\\
1.7136	6.2472e-05\\
1.7172	6.4339e-05\\
1.7208	6.6258e-05\\
1.7244	6.8232e-05\\
1.728	7.0262e-05\\
1.7317	7.2348e-05\\
1.7353	7.4493e-05\\
1.7389	7.6699e-05\\
1.7425	7.8966e-05\\
1.7461	8.1296e-05\\
1.7497	8.3692e-05\\
1.7533	8.6154e-05\\
1.7569	8.8685e-05\\
1.7605	9.1286e-05\\
1.7641	9.396e-05\\
1.7677	9.6707e-05\\
1.7713	9.9531e-05\\
1.7749	0.00010243\\
1.7785	0.00010541\\
1.7822	0.00010848\\
1.7858	0.00011163\\
1.7894	0.00011486\\
1.793	0.00011818\\
1.7966	0.0001216\\
1.8002	0.00012511\\
1.8038	0.00012871\\
1.8074	0.00013241\\
1.811	0.00013622\\
1.8146	0.00014013\\
1.8182	0.00014414\\
1.8218	0.00014826\\
1.8254	0.0001525\\
1.8291	0.00015685\\
1.8327	0.00016131\\
1.8363	0.0001659\\
1.8399	0.00017061\\
1.8435	0.00017545\\
1.8471	0.00018042\\
1.8507	0.00018553\\
1.8543	0.00019077\\
1.8579	0.00019615\\
1.8615	0.00020168\\
1.8651	0.00020736\\
1.8687	0.00021318\\
1.8723	0.00021917\\
1.8759	0.00022531\\
1.8796	0.00023162\\
1.8832	0.0002381\\
1.8868	0.00024475\\
1.8904	0.00025158\\
1.894	0.0002586\\
1.8976	0.00026579\\
1.9012	0.00027319\\
1.9048	0.00028077\\
1.9084	0.00028857\\
1.912	0.00029656\\
1.9156	0.00030478\\
1.9192	0.00031321\\
1.9228	0.00032186\\
1.9264	0.00033075\\
1.9301	0.00033987\\
1.9337	0.00034923\\
1.9373	0.00035884\\
1.9409	0.00036871\\
1.9445	0.00037885\\
1.9481	0.00038925\\
1.9517	0.00039993\\
1.9553	0.00041089\\
1.9589	0.00042214\\
1.9625	0.0004337\\
1.9661	0.00044556\\
1.9697	0.00045774\\
1.9733	0.00047024\\
1.977	0.00048307\\
1.9806	0.00049625\\
1.9842	0.00050978\\
1.9878	0.00052367\\
1.9914	0.00053793\\
1.995	0.00055258\\
1.9986	0.00056761\\
2.0022	0.00058305\\
2.0058	0.0005989\\
2.0094	0.00061518\\
2.013	0.00063189\\
2.0166	0.00064905\\
2.0202	0.00066668\\
2.0238	0.00068478\\
2.0275	0.00070337\\
2.0311	0.00072246\\
2.0347	0.00074206\\
2.0383	0.0007622\\
2.0419	0.00078288\\
2.0455	0.00080412\\
2.0491	0.00082594\\
2.0527	0.00084836\\
2.0563	0.00087138\\
2.0599	0.00089504\\
2.0635	0.00091934\\
2.0671	0.00094431\\
2.0707	0.00096996\\
2.0743	0.00099632\\
2.078	0.0010234\\
2.0816	0.0010512\\
2.0852	0.0010798\\
2.0888	0.0011092\\
2.0924	0.0011395\\
2.096	0.0011705\\
2.0996	0.0012024\\
2.1032	0.0012352\\
2.1068	0.001269\\
2.1104	0.0013036\\
2.114	0.0013393\\
2.1176	0.001376\\
2.1212	0.0014137\\
2.1249	0.0014524\\
2.1285	0.0014923\\
2.1321	0.0015333\\
2.1357	0.0015755\\
2.1393	0.0016188\\
2.1429	0.0016634\\
2.1465	0.0017094\\
2.1501	0.0017566\\
2.1537	0.0018052\\
2.1573	0.0018552\\
2.1609	0.0019066\\
2.1645	0.0019596\\
2.1681	0.0020141\\
2.1717	0.0020702\\
2.1754	0.0021279\\
2.179	0.0021874\\
2.1826	0.0022486\\
2.1862	0.0023116\\
2.1898	0.0023765\\
2.1934	0.0024433\\
2.197	0.0025121\\
2.2006	0.0025829\\
2.2042	0.0026559\\
2.2078	0.0027311\\
2.2114	0.0028085\\
2.215	0.0028883\\
2.2186	0.0029705\\
2.2223	0.0030552\\
2.2259	0.0031424\\
2.2295	0.0032323\\
2.2331	0.003325\\
2.2367	0.0034205\\
2.2403	0.0035189\\
2.2439	0.0036204\\
2.2475	0.003725\\
2.2511	0.0038328\\
2.2547	0.003944\\
2.2583	0.0040586\\
2.2619	0.0041767\\
2.2655	0.0042986\\
2.2691	0.0044242\\
2.2728	0.0045538\\
2.2764	0.0046875\\
2.28	0.0048253\\
2.2836	0.0049674\\
2.2872	0.005114\\
2.2908	0.0052653\\
2.2944	0.0054213\\
2.298	0.0055822\\
2.3016	0.0057482\\
2.3052	0.0059195\\
2.3088	0.0060962\\
2.3124	0.0062785\\
2.316	0.0064666\\
2.3196	0.0066607\\
2.3233	0.006861\\
2.3269	0.0070676\\
2.3305	0.0072808\\
2.3341	0.0075007\\
2.3377	0.0077277\\
2.3413	0.007962\\
2.3449	0.0082037\\
2.3485	0.0084531\\
2.3521	0.0087104\\
2.3557	0.008976\\
2.3593	0.0092501\\
2.3629	0.0095328\\
2.3665	0.0098246\\
2.3702	0.010126\\
2.3738	0.010436\\
2.3774	0.010757\\
2.381	0.011088\\
2.3846	0.011429\\
2.3882	0.011781\\
2.3918	0.012144\\
2.3954	0.012519\\
2.399	0.012906\\
2.4026	0.013305\\
2.4062	0.013716\\
2.4098	0.014141\\
2.4134	0.014579\\
2.417	0.01503\\
2.4207	0.015496\\
2.4243	0.015977\\
2.4279	0.016472\\
2.4315	0.016983\\
2.4351	0.01751\\
2.4387	0.018053\\
2.4423	0.018613\\
2.4459	0.019191\\
2.4495	0.019786\\
2.4531	0.0204\\
2.4567	0.021033\\
2.4603	0.021685\\
2.4639	0.022357\\
2.4675	0.023049\\
2.4712	0.023763\\
2.4748	0.024498\\
2.4784	0.025255\\
2.482	0.026036\\
2.4856	0.02684\\
2.4892	0.027668\\
2.4928	0.02852\\
2.4964	0.029398\\
2.5	0.030303\\
2.5036	0.031234\\
2.5072	0.032192\\
2.5108	0.033178\\
2.5144	0.034194\\
2.5181	0.035239\\
2.5217	0.036314\\
2.5253	0.037421\\
2.5289	0.038559\\
2.5325	0.03973\\
2.5361	0.040935\\
2.5397	0.042173\\
2.5433	0.043447\\
2.5469	0.044757\\
2.5505	0.046103\\
2.5541	0.047487\\
2.5577	0.04891\\
2.5613	0.050372\\
2.5649	0.051874\\
2.5686	0.053418\\
2.5722	0.055003\\
2.5758	0.056632\\
2.5794	0.058304\\
2.583	0.060022\\
2.5866	0.061785\\
2.5902	0.063595\\
2.5938	0.065453\\
2.5974	0.06736\\
2.601	0.069316\\
2.6046	0.071323\\
2.6082	0.073382\\
2.6118	0.075494\\
2.6155	0.07766\\
2.6191	0.079881\\
2.6227	0.082157\\
2.6263	0.08449\\
2.6299	0.086882\\
2.6335	0.089332\\
2.6371	0.091843\\
2.6407	0.094414\\
2.6443	0.097048\\
2.6479	0.099745\\
2.6515	0.10251\\
2.6551	0.10533\\
2.6587	0.10823\\
2.6623	0.11119\\
2.666	0.11422\\
2.6696	0.11731\\
2.6732	0.12048\\
2.6768	0.12372\\
2.6804	0.12704\\
2.684	0.13042\\
2.6876	0.13389\\
2.6912	0.13742\\
2.6948	0.14104\\
2.6984	0.14473\\
2.702	0.1485\\
2.7056	0.15235\\
2.7092	0.15628\\
2.7128	0.16029\\
2.7165	0.16439\\
2.7201	0.16856\\
2.7237	0.17282\\
2.7273	0.17717\\
2.7309	0.1816\\
2.7345	0.18612\\
2.7381	0.19072\\
2.7417	0.19542\\
2.7453	0.2002\\
2.7489	0.20507\\
2.7525	0.21003\\
2.7561	0.21508\\
2.7597	0.22022\\
2.7634	0.22545\\
2.767	0.23078\\
2.7706	0.2362\\
2.7742	0.24171\\
2.7778	0.24731\\
2.7814	0.25301\\
2.785	0.2588\\
2.7886	0.26469\\
2.7922	0.27067\\
2.7958	0.27674\\
2.7994	0.28291\\
2.803	0.28917\\
2.8066	0.29553\\
2.8102	0.30198\\
2.8139	0.30853\\
2.8175	0.31517\\
2.8211	0.3219\\
2.8247	0.32873\\
2.8283	0.33565\\
2.8319	0.34267\\
2.8355	0.34977\\
2.8391	0.35697\\
2.8427	0.36426\\
2.8463	0.37164\\
2.8499	0.37911\\
2.8535	0.38667\\
2.8571	0.39432\\
2.8607	0.40205\\
2.8644	0.40987\\
2.868	0.41778\\
2.8716	0.42576\\
2.8752	0.43383\\
2.8788	0.44199\\
2.8824	0.45022\\
2.886	0.45853\\
2.8896	0.46692\\
2.8932	0.47538\\
2.8968	0.48391\\
2.9004	0.49252\\
2.904	0.5012\\
2.9076	0.50994\\
2.9113	0.51876\\
2.9149	0.52763\\
2.9185	0.53657\\
2.9221	0.54557\\
2.9257	0.55462\\
2.9293	0.56373\\
2.9329	0.57289\\
2.9365	0.5821\\
2.9401	0.59136\\
2.9437	0.60067\\
2.9473	0.61001\\
2.9509	0.6194\\
2.9545	0.62882\\
2.9581	0.63828\\
2.9618	0.64777\\
2.9654	0.65728\\
2.969	0.66682\\
2.9726	0.67638\\
2.9762	0.68596\\
2.9798	0.69556\\
2.9834	0.70516\\
2.987	0.71478\\
2.9906	0.7244\\
2.9942	0.73402\\
2.9978	0.74364\\
3.0014	0.75325\\
3.005	0.76286\\
3.0087	0.77245\\
3.0123	0.78202\\
3.0159	0.79158\\
3.0195	0.80111\\
3.0231	0.81062\\
3.0267	0.82009\\
3.0303	0.82953\\
3.0339	0.83892\\
3.0375	0.84828\\
3.0411	0.85759\\
3.0447	0.86685\\
3.0483	0.87605\\
3.0519	0.8852\\
3.0555	0.89428\\
3.0592	0.9033\\
3.0628	0.91225\\
3.0664	0.92112\\
3.07	0.92992\\
3.0736	0.93864\\
3.0772	0.94727\\
3.0808	0.95581\\
3.0844	0.96425\\
3.088	0.97261\\
3.0916	0.98086\\
3.0952	0.98901\\
3.0988	0.99705\\
3.1024	1.005\\
3.106	1.0128\\
3.1097	1.0205\\
3.1133	1.0281\\
3.1169	1.0355\\
3.1205	1.0428\\
3.1241	1.05\\
3.1277	1.0571\\
3.1313	1.064\\
3.1349	1.0708\\
3.1385	1.0774\\
3.1421	1.0839\\
3.1457	1.0902\\
3.1493	1.0964\\
3.1529	1.1024\\
3.1566	1.1083\\
3.1602	1.114\\
3.1638	1.1196\\
3.1674	1.125\\
3.171	1.1302\\
3.1746	1.1352\\
3.1782	1.1401\\
3.1818	1.1448\\
3.1854	1.1493\\
3.189	1.1537\\
3.1926	1.1579\\
3.1962	1.1619\\
3.1998	1.1657\\
3.2034	1.1694\\
3.2071	1.1728\\
3.2107	1.1761\\
3.2143	1.1792\\
3.2179	1.1822\\
3.2215	1.1849\\
3.2251	1.1875\\
3.2287	1.1899\\
3.2323	1.1921\\
3.2359	1.1941\\
3.2395	1.1959\\
3.2431	1.1976\\
3.2467	1.1991\\
3.2503	1.2004\\
3.254	1.2015\\
3.2576	1.2025\\
3.2612	1.2033\\
3.2648	1.204\\
3.2684	1.2044\\
3.272	1.2047\\
3.2756	1.2049\\
3.2792	1.2049\\
3.2828	1.2047\\
3.2864	1.2044\\
3.29	1.2039\\
3.2936	1.2033\\
3.2972	1.2025\\
3.3008	1.2016\\
3.3045	1.2006\\
3.3081	1.1994\\
3.3117	1.198\\
3.3153	1.1965\\
3.3189	1.1949\\
3.3225	1.1931\\
3.3261	1.1911\\
3.3297	1.1891\\
3.3333	1.1868\\
3.3369	1.1845\\
3.3405	1.1819\\
3.3441	1.1793\\
3.3477	1.1765\\
3.3513	1.1735\\
3.355	1.1704\\
3.3586	1.1672\\
3.3622	1.1638\\
3.3658	1.1603\\
3.3694	1.1567\\
3.373	1.1529\\
3.3766	1.1489\\
3.3802	1.1449\\
3.3838	1.1406\\
3.3874	1.1363\\
3.391	1.1318\\
3.3946	1.1272\\
3.3982	1.1225\\
3.4019	1.1177\\
3.4055	1.1127\\
3.4091	1.1076\\
3.4127	1.1024\\
3.4163	1.0971\\
3.4199	1.0916\\
3.4235	1.0861\\
3.4271	1.0805\\
3.4307	1.0747\\
3.4343	1.0689\\
3.4379	1.0629\\
3.4415	1.0569\\
3.4451	1.0507\\
3.4487	1.0445\\
3.4524	1.0382\\
3.456	1.0318\\
3.4596	1.0254\\
3.4632	1.0188\\
3.4668	1.0122\\
3.4704	1.0055\\
3.474	0.9988\\
3.4776	0.992\\
3.4812	0.98513\\
3.4848	0.97821\\
3.4884	0.97124\\
3.492	0.96422\\
3.4956	0.95715\\
3.4992	0.95004\\
3.5029	0.94289\\
3.5065	0.9357\\
3.5101	0.92848\\
3.5137	0.92122\\
3.5173	0.91394\\
3.5209	0.90664\\
3.5245	0.89931\\
3.5281	0.89197\\
3.5317	0.88461\\
3.5353	0.87723\\
3.5389	0.86985\\
3.5425	0.86246\\
3.5461	0.85506\\
3.5498	0.84767\\
3.5534	0.84027\\
3.557	0.83288\\
3.5606	0.82549\\
3.5642	0.81811\\
3.5678	0.81074\\
3.5714	0.80339\\
3.575	0.79605\\
3.5786	0.78873\\
3.5822	0.78143\\
3.5858	0.77416\\
3.5894	0.7669\\
3.593	0.75968\\
3.5966	0.75248\\
3.6003	0.74532\\
3.6039	0.73818\\
3.6075	0.73108\\
3.6111	0.72402\\
3.6147	0.717\\
3.6183	0.71001\\
3.6219	0.70307\\
3.6255	0.69617\\
3.6291	0.68932\\
3.6327	0.68251\\
3.6363	0.67575\\
3.6399	0.66903\\
3.6435	0.66237\\
3.6472	0.65576\\
3.6508	0.6492\\
3.6544	0.64269\\
3.658	0.63624\\
3.6616	0.62985\\
3.6652	0.62351\\
3.6688	0.61722\\
3.6724	0.611\\
3.676	0.60483\\
3.6796	0.59873\\
3.6832	0.59268\\
3.6868	0.5867\\
3.6904	0.58077\\
3.694	0.57491\\
3.6977	0.56911\\
3.7013	0.56337\\
3.7049	0.5577\\
3.7085	0.55209\\
3.7121	0.54654\\
3.7157	0.54106\\
3.7193	0.53564\\
3.7229	0.53028\\
3.7265	0.52499\\
3.7301	0.51976\\
3.7337	0.5146\\
3.7373	0.5095\\
3.7409	0.50447\\
3.7445	0.4995\\
3.7482	0.49459\\
3.7518	0.48974\\
3.7554	0.48496\\
3.759	0.48024\\
3.7626	0.47559\\
3.7662	0.471\\
3.7698	0.46647\\
3.7734	0.462\\
3.777	0.45759\\
3.7806	0.45324\\
3.7842	0.44895\\
3.7878	0.44473\\
3.7914	0.44056\\
3.7951	0.43645\\
3.7987	0.4324\\
3.8023	0.42841\\
3.8059	0.42447\\
3.8095	0.42059\\
3.8131	0.41676\\
3.8167	0.413\\
3.8203	0.40928\\
3.8239	0.40562\\
3.8275	0.40201\\
3.8311	0.39845\\
3.8347	0.39495\\
3.8383	0.3915\\
3.8419	0.38809\\
3.8456	0.38474\\
3.8492	0.38143\\
3.8528	0.37818\\
3.8564	0.37496\\
3.86	0.3718\\
3.8636	0.36868\\
3.8672	0.36561\\
3.8708	0.36258\\
3.8744	0.35959\\
3.878	0.35664\\
3.8816	0.35374\\
3.8852	0.35088\\
3.8888	0.34806\\
3.8924	0.34527\\
3.8961	0.34253\\
3.8997	0.33982\\
3.9033	0.33715\\
3.9069	0.33451\\
3.9105	0.33191\\
3.9141	0.32935\\
3.9177	0.32681\\
3.9213	0.32432\\
3.9249	0.32185\\
3.9285	0.31941\\
3.9321	0.31701\\
3.9357	0.31463\\
3.9393	0.31229\\
3.943	0.30997\\
3.9466	0.30768\\
3.9502	0.30542\\
3.9538	0.30318\\
3.9574	0.30097\\
3.961	0.29878\\
3.9646	0.29662\\
3.9682	0.29448\\
3.9718	0.29237\\
3.9754	0.29027\\
3.979	0.2882\\
3.9826	0.28615\\
3.9862	0.28412\\
3.9898	0.28211\\
3.9935	0.28012\\
3.9971	0.27815\\
4.0007	0.27619\\
4.0043	0.27426\\
4.0079	0.27234\\
4.0115	0.27043\\
4.0151	0.26854\\
4.0187	0.26667\\
4.0223	0.26481\\
4.0259	0.26297\\
4.0295	0.26114\\
4.0331	0.25932\\
4.0367	0.25752\\
4.0404	0.25573\\
4.044	0.25395\\
4.0476	0.25218\\
4.0512	0.25043\\
4.0548	0.24868\\
4.0584	0.24695\\
4.062	0.24522\\
4.0656	0.24351\\
4.0692	0.2418\\
4.0728	0.24011\\
4.0764	0.23842\\
4.08	0.23674\\
4.0836	0.23507\\
4.0872	0.2334\\
4.0909	0.23175\\
4.0945	0.2301\\
4.0981	0.22846\\
4.1017	0.22682\\
4.1053	0.22519\\
4.1089	0.22357\\
4.1125	0.22195\\
4.1161	0.22034\\
4.1197	0.21874\\
4.1233	0.21714\\
4.1269	0.21554\\
4.1305	0.21395\\
4.1341	0.21237\\
4.1377	0.21079\\
4.1414	0.20921\\
4.145	0.20764\\
4.1486	0.20608\\
4.1522	0.20451\\
4.1558	0.20296\\
4.1594	0.2014\\
4.163	0.19985\\
4.1666	0.1983\\
4.1702	0.19676\\
4.1738	0.19522\\
4.1774	0.19369\\
4.181	0.19216\\
4.1846	0.19063\\
4.1883	0.18911\\
4.1919	0.18758\\
4.1955	0.18607\\
4.1991	0.18455\\
4.2027	0.18304\\
4.2063	0.18154\\
4.2099	0.18003\\
4.2135	0.17853\\
4.2171	0.17704\\
4.2207	0.17554\\
4.2243	0.17405\\
4.2279	0.17257\\
4.2315	0.17109\\
4.2351	0.16961\\
4.2388	0.16813\\
4.2424	0.16666\\
4.246	0.16519\\
4.2496	0.16373\\
4.2532	0.16227\\
4.2568	0.16081\\
4.2604	0.15936\\
4.264	0.15791\\
4.2676	0.15646\\
4.2712	0.15502\\
4.2748	0.15359\\
4.2784	0.15215\\
4.282	0.15072\\
4.2856	0.1493\\
4.2893	0.14788\\
4.2929	0.14647\\
4.2965	0.14506\\
4.3001	0.14365\\
4.3037	0.14225\\
4.3073	0.14085\\
4.3109	0.13946\\
4.3145	0.13807\\
4.3181	0.13669\\
4.3217	0.13531\\
4.3253	0.13394\\
4.3289	0.13258\\
4.3325	0.13122\\
4.3362	0.12986\\
4.3398	0.12851\\
4.3434	0.12716\\
4.347	0.12583\\
4.3506	0.12449\\
4.3542	0.12316\\
4.3578	0.12184\\
4.3614	0.12053\\
4.365	0.11922\\
4.3686	0.11791\\
4.3722	0.11662\\
4.3758	0.11532\\
4.3794	0.11404\\
4.383	0.11276\\
4.3867	0.11149\\
4.3903	0.11022\\
4.3939	0.10896\\
4.3975	0.10771\\
4.4011	0.10646\\
4.4047	0.10522\\
4.4083	0.10399\\
4.4119	0.10277\\
4.4155	0.10155\\
4.4191	0.10034\\
4.4227	0.099131\\
4.4263	0.097933\\
4.4299	0.096743\\
4.4336	0.09556\\
4.4372	0.094384\\
4.4408	0.093215\\
4.4444	0.092054\\
4.448	0.090901\\
4.4516	0.089754\\
4.4552	0.088616\\
4.4588	0.087485\\
4.4624	0.086361\\
4.466	0.085245\\
4.4696	0.084137\\
4.4732	0.083037\\
4.4768	0.081945\\
4.4804	0.08086\\
4.4841	0.079783\\
4.4877	0.078715\\
4.4913	0.077654\\
4.4949	0.076601\\
4.4985	0.075556\\
4.5021	0.074519\\
4.5057	0.073491\\
4.5093	0.07247\\
4.5129	0.071458\\
4.5165	0.070453\\
4.5201	0.069457\\
4.5237	0.06847\\
4.5273	0.06749\\
4.5309	0.066519\\
4.5346	0.065556\\
4.5382	0.064601\\
4.5418	0.063655\\
4.5454	0.062717\\
4.549	0.061788\\
4.5526	0.060867\\
4.5562	0.059954\\
4.5598	0.05905\\
4.5634	0.058154\\
4.567	0.057267\\
4.5706	0.056388\\
4.5742	0.055518\\
4.5778	0.054656\\
4.5815	0.053803\\
4.5851	0.052958\\
4.5887	0.052122\\
4.5923	0.051294\\
4.5959	0.050475\\
4.5995	0.049664\\
4.6031	0.048862\\
4.6067	0.048068\\
4.6103	0.047283\\
4.6139	0.046507\\
4.6175	0.045739\\
4.6211	0.044979\\
4.6247	0.044228\\
4.6283	0.043485\\
4.632	0.042751\\
4.6356	0.042026\\
4.6392	0.041309\\
4.6428	0.0406\\
4.6464	0.0399\\
4.65	0.039208\\
4.6536	0.038525\\
4.6572	0.03785\\
4.6608	0.037183\\
4.6644	0.036525\\
4.668	0.035875\\
4.6716	0.035233\\
4.6752	0.0346\\
4.6788	0.033974\\
4.6825	0.033357\\
4.6861	0.032749\\
4.6897	0.032148\\
4.6933	0.031555\\
4.6969	0.030971\\
4.7005	0.030394\\
4.7041	0.029826\\
4.7077	0.029265\\
4.7113	0.028713\\
4.7149	0.028168\\
4.7185	0.027631\\
4.7221	0.027102\\
4.7257	0.026581\\
4.7294	0.026067\\
4.733	0.025561\\
4.7366	0.025062\\
4.7402	0.024572\\
4.7438	0.024088\\
4.7474	0.023612\\
4.751	0.023144\\
4.7546	0.022683\\
4.7582	0.022229\\
4.7618	0.021782\\
4.7654	0.021343\\
4.769	0.02091\\
4.7726	0.020485\\
4.7762	0.020066\\
4.7799	0.019655\\
4.7835	0.01925\\
4.7871	0.018852\\
4.7907	0.018461\\
4.7943	0.018077\\
4.7979	0.017699\\
4.8015	0.017328\\
4.8051	0.016963\\
4.8087	0.016604\\
4.8123	0.016252\\
4.8159	0.015906\\
4.8195	0.015567\\
4.8231	0.015233\\
4.8268	0.014906\\
4.8304	0.014584\\
4.834	0.014268\\
4.8376	0.013958\\
4.8412	0.013654\\
4.8448	0.013356\\
4.8484	0.013063\\
4.852	0.012776\\
4.8556	0.012494\\
4.8592	0.012217\\
4.8628	0.011946\\
4.8664	0.01168\\
4.87	0.01142\\
4.8736	0.011164\\
4.8773	0.010913\\
4.8809	0.010668\\
4.8845	0.010427\\
4.8881	0.010191\\
4.8917	0.0099595\\
4.8953	0.0097329\\
4.8989	0.0095108\\
4.9025	0.0092932\\
4.9061	0.0090801\\
4.9097	0.0088714\\
4.9133	0.0086669\\
4.9169	0.0084667\\
4.9205	0.0082706\\
4.9241	0.0080786\\
4.9278	0.0078906\\
4.9314	0.0077066\\
4.935	0.0075264\\
4.9386	0.0073501\\
4.9422	0.0071775\\
4.9458	0.0070085\\
4.9494	0.0068432\\
4.953	0.0066815\\
4.9566	0.0065232\\
4.9602	0.0063683\\
4.9638	0.0062168\\
4.9674	0.0060686\\
4.971	0.0059237\\
4.9747	0.0057819\\
4.9783	0.0056432\\
4.9819	0.0055076\\
4.9855	0.0053749\\
4.9891	0.0052452\\
4.9927	0.0051184\\
4.9963	0.0049945\\
4.9999	0.0048733\\
5.0035	0.0047548\\
5.0071	0.0046389\\
5.0107	0.0045257\\
5.0143	0.0044151\\
5.0179	0.0043069\\
5.0215	0.0042012\\
5.0252	0.0040979\\
5.0288	0.003997\\
5.0324	0.0038984\\
5.036	0.003802\\
5.0396	0.0037079\\
5.0432	0.0036159\\
5.0468	0.003526\\
5.0504	0.0034383\\
5.054	0.0033525\\
5.0576	0.0032688\\
5.0612	0.003187\\
5.0648	0.0031071\\
5.0684	0.0030291\\
5.072	0.0029529\\
5.0757	0.0028785\\
5.0793	0.0028059\\
5.0829	0.002735\\
5.0865	0.0026657\\
5.0901	0.0025981\\
5.0937	0.0025321\\
5.0973	0.0024677\\
5.1009	0.0024048\\
5.1045	0.0023434\\
5.1081	0.0022835\\
5.1117	0.002225\\
5.1153	0.0021679\\
5.1189	0.0021122\\
5.1226	0.0020579\\
5.1262	0.0020048\\
5.1298	0.0019531\\
5.1334	0.0019026\\
5.137	0.0018533\\
5.1406	0.0018052\\
5.1442	0.0017583\\
5.1478	0.0017126\\
5.1514	0.0016679\\
5.155	0.0016244\\
5.1586	0.0015819\\
5.1622	0.0015405\\
5.1658	0.0015001\\
5.1694	0.0014607\\
5.1731	0.0014222\\
5.1767	0.0013848\\
5.1803	0.0013482\\
5.1839	0.0013126\\
5.1875	0.0012778\\
5.1911	0.0012439\\
5.1947	0.0012109\\
5.1983	0.0011787\\
5.2019	0.0011473\\
5.2055	0.0011167\\
5.2091	0.0010868\\
5.2127	0.0010577\\
5.2163	0.0010294\\
5.22	0.0010018\\
5.2236	0.00097482\\
5.2272	0.00094857\\
5.2308	0.00092299\\
5.2344	0.00089806\\
5.238	0.00087376\\
5.2416	0.00085009\\
5.2452	0.00082702\\
5.2488	0.00080454\\
5.2524	0.00078264\\
5.256	0.0007613\\
5.2596	0.00074052\\
5.2632	0.00072027\\
5.2668	0.00070054\\
5.2705	0.00068133\\
5.2741	0.00066261\\
5.2777	0.00064438\\
5.2813	0.00062662\\
5.2849	0.00060933\\
5.2885	0.00059249\\
5.2921	0.00057609\\
5.2957	0.00056012\\
5.2993	0.00054457\\
5.3029	0.00052942\\
5.3065	0.00051468\\
5.3101	0.00050032\\
5.3101	0.00031583\\
5.3065	0.00032554\\
5.3029	0.00033553\\
5.2993	0.00034581\\
5.2957	0.00035639\\
5.2921	0.00036728\\
5.2885	0.0003785\\
5.2849	0.00039003\\
5.2813	0.00040191\\
5.2777	0.00041412\\
5.2741	0.0004267\\
5.2705	0.00043963\\
5.2668	0.00045294\\
5.2632	0.00046664\\
5.2596	0.00048073\\
5.256	0.00049523\\
5.2524	0.00051014\\
5.2488	0.00052548\\
5.2452	0.00054126\\
5.2416	0.00055749\\
5.238	0.00057418\\
5.2344	0.00059135\\
5.2308	0.00060901\\
5.2272	0.00062717\\
5.2236	0.00064585\\
5.22	0.00066506\\
5.2163	0.0006848\\
5.2127	0.00070511\\
5.2091	0.00072599\\
5.2055	0.00074745\\
5.2019	0.00076952\\
5.1983	0.00079221\\
5.1947	0.00081553\\
5.1911	0.0008395\\
5.1875	0.00086414\\
5.1839	0.00088946\\
5.1803	0.00091549\\
5.1767	0.00094224\\
5.1731	0.00096973\\
5.1694	0.00099797\\
5.1658	0.001027\\
5.1622	0.0010568\\
5.1586	0.0010875\\
5.155	0.0011189\\
5.1514	0.0011513\\
5.1478	0.0011845\\
5.1442	0.0012186\\
5.1406	0.0012537\\
5.137	0.0012897\\
5.1334	0.0013266\\
5.1298	0.0013646\\
5.1262	0.0014036\\
5.1226	0.0014436\\
5.1189	0.0014847\\
5.1153	0.0015269\\
5.1117	0.0015702\\
5.1081	0.0016147\\
5.1045	0.0016603\\
5.1009	0.0017072\\
5.0973	0.0017553\\
5.0937	0.0018046\\
5.0901	0.0018553\\
5.0865	0.0019072\\
5.0829	0.0019605\\
5.0793	0.0020152\\
5.0757	0.0020713\\
5.072	0.0021289\\
5.0684	0.0021879\\
5.0648	0.0022485\\
5.0612	0.0023106\\
5.0576	0.0023743\\
5.054	0.0024396\\
5.0504	0.0025065\\
5.0468	0.0025752\\
5.0432	0.0026456\\
5.0396	0.0027177\\
5.036	0.0027916\\
5.0324	0.0028674\\
5.0288	0.0029451\\
5.0252	0.0030247\\
5.0215	0.0031062\\
5.0179	0.0031897\\
5.0143	0.0032753\\
5.0107	0.003363\\
5.0071	0.0034528\\
5.0035	0.0035447\\
4.9999	0.0036389\\
4.9963	0.0037353\\
4.9927	0.0038341\\
4.9891	0.0039352\\
4.9855	0.0040386\\
4.9819	0.0041446\\
4.9783	0.004253\\
4.9747	0.0043639\\
4.971	0.0044774\\
4.9674	0.0045936\\
4.9638	0.0047125\\
4.9602	0.004834\\
4.9566	0.0049584\\
4.953	0.0050856\\
4.9494	0.0052157\\
4.9458	0.0053487\\
4.9422	0.0054847\\
4.9386	0.0056238\\
4.935	0.0057659\\
4.9314	0.0059112\\
4.9278	0.0060598\\
4.9241	0.0062115\\
4.9205	0.0063667\\
4.9169	0.0065251\\
4.9133	0.0066871\\
4.9097	0.0068525\\
4.9061	0.0070214\\
4.9025	0.007194\\
4.8989	0.0073702\\
4.8953	0.0075502\\
4.8917	0.0077339\\
4.8881	0.0079215\\
4.8845	0.0081131\\
4.8809	0.0083085\\
4.8773	0.0085081\\
4.8736	0.0087117\\
4.87	0.0089195\\
4.8664	0.0091315\\
4.8628	0.0093478\\
4.8592	0.0095684\\
4.8556	0.0097935\\
4.852	0.010023\\
4.8484	0.010257\\
4.8448	0.010496\\
4.8412	0.010739\\
4.8376	0.010987\\
4.834	0.01124\\
4.8304	0.011498\\
4.8268	0.011761\\
4.8231	0.012029\\
4.8195	0.012302\\
4.8159	0.01258\\
4.8123	0.012863\\
4.8087	0.013152\\
4.8051	0.013446\\
4.8015	0.013745\\
4.7979	0.01405\\
4.7943	0.014361\\
4.7907	0.014677\\
4.7871	0.014999\\
4.7835	0.015327\\
4.7799	0.01566\\
4.7762	0.016\\
4.7726	0.016345\\
4.769	0.016697\\
4.7654	0.017055\\
4.7618	0.017419\\
4.7582	0.017789\\
4.7546	0.018166\\
4.751	0.018549\\
4.7474	0.018939\\
4.7438	0.019335\\
4.7402	0.019738\\
4.7366	0.020147\\
4.733	0.020564\\
4.7294	0.020987\\
4.7257	0.021417\\
4.7221	0.021855\\
4.7185	0.022299\\
4.7149	0.022751\\
4.7113	0.023209\\
4.7077	0.023675\\
4.7041	0.024149\\
4.7005	0.024629\\
4.6969	0.025118\\
4.6933	0.025614\\
4.6897	0.026117\\
4.6861	0.026628\\
4.6825	0.027147\\
4.6788	0.027674\\
4.6752	0.028208\\
4.6716	0.028751\\
4.668	0.029301\\
4.6644	0.029859\\
4.6608	0.030426\\
4.6572	0.031001\\
4.6536	0.031584\\
4.65	0.032175\\
4.6464	0.032774\\
4.6428	0.033382\\
4.6392	0.033999\\
4.6356	0.034623\\
4.632	0.035257\\
4.6283	0.035898\\
4.6247	0.036549\\
4.6211	0.037208\\
4.6175	0.037876\\
4.6139	0.038552\\
4.6103	0.039237\\
4.6067	0.039931\\
4.6031	0.040634\\
4.5995	0.041346\\
4.5959	0.042067\\
4.5923	0.042796\\
4.5887	0.043535\\
4.5851	0.044282\\
4.5815	0.045039\\
4.5778	0.045804\\
4.5742	0.046579\\
4.5706	0.047362\\
4.567	0.048155\\
4.5634	0.048957\\
4.5598	0.049767\\
4.5562	0.050587\\
4.5526	0.051416\\
4.549	0.052254\\
4.5454	0.053102\\
4.5418	0.053958\\
4.5382	0.054823\\
4.5346	0.055698\\
4.5309	0.056582\\
4.5273	0.057474\\
4.5237	0.058376\\
4.5201	0.059287\\
4.5165	0.060207\\
4.5129	0.061136\\
4.5093	0.062073\\
4.5057	0.06302\\
4.5021	0.063976\\
4.4985	0.064941\\
4.4949	0.065914\\
4.4913	0.066897\\
4.4877	0.067888\\
4.4841	0.068888\\
4.4804	0.069897\\
4.4768	0.070915\\
4.4732	0.071941\\
4.4696	0.072976\\
4.466	0.074019\\
4.4624	0.075071\\
4.4588	0.076131\\
4.4552	0.0772\\
4.4516	0.078277\\
4.448	0.079363\\
4.4444	0.080456\\
4.4408	0.081558\\
4.4372	0.082668\\
4.4336	0.083786\\
4.4299	0.084912\\
4.4263	0.086046\\
4.4227	0.087188\\
4.4191	0.088337\\
4.4155	0.089494\\
4.4119	0.090659\\
4.4083	0.091832\\
4.4047	0.093011\\
4.4011	0.094199\\
4.3975	0.095393\\
4.3939	0.096595\\
4.3903	0.097804\\
4.3867	0.099019\\
4.383	0.10024\\
4.3794	0.10147\\
4.3758	0.10271\\
4.3722	0.10395\\
4.3686	0.1052\\
4.365	0.10646\\
4.3614	0.10772\\
4.3578	0.10899\\
4.3542	0.11026\\
4.3506	0.11155\\
4.347	0.11283\\
4.3434	0.11413\\
4.3398	0.11543\\
4.3362	0.11673\\
4.3325	0.11804\\
4.3289	0.11936\\
4.3253	0.12068\\
4.3217	0.12201\\
4.3181	0.12334\\
4.3145	0.12468\\
4.3109	0.12602\\
4.3073	0.12737\\
4.3037	0.12872\\
4.3001	0.13008\\
4.2965	0.13144\\
4.2929	0.13281\\
4.2893	0.13418\\
4.2856	0.13556\\
4.282	0.13694\\
4.2784	0.13832\\
4.2748	0.13971\\
4.2712	0.14111\\
4.2676	0.1425\\
4.264	0.1439\\
4.2604	0.14531\\
4.2568	0.14672\\
4.2532	0.14813\\
4.2496	0.14955\\
4.246	0.15097\\
4.2424	0.15239\\
4.2388	0.15382\\
4.2351	0.15525\\
4.2315	0.15668\\
4.2279	0.15812\\
4.2243	0.15956\\
4.2207	0.161\\
4.2171	0.16245\\
4.2135	0.1639\\
4.2099	0.16535\\
4.2063	0.16681\\
4.2027	0.16826\\
4.1991	0.16973\\
4.1955	0.17119\\
4.1919	0.17266\\
4.1883	0.17413\\
4.1846	0.1756\\
4.181	0.17707\\
4.1774	0.17855\\
4.1738	0.18003\\
4.1702	0.18152\\
4.1666	0.183\\
4.163	0.18449\\
4.1594	0.18598\\
4.1558	0.18748\\
4.1522	0.18897\\
4.1486	0.19047\\
4.145	0.19198\\
4.1414	0.19348\\
4.1377	0.19499\\
4.1341	0.1965\\
4.1305	0.19802\\
4.1269	0.19954\\
4.1233	0.20106\\
4.1197	0.20258\\
4.1161	0.20411\\
4.1125	0.20564\\
4.1089	0.20718\\
4.1053	0.20871\\
4.1017	0.21026\\
4.0981	0.2118\\
4.0945	0.21336\\
4.0909	0.21491\\
4.0872	0.21647\\
4.0836	0.21803\\
4.08	0.2196\\
4.0764	0.22118\\
4.0728	0.22276\\
4.0692	0.22434\\
4.0656	0.22593\\
4.062	0.22753\\
4.0584	0.22913\\
4.0548	0.23074\\
4.0512	0.23236\\
4.0476	0.23398\\
4.044	0.23561\\
4.0404	0.23725\\
4.0367	0.23889\\
4.0331	0.24054\\
4.0295	0.24221\\
4.0259	0.24388\\
4.0223	0.24556\\
4.0187	0.24725\\
4.0151	0.24895\\
4.0115	0.25066\\
4.0079	0.25238\\
4.0043	0.25411\\
4.0007	0.25585\\
3.9971	0.25761\\
3.9935	0.25938\\
3.9898	0.26116\\
3.9862	0.26296\\
3.9826	0.26476\\
3.979	0.26659\\
3.9754	0.26843\\
3.9718	0.27028\\
3.9682	0.27215\\
3.9646	0.27404\\
3.961	0.27595\\
3.9574	0.27787\\
3.9538	0.27981\\
3.9502	0.28178\\
3.9466	0.28376\\
3.943	0.28576\\
3.9393	0.28778\\
3.9357	0.28983\\
3.9321	0.2919\\
3.9285	0.29399\\
3.9249	0.2961\\
3.9213	0.29824\\
3.9177	0.30041\\
3.9141	0.3026\\
3.9105	0.30482\\
3.9069	0.30706\\
3.9033	0.30934\\
3.8997	0.31164\\
3.8961	0.31398\\
3.8924	0.31634\\
3.8888	0.31874\\
3.8852	0.32117\\
3.8816	0.32363\\
3.878	0.32612\\
3.8744	0.32866\\
3.8708	0.33122\\
3.8672	0.33382\\
3.8636	0.33646\\
3.86	0.33914\\
3.8564	0.34186\\
3.8528	0.34462\\
3.8492	0.34742\\
3.8456	0.35026\\
3.8419	0.35314\\
3.8383	0.35606\\
3.8347	0.35903\\
3.8311	0.36204\\
3.8275	0.3651\\
3.8239	0.36821\\
3.8203	0.37136\\
3.8167	0.37457\\
3.8131	0.37782\\
3.8095	0.38112\\
3.8059	0.38447\\
3.8023	0.38787\\
3.7987	0.39132\\
3.7951	0.39483\\
3.7914	0.39839\\
3.7878	0.40201\\
3.7842	0.40568\\
3.7806	0.4094\\
3.777	0.41319\\
3.7734	0.41703\\
3.7698	0.42092\\
3.7662	0.42488\\
3.7626	0.4289\\
3.759	0.43297\\
3.7554	0.43711\\
3.7518	0.4413\\
3.7482	0.44556\\
3.7445	0.44988\\
3.7409	0.45426\\
3.7373	0.45871\\
3.7337	0.46322\\
3.7301	0.46779\\
3.7265	0.47243\\
3.7229	0.47713\\
3.7193	0.48189\\
3.7157	0.48673\\
3.7121	0.49162\\
3.7085	0.49658\\
3.7049	0.50161\\
3.7013	0.5067\\
3.6977	0.51186\\
3.694	0.51709\\
3.6904	0.52238\\
3.6868	0.52774\\
3.6832	0.53316\\
3.6796	0.53865\\
3.676	0.54421\\
3.6724	0.54983\\
3.6688	0.55552\\
3.6652	0.56127\\
3.6616	0.56708\\
3.658	0.57296\\
3.6544	0.57891\\
3.6508	0.58492\\
3.6472	0.59099\\
3.6435	0.59712\\
3.6399	0.60331\\
3.6363	0.60957\\
3.6327	0.61588\\
3.6291	0.62226\\
3.6255	0.62869\\
3.6219	0.63518\\
3.6183	0.64173\\
3.6147	0.64833\\
3.6111	0.65499\\
3.6075	0.6617\\
3.6039	0.66846\\
3.6003	0.67528\\
3.5966	0.68214\\
3.593	0.68905\\
3.5894	0.69601\\
3.5858	0.70301\\
3.5822	0.71006\\
3.5786	0.71714\\
3.575	0.72427\\
3.5714	0.73144\\
3.5678	0.73864\\
3.5642	0.74587\\
3.5606	0.75314\\
3.557	0.76044\\
3.5534	0.76777\\
3.5498	0.77512\\
3.5461	0.7825\\
3.5425	0.78989\\
3.5389	0.79731\\
3.5353	0.80474\\
3.5317	0.81219\\
3.5281	0.81965\\
3.5245	0.82712\\
3.5209	0.8346\\
3.5173	0.84207\\
3.5137	0.84955\\
3.5101	0.85703\\
3.5065	0.8645\\
3.5029	0.87197\\
3.4992	0.87942\\
3.4956	0.88686\\
3.492	0.89428\\
3.4884	0.90169\\
3.4848	0.90906\\
3.4812	0.91642\\
3.4776	0.92374\\
3.474	0.93102\\
3.4704	0.93828\\
3.4668	0.94549\\
3.4632	0.95265\\
3.4596	0.95977\\
3.456	0.96684\\
3.4524	0.97386\\
3.4487	0.98081\\
3.4451	0.98771\\
3.4415	0.99453\\
3.4379	1.0013\\
3.4343	1.008\\
3.4307	1.0146\\
3.4271	1.0211\\
3.4235	1.0276\\
3.4199	1.0339\\
3.4163	1.0402\\
3.4127	1.0463\\
3.4091	1.0524\\
3.4055	1.0584\\
3.4019	1.0642\\
3.3982	1.07\\
3.3946	1.0756\\
3.391	1.0811\\
3.3874	1.0865\\
3.3838	1.0917\\
3.3802	1.0968\\
3.3766	1.1018\\
3.373	1.1066\\
3.3694	1.1113\\
3.3658	1.1159\\
3.3622	1.1203\\
3.3586	1.1245\\
3.355	1.1285\\
3.3513	1.1324\\
3.3477	1.1362\\
3.3441	1.1397\\
3.3405	1.1431\\
3.3369	1.1463\\
3.3333	1.1492\\
3.3297	1.152\\
3.3261	1.1546\\
3.3225	1.157\\
3.3189	1.1592\\
3.3153	1.1612\\
3.3117	1.1629\\
3.3081	1.1645\\
3.3045	1.1658\\
3.3008	1.1669\\
3.2972	1.1678\\
3.2936	1.1684\\
3.29	1.1688\\
3.2864	1.169\\
3.2828	1.169\\
3.2792	1.1687\\
3.2756	1.1682\\
3.272	1.1675\\
3.2684	1.1666\\
3.2648	1.1654\\
3.2612	1.164\\
3.2576	1.1624\\
3.254	1.1606\\
3.2503	1.1586\\
3.2467	1.1563\\
3.2431	1.1539\\
3.2395	1.1512\\
3.2359	1.1483\\
3.2323	1.1452\\
3.2287	1.142\\
3.2251	1.1385\\
3.2215	1.1348\\
3.2179	1.1309\\
3.2143	1.1269\\
3.2107	1.1226\\
3.2071	1.1182\\
3.2034	1.1136\\
3.1998	1.1088\\
3.1962	1.1038\\
3.1926	1.0987\\
3.189	1.0934\\
3.1854	1.0879\\
3.1818	1.0822\\
3.1782	1.0764\\
3.1746	1.0704\\
3.171	1.0643\\
3.1674	1.058\\
3.1638	1.0515\\
3.1602	1.045\\
3.1566	1.0382\\
3.1529	1.0314\\
3.1493	1.0243\\
3.1457	1.0172\\
3.1421	1.0099\\
3.1385	1.0025\\
3.1349	0.995\\
3.1313	0.98736\\
3.1277	0.9796\\
3.1241	0.97173\\
3.1205	0.96375\\
3.1169	0.95567\\
3.1133	0.94748\\
3.1097	0.9392\\
3.106	0.93083\\
3.1024	0.92236\\
3.0988	0.91381\\
3.0952	0.90518\\
3.0916	0.89648\\
3.088	0.88769\\
3.0844	0.87884\\
3.0808	0.86992\\
3.0772	0.86094\\
3.0736	0.8519\\
3.07	0.84281\\
3.0664	0.83367\\
3.0628	0.82448\\
3.0592	0.81525\\
3.0555	0.80597\\
3.0519	0.79667\\
3.0483	0.78733\\
3.0447	0.77796\\
3.0411	0.76857\\
3.0375	0.75916\\
3.0339	0.74974\\
3.0303	0.7403\\
3.0267	0.73085\\
3.0231	0.7214\\
3.0195	0.71194\\
3.0159	0.70249\\
3.0123	0.69304\\
3.0087	0.68359\\
3.005	0.67416\\
3.0014	0.66474\\
2.9978	0.65535\\
2.9942	0.64597\\
2.9906	0.63661\\
2.987	0.62729\\
2.9834	0.61799\\
2.9798	0.60872\\
2.9762	0.59949\\
2.9726	0.5903\\
2.969	0.58115\\
2.9654	0.57205\\
2.9618	0.56299\\
2.9581	0.55398\\
2.9545	0.54502\\
2.9509	0.53611\\
2.9473	0.52726\\
2.9437	0.51847\\
2.9401	0.50974\\
2.9365	0.50107\\
2.9329	0.49247\\
2.9293	0.48393\\
2.9257	0.47546\\
2.9221	0.46706\\
2.9185	0.45873\\
2.9149	0.45048\\
2.9113	0.4423\\
2.9076	0.4342\\
2.904	0.42617\\
2.9004	0.41823\\
2.8968	0.41036\\
2.8932	0.40258\\
2.8896	0.39488\\
2.886	0.38726\\
2.8824	0.37973\\
2.8788	0.37229\\
2.8752	0.36493\\
2.8716	0.35766\\
2.868	0.35048\\
2.8644	0.34338\\
2.8607	0.33638\\
2.8571	0.32947\\
2.8535	0.32265\\
2.8499	0.31592\\
2.8463	0.30928\\
2.8427	0.30273\\
2.8391	0.29627\\
2.8355	0.28991\\
2.8319	0.28364\\
2.8283	0.27746\\
2.8247	0.27138\\
2.8211	0.26539\\
2.8175	0.25949\\
2.8139	0.25368\\
2.8102	0.24797\\
2.8066	0.24235\\
2.803	0.23682\\
2.7994	0.23138\\
2.7958	0.22603\\
2.7922	0.22077\\
2.7886	0.21561\\
2.785	0.21053\\
2.7814	0.20554\\
2.7778	0.20065\\
2.7742	0.19584\\
2.7706	0.19112\\
2.767	0.18648\\
2.7634	0.18193\\
2.7597	0.17747\\
2.7561	0.1731\\
2.7525	0.1688\\
2.7489	0.1646\\
2.7453	0.16047\\
2.7417	0.15643\\
2.7381	0.15247\\
2.7345	0.14858\\
2.7309	0.14478\\
2.7273	0.14106\\
2.7237	0.13741\\
2.7201	0.13384\\
2.7165	0.13035\\
2.7128	0.12693\\
2.7092	0.12358\\
2.7056	0.12031\\
2.702	0.11711\\
2.6984	0.11398\\
2.6948	0.11092\\
2.6912	0.10793\\
2.6876	0.10501\\
2.684	0.10216\\
2.6804	0.099365\\
2.6768	0.09664\\
2.6732	0.093979\\
2.6696	0.091381\\
2.666	0.088844\\
2.6623	0.086368\\
2.6587	0.083951\\
2.6551	0.081593\\
2.6515	0.079293\\
2.6479	0.077049\\
2.6443	0.074861\\
2.6407	0.072727\\
2.6371	0.070647\\
2.6335	0.068619\\
2.6299	0.066642\\
2.6263	0.064716\\
2.6227	0.06284\\
2.6191	0.061011\\
2.6155	0.059231\\
2.6118	0.057497\\
2.6082	0.055808\\
2.6046	0.054164\\
2.601	0.052564\\
2.5974	0.051006\\
2.5938	0.049491\\
2.5902	0.048016\\
2.5866	0.046581\\
2.583	0.045185\\
2.5794	0.043828\\
2.5758	0.042508\\
2.5722	0.041224\\
2.5686	0.039976\\
2.5649	0.038763\\
2.5613	0.037584\\
2.5577	0.036438\\
2.5541	0.035325\\
2.5505	0.034243\\
2.5469	0.033192\\
2.5433	0.032171\\
2.5397	0.031179\\
2.5361	0.030216\\
2.5325	0.029281\\
2.5289	0.028374\\
2.5253	0.027492\\
2.5217	0.026637\\
2.5181	0.025807\\
2.5144	0.025001\\
2.5108	0.024219\\
2.5072	0.02346\\
2.5036	0.022724\\
2.5	0.02201\\
2.4964	0.021317\\
2.4928	0.020645\\
2.4892	0.019994\\
2.4856	0.019362\\
2.482	0.018749\\
2.4784	0.018155\\
2.4748	0.017579\\
2.4712	0.017021\\
2.4675	0.01648\\
2.4639	0.015956\\
2.4603	0.015447\\
2.4567	0.014955\\
2.4531	0.014477\\
2.4495	0.014015\\
2.4459	0.013567\\
2.4423	0.013133\\
2.4387	0.012712\\
2.4351	0.012305\\
2.4315	0.01191\\
2.4279	0.011528\\
2.4243	0.011158\\
2.4207	0.010799\\
2.417	0.010452\\
2.4134	0.010116\\
2.4098	0.0097907\\
2.4062	0.0094756\\
2.4026	0.0091705\\
2.399	0.0088751\\
2.3954	0.0085892\\
2.3918	0.0083124\\
2.3882	0.0080444\\
2.3846	0.0077851\\
2.381	0.007534\\
2.3774	0.007291\\
2.3738	0.0070559\\
2.3702	0.0068283\\
2.3665	0.006608\\
2.3629	0.0063949\\
2.3593	0.0061886\\
2.3557	0.005989\\
2.3521	0.0057959\\
2.3485	0.005609\\
2.3449	0.0054282\\
2.3413	0.0052533\\
2.3377	0.005084\\
2.3341	0.0049202\\
2.3305	0.0047618\\
2.3269	0.0046085\\
2.3233	0.0044602\\
2.3196	0.0043167\\
2.316	0.0041779\\
2.3124	0.0040436\\
2.3088	0.0039137\\
2.3052	0.0037881\\
2.3016	0.0036665\\
2.298	0.0035489\\
2.2944	0.0034352\\
2.2908	0.0033252\\
2.2872	0.0032187\\
2.2836	0.0031158\\
2.28	0.0030162\\
2.2764	0.0029198\\
2.2728	0.0028266\\
2.2691	0.0027365\\
2.2655	0.0026493\\
2.2619	0.002565\\
2.2583	0.0024834\\
2.2547	0.0024044\\
2.2511	0.0023281\\
2.2475	0.0022542\\
2.2439	0.0021828\\
2.2403	0.0021137\\
2.2367	0.0020468\\
2.2331	0.0019822\\
2.2295	0.0019196\\
2.2259	0.0018591\\
2.2223	0.0018005\\
2.2186	0.0017439\\
2.215	0.0016891\\
2.2114	0.0016361\\
2.2078	0.0015848\\
2.2042	0.0015351\\
2.2006	0.0014871\\
2.197	0.0014406\\
2.1934	0.0013957\\
2.1898	0.0013522\\
2.1862	0.0013101\\
2.1826	0.0012693\\
2.179	0.0012299\\
2.1754	0.0011918\\
2.1717	0.0011548\\
2.1681	0.0011191\\
2.1645	0.0010845\\
2.1609	0.001051\\
2.1573	0.0010186\\
2.1537	0.00098728\\
2.1501	0.00095692\\
2.1465	0.00092753\\
2.1429	0.00089908\\
2.1393	0.00087154\\
2.1357	0.00084487\\
2.1321	0.00081905\\
2.1285	0.00079405\\
2.1249	0.00076984\\
2.1212	0.0007464\\
2.1176	0.0007237\\
2.114	0.00070172\\
2.1104	0.00068043\\
2.1068	0.00065981\\
2.1032	0.00063984\\
2.0996	0.00062049\\
2.096	0.00060175\\
2.0924	0.0005836\\
2.0888	0.00056601\\
2.0852	0.00054897\\
2.0816	0.00053246\\
2.078	0.00051647\\
2.0743	0.00050097\\
2.0707	0.00048595\\
2.0671	0.00047139\\
2.0635	0.00045729\\
2.0599	0.00044362\\
2.0563	0.00043037\\
2.0527	0.00041753\\
2.0491	0.00040508\\
2.0455	0.00039301\\
2.0419	0.00038131\\
2.0383	0.00036997\\
2.0347	0.00035897\\
2.0311	0.00034831\\
2.0275	0.00033798\\
2.0238	0.00032795\\
2.0202	0.00031823\\
2.0166	0.00030881\\
2.013	0.00029967\\
2.0094	0.0002908\\
2.0058	0.0002822\\
2.0022	0.00027386\\
1.9986	0.00026577\\
1.995	0.00025792\\
1.9914	0.00025031\\
1.9878	0.00024292\\
1.9842	0.00023576\\
1.9806	0.00022881\\
1.977	0.00022206\\
1.9733	0.00021552\\
1.9697	0.00020917\\
1.9661	0.00020301\\
1.9625	0.00019703\\
1.9589	0.00019123\\
1.9553	0.0001856\\
1.9517	0.00018013\\
1.9481	0.00017483\\
1.9445	0.00016968\\
1.9409	0.00016469\\
1.9373	0.00015984\\
1.9337	0.00015513\\
1.9301	0.00015057\\
1.9264	0.00014613\\
1.9228	0.00014183\\
1.9192	0.00013765\\
1.9156	0.00013359\\
1.912	0.00012966\\
1.9084	0.00012584\\
1.9048	0.00012212\\
1.9012	0.00011852\\
1.8976	0.00011502\\
1.894	0.00011163\\
1.8904	0.00010833\\
1.8868	0.00010513\\
1.8832	0.00010202\\
1.8796	9.9008e-05\\
1.8759	9.6079e-05\\
1.8723	9.3235e-05\\
1.8687	9.0473e-05\\
1.8651	8.7792e-05\\
1.8615	8.5189e-05\\
1.8579	8.2661e-05\\
1.8543	8.0207e-05\\
1.8507	7.7824e-05\\
1.8471	7.551e-05\\
1.8435	7.3263e-05\\
1.8399	7.1082e-05\\
1.8363	6.8964e-05\\
1.8327	6.6907e-05\\
1.8291	6.491e-05\\
1.8254	6.2972e-05\\
1.8218	6.1089e-05\\
1.8182	5.9261e-05\\
1.8146	5.7487e-05\\
1.811	5.5764e-05\\
1.8074	5.4091e-05\\
1.8038	5.2467e-05\\
1.8002	5.089e-05\\
1.7966	4.9359e-05\\
1.793	4.7872e-05\\
1.7894	4.6429e-05\\
1.7858	4.5028e-05\\
1.7822	4.3668e-05\\
1.7785	4.2348e-05\\
1.7749	4.1066e-05\\
1.7713	3.9821e-05\\
1.7677	3.8613e-05\\
1.7641	3.7441e-05\\
1.7605	3.6302e-05\\
1.7569	3.5197e-05\\
1.7533	3.4125e-05\\
1.7497	3.3084e-05\\
1.7461	3.2073e-05\\
1.7425	3.1093e-05\\
1.7389	3.0141e-05\\
1.7353	2.9217e-05\\
1.7317	2.832e-05\\
1.728	2.745e-05\\
1.7244	2.6605e-05\\
1.7208	2.5786e-05\\
1.7172	2.4991e-05\\
1.7136	2.4219e-05\\
1.71	2.347e-05\\
1.7064	2.2743e-05\\
}--cycle;
\addplot[ybar interval, fill=black, fill opacity=0.35, area legend, draw=none] table[row sep=crcr, x=Lower, y=Count] {%
Lower	Upper	Count\\
1.7064	1.9836	0\\
1.9836	2.2608	0.019239\\
2.2608	2.538	0.15391\\
2.538	2.8152	0.31745\\
2.8152	3.0925	0.81767\\
3.0925	3.3697	0.84653\\
3.3697	3.6469	0.75033\\
3.6469	3.9241	0.37517\\
3.9241	4.2013	0.20201\\
4.2013	4.4785	0.038479\\
4.4785	4.7557	0.057718\\
4.7557	5.0329	0.0096197\\
5.0329	5.3101	0.019239\\
5.3101	5.3101	0.019239\\
};
\addlegendentry{Istogramma reale}

\addplot [color=black]
  table[row sep=crcr]{%
1.7064	0.00086709\\
1.8146	0.0022528\\
1.8904	0.0042154\\
1.9517	0.0068258\\
2.0022	0.0099809\\
2.0455	0.013656\\
2.0852	0.018023\\
2.1212	0.023004\\
2.1537	0.028463\\
2.1862	0.034996\\
2.215	0.041828\\
2.2439	0.049744\\
2.2728	0.058864\\
2.298	0.067926\\
2.3233	0.078083\\
2.3485	0.089415\\
2.3738	0.102\\
2.399	0.11591\\
2.4243	0.13122\\
2.4495	0.14798\\
2.4712	0.16354\\
2.4964	0.18312\\
2.5217	0.20426\\
2.5469	0.22697\\
2.5722	0.25124\\
2.5974	0.27705\\
2.6227	0.30433\\
2.6515	0.33724\\
2.6804	0.37185\\
2.7128	0.41258\\
2.7489	0.45966\\
2.7958	0.52286\\
2.9004	0.66466\\
2.9329	0.70659\\
2.9618	0.74215\\
2.987	0.77158\\
3.0087	0.79533\\
3.0303	0.81752\\
3.0519	0.838\\
3.07	0.85364\\
3.088	0.8679\\
3.106	0.88069\\
3.1241	0.89196\\
3.1385	0.89983\\
3.1529	0.90666\\
3.1674	0.91241\\
3.1818	0.91708\\
3.1962	0.92065\\
3.2107	0.92309\\
3.2251	0.92442\\
3.2395	0.92462\\
3.254	0.9237\\
3.2684	0.92166\\
3.2828	0.91851\\
3.2972	0.91427\\
3.3117	0.90895\\
3.3261	0.90257\\
3.3405	0.89517\\
3.355	0.88676\\
3.373	0.8749\\
3.391	0.86159\\
3.4091	0.84693\\
3.4271	0.83098\\
3.4487	0.81029\\
3.4704	0.78805\\
3.4956	0.76039\\
3.5209	0.73115\\
3.5498	0.69617\\
3.5858	0.65074\\
3.6472	0.57138\\
3.7085	0.49257\\
3.7482	0.44339\\
3.7806	0.40482\\
3.8095	0.37207\\
3.8383	0.34094\\
3.8636	0.31513\\
3.8888	0.29073\\
3.9141	0.26777\\
3.9393	0.24628\\
3.9646	0.22625\\
3.9898	0.20766\\
4.0151	0.1905\\
4.0404	0.17471\\
4.0656	0.16023\\
4.0909	0.14701\\
4.1161	0.13498\\
4.1414	0.12406\\
4.1702	0.11284\\
4.1991	0.10285\\
4.2279	0.09398\\
4.2604	0.085187\\
4.2929	0.077503\\
4.3289	0.070088\\
4.3686	0.063083\\
4.4119	0.05657\\
4.4624	0.050147\\
4.5201	0.04398\\
4.5887	0.03782\\
4.6716	0.031519\\
4.769	0.02522\\
4.8809	0.019087\\
4.9963	0.013844\\
5.1117	0.0096573\\
5.2344	0.0062955\\
5.3101	0.0047195\\
};
\addlegendentry{Adatt. BLN reale}

            % \input{../../../.tex/20/KS/SizeDistributionFittingsfig3.tex}
    
            \nextgroupplot[%
                xmin=2.0755,xmax=5.4719,
                % xmin=1.9138,xmax=5.4719,
                % ymin=0,ymax=1.4,
            ] % This file was created by matlab2tikz.
%
\definecolor{mycolor1}{rgb}{0.83529,0.36863,0.00000}%
%
\addplot[ybar interval, fill=mycolor1, fill opacity=0.35, area legend, draw=none] table[row sep=crcr, x=Lower, y=Count] {%
Lower	Upper	Count\\
2.0755	2.3244	0.00075022\\
2.3244	2.5732	0.0097529\\
2.5732	2.822	0.099566\\
2.822	3.0708	0.55141\\
3.0708	3.3196	1.1455\\
3.3196	3.5684	1.0567\\
3.5684	3.8172	0.55099\\
3.8172	4.0661	0.29344\\
4.0661	4.3149	0.17352\\
4.3149	4.5637	0.10042\\
4.5637	4.8125	0.03451\\
4.8125	5.0613	0.002465\\
5.0613	5.3101	0\\
5.3101	5.3101	0\\
};
\addlegendentry{Istogramma appr.}

\addplot [color=mycolor1, only marks, every error bar/.append style={opacity=0.45}, mark=*, mark size=0pt, draw=none, forget plot]
 plot [error bars/.cd, y dir=both, y explicit, error bar style={line width=1pt, color=mycolor1}, error mark options={mark=none,mark size=0pt}]
 table[row sep=crcr, y error plus index=2, y error minus index=3]{%
2.2	0.00075022	0.00054532	0.00054532\\
2.4488	0.0097529	0.0020071	0.0020071\\
2.6976	0.099566	0.0056491	0.0056491\\
2.9464	0.55141	0.011062	0.011062\\
3.1952	1.1455	0.01749	0.01749\\
3.444	1.0567	0.014334	0.014334\\
3.6928	0.55099	0.012038	0.012038\\
3.9417	0.29344	0.0079672	0.0079672\\
4.1905	0.17352	0.0058038	0.0058038\\
4.4393	0.10042	0.0042279	0.0042279\\
4.6881	0.03451	0.0022534	0.0022534\\
4.9369	0.002465	0.00089944	0.00089944\\
5.1857	0	0	0\\
};
\addplot [color=mycolor1]
  table[row sep=crcr]{%
2.0755	0.00032936\\
2.2407	0.0012963\\
2.3281	0.0028174\\
2.3864	0.0048427\\
2.4317	0.007457\\
2.4673	0.010508\\
2.4997	0.014366\\
2.5256	0.018439\\
2.5515	0.02363\\
2.5742	0.029301\\
2.5936	0.035166\\
2.613	0.042115\\
2.6325	0.050312\\
2.6486	0.058227\\
2.6648	0.067246\\
2.681	0.077486\\
2.6972	0.089071\\
2.7134	0.10213\\
2.7296	0.11678\\
2.7458	0.13316\\
2.762	0.15139\\
2.7782	0.17159\\
2.7943	0.19387\\
2.8105	0.21834\\
2.8267	0.24507\\
2.8429	0.27413\\
2.8591	0.30556\\
2.8753	0.33937\\
2.8915	0.37554\\
2.9077	0.41403\\
2.9271	0.46312\\
2.9465	0.51517\\
2.9692	0.57918\\
2.9951	0.65584\\
3.0307	0.76502\\
3.0793	0.91396\\
3.1019	0.98043\\
3.1214	1.0345\\
3.1376	1.0768\\
3.1537	1.1162\\
3.1667	1.1453\\
3.1796	1.172\\
3.1926	1.1961\\
3.2055	1.2175\\
3.2153	1.2317\\
3.225	1.2442\\
3.2347	1.255\\
3.2444	1.264\\
3.2541	1.2712\\
3.2638	1.2766\\
3.2703	1.2792\\
3.2768	1.281\\
3.2833	1.282\\
3.2897	1.2822\\
3.2962	1.2816\\
3.3027	1.2802\\
3.3092	1.278\\
3.3189	1.2732\\
3.3286	1.2668\\
3.3383	1.2586\\
3.348	1.2489\\
3.3577	1.2377\\
3.3674	1.2249\\
3.3804	1.2058\\
3.3933	1.1843\\
3.4063	1.1608\\
3.4225	1.1288\\
3.4387	1.0942\\
3.4581	1.0501\\
3.4808	0.99591\\
3.5164	0.90758\\
3.5649	0.78753\\
3.5908	0.72639\\
3.6135	0.67558\\
3.6329	0.63441\\
3.6524	0.59568\\
3.6718	0.55952\\
3.688	0.5314\\
3.7042	0.50513\\
3.7204	0.48069\\
3.7366	0.45805\\
3.7527	0.43713\\
3.7689	0.41787\\
3.7851	0.40017\\
3.8013	0.38394\\
3.8175	0.36907\\
3.8337	0.35545\\
3.8531	0.3406\\
3.8725	0.3272\\
3.892	0.31505\\
3.9146	0.30222\\
3.9373	0.29058\\
3.9632	0.27842\\
3.9956	0.26445\\
4.0377	0.24755\\
4.2319	0.17104\\
4.3647	0.11729\\
4.4133	0.098923\\
4.4553	0.084097\\
4.4942	0.071508\\
4.5331	0.060088\\
4.5687	0.050705\\
4.6043	0.042373\\
4.6399	0.035073\\
4.6788	0.028235\\
4.7176	0.022488\\
4.7597	0.017371\\
4.805	0.012986\\
4.8536	0.0093766\\
4.9086	0.0063781\\
4.9734	0.0039718\\
5.0511	0.0021917\\
5.1547	0.00095392\\
5.3101	0.00025468\\
};
\addlegendentry{Adatt. BLN appr.}


\addplot[area legend, draw=none, fill=mycolor1, fill opacity=0.15, forget plot]
table[row sep=crcr] {%
x	y\\
2.0755	0.00043465\\
2.0788	0.00044574\\
2.082	0.0004571\\
2.0853	0.00046876\\
2.0885	0.00048071\\
2.0917	0.00049297\\
2.095	0.00050554\\
2.0982	0.00051843\\
2.1014	0.00053165\\
2.1047	0.00054521\\
2.1079	0.00055911\\
2.1112	0.00057338\\
2.1144	0.00058801\\
2.1176	0.00060302\\
2.1209	0.00061841\\
2.1241	0.00063421\\
2.1274	0.00065041\\
2.1306	0.00066704\\
2.1338	0.00068409\\
2.1371	0.0007016\\
2.1403	0.00071956\\
2.1435	0.00073799\\
2.1468	0.0007569\\
2.15	0.00077631\\
2.1533	0.00079624\\
2.1565	0.00081669\\
2.1597	0.00083768\\
2.163	0.00085923\\
2.1662	0.00088135\\
2.1694	0.00090406\\
2.1727	0.00092738\\
2.1759	0.00095133\\
2.1792	0.00097592\\
2.1824	0.0010012\\
2.1856	0.0010271\\
2.1889	0.0010537\\
2.1921	0.0010811\\
2.1953	0.0011092\\
2.1986	0.0011381\\
2.2018	0.0011678\\
2.2051	0.0011982\\
2.2083	0.0012296\\
2.2115	0.0012617\\
2.2148	0.0012948\\
2.218	0.0013288\\
2.2212	0.0013638\\
2.2245	0.0013997\\
2.2277	0.0014366\\
2.231	0.0014746\\
2.2342	0.0015136\\
2.2374	0.0015538\\
2.2407	0.001595\\
2.2439	0.0016375\\
2.2472	0.0016812\\
2.2504	0.0017261\\
2.2536	0.0017723\\
2.2569	0.0018198\\
2.2601	0.0018687\\
2.2633	0.001919\\
2.2666	0.0019708\\
2.2698	0.0020241\\
2.2731	0.002079\\
2.2763	0.0021354\\
2.2795	0.0021935\\
2.2828	0.0022533\\
2.286	0.0023149\\
2.2892	0.0023783\\
2.2925	0.0024436\\
2.2957	0.0025108\\
2.299	0.0025801\\
2.3022	0.0026514\\
2.3054	0.0027248\\
2.3087	0.0028005\\
2.3119	0.0028784\\
2.3151	0.0029587\\
2.3184	0.0030415\\
2.3216	0.0031267\\
2.3249	0.0032145\\
2.3281	0.003305\\
2.3313	0.0033983\\
2.3346	0.0034944\\
2.3378	0.0035935\\
2.341	0.0036957\\
2.3443	0.0038009\\
2.3475	0.0039095\\
2.3508	0.0040214\\
2.354	0.0041367\\
2.3572	0.0042556\\
2.3605	0.0043783\\
2.3637	0.0045047\\
2.367	0.0046351\\
2.3702	0.0047696\\
2.3734	0.0049083\\
2.3767	0.0050513\\
2.3799	0.0051988\\
2.3831	0.005351\\
2.3864	0.0055079\\
2.3896	0.0056698\\
2.3929	0.0058367\\
2.3961	0.006009\\
2.3993	0.0061866\\
2.4026	0.0063699\\
2.4058	0.006559\\
2.409	0.006754\\
2.4123	0.0069552\\
2.4155	0.0071628\\
2.4188	0.007377\\
2.422	0.0075979\\
2.4252	0.0078259\\
2.4285	0.008061\\
2.4317	0.0083037\\
2.4349	0.008554\\
2.4382	0.0088122\\
2.4414	0.0090786\\
2.4447	0.0093534\\
2.4479	0.009637\\
2.4511	0.0099295\\
2.4544	0.010231\\
2.4576	0.010542\\
2.4608	0.010864\\
2.4641	0.011195\\
2.4673	0.011536\\
2.4706	0.011889\\
2.4738	0.012252\\
2.477	0.012627\\
2.4803	0.013014\\
2.4835	0.013413\\
2.4868	0.013824\\
2.49	0.014248\\
2.4932	0.014685\\
2.4965	0.015136\\
2.4997	0.015601\\
2.5029	0.016081\\
2.5062	0.016575\\
2.5094	0.017084\\
2.5127	0.017609\\
2.5159	0.018151\\
2.5191	0.018708\\
2.5224	0.019283\\
2.5256	0.019876\\
2.5288	0.020486\\
2.5321	0.021116\\
2.5353	0.021764\\
2.5386	0.022431\\
2.5418	0.023119\\
2.545	0.023828\\
2.5483	0.024557\\
2.5515	0.025309\\
2.5547	0.026082\\
2.558	0.026879\\
2.5612	0.027699\\
2.5645	0.028544\\
2.5677	0.029413\\
2.5709	0.030307\\
2.5742	0.031228\\
2.5774	0.032175\\
2.5806	0.03315\\
2.5839	0.034152\\
2.5871	0.035184\\
2.5904	0.036245\\
2.5936	0.037336\\
2.5968	0.038457\\
2.6001	0.039611\\
2.6033	0.040797\\
2.6066	0.042016\\
2.6098	0.043269\\
2.613	0.044556\\
2.6163	0.045879\\
2.6195	0.047239\\
2.6227	0.048635\\
2.626	0.050069\\
2.6292	0.051542\\
2.6325	0.053055\\
2.6357	0.054608\\
2.6389	0.056203\\
2.6422	0.057839\\
2.6454	0.059519\\
2.6486	0.061243\\
2.6519	0.063011\\
2.6551	0.064825\\
2.6584	0.066686\\
2.6616	0.068595\\
2.6648	0.070552\\
2.6681	0.072559\\
2.6713	0.074616\\
2.6745	0.076724\\
2.6778	0.078885\\
2.681	0.081099\\
2.6843	0.083368\\
2.6875	0.085691\\
2.6907	0.088071\\
2.694	0.090509\\
2.6972	0.093004\\
2.7004	0.095558\\
2.7037	0.098173\\
2.7069	0.10085\\
2.7102	0.10359\\
2.7134	0.10639\\
2.7166	0.10925\\
2.7199	0.11218\\
2.7231	0.11518\\
2.7263	0.11824\\
2.7296	0.12138\\
2.7328	0.12458\\
2.7361	0.12785\\
2.7393	0.13119\\
2.7425	0.1346\\
2.7458	0.13809\\
2.749	0.14165\\
2.7523	0.14528\\
2.7555	0.14899\\
2.7587	0.15278\\
2.762	0.15664\\
2.7652	0.16059\\
2.7684	0.16461\\
2.7717	0.16871\\
2.7749	0.17289\\
2.7782	0.17716\\
2.7814	0.18151\\
2.7846	0.18594\\
2.7879	0.19045\\
2.7911	0.19505\\
2.7943	0.19973\\
2.7976	0.2045\\
2.8008	0.20936\\
2.8041	0.21431\\
2.8073	0.21934\\
2.8105	0.22446\\
2.8138	0.22967\\
2.817	0.23497\\
2.8202	0.24036\\
2.8235	0.24584\\
2.8267	0.25141\\
2.83	0.25708\\
2.8332	0.26283\\
2.8364	0.26868\\
2.8397	0.27462\\
2.8429	0.28065\\
2.8461	0.28678\\
2.8494	0.29299\\
2.8526	0.2993\\
2.8559	0.30571\\
2.8591	0.31221\\
2.8623	0.3188\\
2.8656	0.32548\\
2.8688	0.33226\\
2.8721	0.33913\\
2.8753	0.34609\\
2.8785	0.35314\\
2.8818	0.36029\\
2.885	0.36753\\
2.8882	0.37485\\
2.8915	0.38227\\
2.8947	0.38978\\
2.898	0.39738\\
2.9012	0.40507\\
2.9044	0.41285\\
2.9077	0.42071\\
2.9109	0.42866\\
2.9141	0.4367\\
2.9174	0.44482\\
2.9206	0.45303\\
2.9239	0.46131\\
2.9271	0.46969\\
2.9303	0.47814\\
2.9336	0.48667\\
2.9368	0.49528\\
2.94	0.50396\\
2.9433	0.51272\\
2.9465	0.52156\\
2.9498	0.53047\\
2.953	0.53945\\
2.9562	0.5485\\
2.9595	0.55762\\
2.9627	0.5668\\
2.9659	0.57605\\
2.9692	0.58536\\
2.9724	0.59473\\
2.9757	0.60416\\
2.9789	0.61365\\
2.9821	0.62319\\
2.9854	0.63278\\
2.9886	0.64243\\
2.9919	0.65212\\
2.9951	0.66186\\
2.9983	0.67164\\
3.0016	0.68146\\
3.0048	0.69132\\
3.008	0.70122\\
3.0113	0.71115\\
3.0145	0.72111\\
3.0178	0.7311\\
3.021	0.74111\\
3.0242	0.75114\\
3.0275	0.76119\\
3.0307	0.77126\\
3.0339	0.78133\\
3.0372	0.79142\\
3.0404	0.80151\\
3.0437	0.8116\\
3.0469	0.82168\\
3.0501	0.83177\\
3.0534	0.84184\\
3.0566	0.8519\\
3.0598	0.86193\\
3.0631	0.87195\\
3.0663	0.88194\\
3.0696	0.89191\\
3.0728	0.90183\\
3.076	0.91172\\
3.0793	0.92157\\
3.0825	0.93137\\
3.0857	0.94112\\
3.089	0.95082\\
3.0922	0.96046\\
3.0955	0.97003\\
3.0987	0.97954\\
3.1019	0.98897\\
3.1052	0.99833\\
3.1084	1.0076\\
3.1117	1.0168\\
3.1149	1.0259\\
3.1181	1.0349\\
3.1214	1.0438\\
3.1246	1.0527\\
3.1278	1.0614\\
3.1311	1.07\\
3.1343	1.0785\\
3.1376	1.0868\\
3.1408	1.0951\\
3.144	1.1032\\
3.1473	1.1112\\
3.1505	1.1191\\
3.1537	1.1268\\
3.157	1.1344\\
3.1602	1.1419\\
3.1635	1.1492\\
3.1667	1.1563\\
3.1699	1.1633\\
3.1732	1.1702\\
3.1764	1.1769\\
3.1796	1.1834\\
3.1829	1.1897\\
3.1861	1.1959\\
3.1894	1.202\\
3.1926	1.2078\\
3.1958	1.2135\\
3.1991	1.219\\
3.2023	1.2243\\
3.2055	1.2294\\
3.2088	1.2343\\
3.212	1.2391\\
3.2153	1.2436\\
3.2185	1.248\\
3.2217	1.2522\\
3.225	1.2562\\
3.2282	1.2599\\
3.2315	1.2635\\
3.2347	1.2669\\
3.2379	1.2701\\
3.2412	1.2731\\
3.2444	1.2758\\
3.2476	1.2784\\
3.2509	1.2808\\
3.2541	1.2829\\
3.2574	1.2849\\
3.2606	1.2866\\
3.2638	1.2882\\
3.2671	1.2895\\
3.2703	1.2907\\
3.2735	1.2916\\
3.2768	1.2923\\
3.28	1.2928\\
3.2833	1.2931\\
3.2865	1.2932\\
3.2897	1.2931\\
3.293	1.2928\\
3.2962	1.2923\\
3.2994	1.2916\\
3.3027	1.2908\\
3.3059	1.2897\\
3.3092	1.2884\\
3.3124	1.2869\\
3.3156	1.2852\\
3.3189	1.2833\\
3.3221	1.2813\\
3.3253	1.279\\
3.3286	1.2766\\
3.3318	1.274\\
3.3351	1.2712\\
3.3383	1.2682\\
3.3415	1.2651\\
3.3448	1.2618\\
3.348	1.2583\\
3.3513	1.2547\\
3.3545	1.2509\\
3.3577	1.2469\\
3.361	1.2428\\
3.3642	1.2385\\
3.3674	1.234\\
3.3707	1.2294\\
3.3739	1.2247\\
3.3772	1.2198\\
3.3804	1.2148\\
3.3836	1.2097\\
3.3869	1.2044\\
3.3901	1.199\\
3.3933	1.1934\\
3.3966	1.1878\\
3.3998	1.182\\
3.4031	1.1761\\
3.4063	1.17\\
3.4095	1.1639\\
3.4128	1.1576\\
3.416	1.1513\\
3.4192	1.1448\\
3.4225	1.1383\\
3.4257	1.1316\\
3.429	1.1249\\
3.4322	1.118\\
3.4354	1.1111\\
3.4387	1.1041\\
3.4419	1.097\\
3.4451	1.0898\\
3.4484	1.0826\\
3.4516	1.0753\\
3.4549	1.0679\\
3.4581	1.0605\\
3.4613	1.053\\
3.4646	1.0454\\
3.4678	1.0378\\
3.471	1.0301\\
3.4743	1.0224\\
3.4775	1.0146\\
3.4808	1.0068\\
3.484	0.99895\\
3.4872	0.99106\\
3.4905	0.98314\\
3.4937	0.97519\\
3.497	0.96722\\
3.5002	0.95922\\
3.5034	0.9512\\
3.5067	0.94316\\
3.5099	0.9351\\
3.5131	0.92704\\
3.5164	0.91897\\
3.5196	0.91089\\
3.5229	0.9028\\
3.5261	0.89472\\
3.5293	0.88664\\
3.5326	0.87856\\
3.5358	0.8705\\
3.539	0.86244\\
3.5423	0.85439\\
3.5455	0.84637\\
3.5488	0.83836\\
3.552	0.83037\\
3.5552	0.8224\\
3.5585	0.81446\\
3.5617	0.80654\\
3.5649	0.79866\\
3.5682	0.79081\\
3.5714	0.78299\\
3.5747	0.77521\\
3.5779	0.76747\\
3.5811	0.75977\\
3.5844	0.75211\\
3.5876	0.74449\\
3.5908	0.73692\\
3.5941	0.7294\\
3.5973	0.72193\\
3.6006	0.71451\\
3.6038	0.70714\\
3.607	0.69983\\
3.6103	0.69257\\
3.6135	0.68537\\
3.6168	0.67823\\
3.62	0.67114\\
3.6232	0.66412\\
3.6265	0.65716\\
3.6297	0.65027\\
3.6329	0.64343\\
3.6362	0.63667\\
3.6394	0.62997\\
3.6427	0.62334\\
3.6459	0.61678\\
3.6491	0.61028\\
3.6524	0.60386\\
3.6556	0.59751\\
3.6588	0.59123\\
3.6621	0.58502\\
3.6653	0.57888\\
3.6686	0.57282\\
3.6718	0.56683\\
3.675	0.56091\\
3.6783	0.55507\\
3.6815	0.54931\\
3.6847	0.54362\\
3.688	0.53801\\
3.6912	0.53247\\
3.6945	0.52701\\
3.6977	0.52163\\
3.7009	0.51632\\
3.7042	0.51109\\
3.7074	0.50594\\
3.7106	0.50086\\
3.7139	0.49587\\
3.7171	0.49094\\
3.7204	0.4861\\
3.7236	0.48133\\
3.7268	0.47664\\
3.7301	0.47203\\
3.7333	0.46749\\
3.7366	0.46303\\
3.7398	0.45865\\
3.743	0.45434\\
3.7463	0.45011\\
3.7495	0.44595\\
3.7527	0.44186\\
3.756	0.43785\\
3.7592	0.43391\\
3.7625	0.43004\\
3.7657	0.42624\\
3.7689	0.42251\\
3.7722	0.41885\\
3.7754	0.41526\\
3.7786	0.41173\\
3.7819	0.40827\\
3.7851	0.40488\\
3.7884	0.40155\\
3.7916	0.39828\\
3.7948	0.39507\\
3.7981	0.39192\\
3.8013	0.38883\\
3.8045	0.38579\\
3.8078	0.38281\\
3.811	0.37989\\
3.8143	0.37702\\
3.8175	0.3742\\
3.8207	0.37143\\
3.824	0.36871\\
3.8272	0.36605\\
3.8304	0.36342\\
3.8337	0.36085\\
3.8369	0.35832\\
3.8402	0.35583\\
3.8434	0.35339\\
3.8466	0.35099\\
3.8499	0.34863\\
3.8531	0.34631\\
3.8564	0.34403\\
3.8596	0.34179\\
3.8628	0.33958\\
3.8661	0.33741\\
3.8693	0.33527\\
3.8725	0.33317\\
3.8758	0.3311\\
3.879	0.32907\\
3.8823	0.32706\\
3.8855	0.32508\\
3.8887	0.32314\\
3.892	0.32122\\
3.8952	0.31933\\
3.8984	0.31746\\
3.9017	0.31562\\
3.9049	0.31381\\
3.9082	0.31202\\
3.9114	0.31025\\
3.9146	0.30851\\
3.9179	0.30678\\
3.9211	0.30508\\
3.9243	0.3034\\
3.9276	0.30174\\
3.9308	0.3001\\
3.9341	0.29847\\
3.9373	0.29686\\
3.9405	0.29527\\
3.9438	0.2937\\
3.947	0.29214\\
3.9502	0.2906\\
3.9535	0.28907\\
3.9567	0.28755\\
3.96	0.28605\\
3.9632	0.28456\\
3.9664	0.28308\\
3.9697	0.28161\\
3.9729	0.28016\\
3.9762	0.27871\\
3.9794	0.27727\\
3.9826	0.27585\\
3.9859	0.27443\\
3.9891	0.27302\\
3.9923	0.27162\\
3.9956	0.27022\\
3.9988	0.26884\\
4.0021	0.26746\\
4.0053	0.26608\\
4.0085	0.26471\\
4.0118	0.26335\\
4.015	0.26199\\
4.0182	0.26064\\
4.0215	0.25929\\
4.0247	0.25795\\
4.028	0.25661\\
4.0312	0.25527\\
4.0344	0.25394\\
4.0377	0.25261\\
4.0409	0.25128\\
4.0441	0.24995\\
4.0474	0.24863\\
4.0506	0.24731\\
4.0539	0.24599\\
4.0571	0.24467\\
4.0603	0.24335\\
4.0636	0.24203\\
4.0668	0.24072\\
4.07	0.2394\\
4.0733	0.23809\\
4.0765	0.23678\\
4.0798	0.23546\\
4.083	0.23415\\
4.0862	0.23284\\
4.0895	0.23152\\
4.0927	0.23021\\
4.0959	0.2289\\
4.0992	0.22758\\
4.1024	0.22627\\
4.1057	0.22495\\
4.1089	0.22364\\
4.1121	0.22232\\
4.1154	0.221\\
4.1186	0.21968\\
4.1219	0.21837\\
4.1251	0.21705\\
4.1283	0.21573\\
4.1316	0.21441\\
4.1348	0.21308\\
4.138	0.21176\\
4.1413	0.21044\\
4.1445	0.20912\\
4.1478	0.20779\\
4.151	0.20647\\
4.1542	0.20514\\
4.1575	0.20381\\
4.1607	0.20249\\
4.1639	0.20116\\
4.1672	0.19983\\
4.1704	0.1985\\
4.1737	0.19717\\
4.1769	0.19584\\
4.1801	0.19451\\
4.1834	0.19318\\
4.1866	0.19185\\
4.1898	0.19052\\
4.1931	0.18919\\
4.1963	0.18786\\
4.1996	0.18653\\
4.2028	0.18519\\
4.206	0.18386\\
4.2093	0.18253\\
4.2125	0.1812\\
4.2157	0.17987\\
4.219	0.17854\\
4.2222	0.1772\\
4.2255	0.17587\\
4.2287	0.17454\\
4.2319	0.17321\\
4.2352	0.17188\\
4.2384	0.17055\\
4.2417	0.16922\\
4.2449	0.16789\\
4.2481	0.16656\\
4.2514	0.16523\\
4.2546	0.1639\\
4.2578	0.16257\\
4.2611	0.16124\\
4.2643	0.15991\\
4.2676	0.15858\\
4.2708	0.15725\\
4.274	0.15593\\
4.2773	0.1546\\
4.2805	0.15328\\
4.2837	0.15195\\
4.287	0.15063\\
4.2902	0.1493\\
4.2935	0.14798\\
4.2967	0.14666\\
4.2999	0.14534\\
4.3032	0.14402\\
4.3064	0.14271\\
4.3096	0.14139\\
4.3129	0.14008\\
4.3161	0.13877\\
4.3194	0.13746\\
4.3226	0.13615\\
4.3258	0.13484\\
4.3291	0.13354\\
4.3323	0.13224\\
4.3355	0.13094\\
4.3388	0.12964\\
4.342	0.12835\\
4.3453	0.12706\\
4.3485	0.12577\\
4.3517	0.12448\\
4.355	0.1232\\
4.3582	0.12192\\
4.3615	0.12065\\
4.3647	0.11938\\
4.3679	0.11811\\
4.3712	0.11685\\
4.3744	0.11559\\
4.3776	0.11433\\
4.3809	0.11308\\
4.3841	0.11183\\
4.3874	0.11059\\
4.3906	0.10935\\
4.3938	0.10812\\
4.3971	0.10689\\
4.4003	0.10567\\
4.4035	0.10445\\
4.4068	0.10323\\
4.41	0.10203\\
4.4133	0.10083\\
4.4165	0.099629\\
4.4197	0.098439\\
4.423	0.097254\\
4.4262	0.096076\\
4.4294	0.094904\\
4.4327	0.093738\\
4.4359	0.092578\\
4.4392	0.091425\\
4.4424	0.090278\\
4.4456	0.089139\\
4.4489	0.088006\\
4.4521	0.08688\\
4.4553	0.085761\\
4.4586	0.084649\\
4.4618	0.083544\\
4.4651	0.082447\\
4.4683	0.081357\\
4.4715	0.080275\\
4.4748	0.0792\\
4.478	0.078133\\
4.4813	0.077074\\
4.4845	0.076022\\
4.4877	0.074979\\
4.491	0.073943\\
4.4942	0.072916\\
4.4974	0.071896\\
4.5007	0.070885\\
4.5039	0.069882\\
4.5072	0.068888\\
4.5104	0.067902\\
4.5136	0.066924\\
4.5169	0.065955\\
4.5201	0.064995\\
4.5233	0.064043\\
4.5266	0.063099\\
4.5298	0.062165\\
4.5331	0.061239\\
4.5363	0.060322\\
4.5395	0.059414\\
4.5428	0.058514\\
4.546	0.057624\\
4.5492	0.056742\\
4.5525	0.05587\\
4.5557	0.055006\\
4.559	0.054152\\
4.5622	0.053306\\
4.5654	0.052469\\
4.5687	0.051642\\
4.5719	0.050823\\
4.5751	0.050014\\
4.5784	0.049213\\
4.5816	0.048422\\
4.5849	0.047639\\
4.5881	0.046866\\
4.5913	0.046102\\
4.5946	0.045347\\
4.5978	0.044601\\
4.6011	0.043863\\
4.6043	0.043135\\
4.6075	0.042416\\
4.6108	0.041706\\
4.614	0.041005\\
4.6172	0.040312\\
4.6205	0.039629\\
4.6237	0.038954\\
4.627	0.038288\\
4.6302	0.037631\\
4.6334	0.036983\\
4.6367	0.036343\\
4.6399	0.035712\\
4.6431	0.03509\\
4.6464	0.034476\\
4.6496	0.033871\\
4.6529	0.033274\\
4.6561	0.032685\\
4.6593	0.032105\\
4.6626	0.031533\\
4.6658	0.030969\\
4.669	0.030414\\
4.6723	0.029866\\
4.6755	0.029327\\
4.6788	0.028795\\
4.682	0.028272\\
4.6852	0.027756\\
4.6885	0.027248\\
4.6917	0.026747\\
4.6949	0.026254\\
4.6982	0.025769\\
4.7014	0.025291\\
4.7047	0.02482\\
4.7079	0.024357\\
4.7111	0.023901\\
4.7144	0.023452\\
4.7176	0.02301\\
4.7209	0.022575\\
4.7241	0.022147\\
4.7273	0.021726\\
4.7306	0.021311\\
4.7338	0.020903\\
4.737	0.020502\\
4.7403	0.020107\\
4.7435	0.019719\\
4.7468	0.019337\\
4.75	0.018961\\
4.7532	0.018591\\
4.7565	0.018228\\
4.7597	0.01787\\
4.7629	0.017519\\
4.7662	0.017173\\
4.7694	0.016834\\
4.7727	0.0165\\
4.7759	0.016171\\
4.7791	0.015849\\
4.7824	0.015531\\
4.7856	0.01522\\
4.7888	0.014913\\
4.7921	0.014612\\
4.7953	0.014316\\
4.7986	0.014026\\
4.8018	0.01374\\
4.805	0.01346\\
4.8083	0.013184\\
4.8115	0.012914\\
4.8147	0.012648\\
4.818	0.012387\\
4.8212	0.01213\\
4.8245	0.011879\\
4.8277	0.011631\\
4.8309	0.011389\\
4.8342	0.011151\\
4.8374	0.010917\\
4.8406	0.010687\\
4.8439	0.010462\\
4.8471	0.010241\\
4.8504	0.010024\\
4.8536	0.0098107\\
4.8568	0.0096017\\
4.8601	0.0093966\\
4.8633	0.0091955\\
4.8666	0.0089981\\
4.8698	0.0088045\\
4.873	0.0086146\\
4.8763	0.0084283\\
4.8795	0.0082457\\
4.8827	0.0080665\\
4.886	0.0078908\\
4.8892	0.0077185\\
4.8925	0.0075496\\
4.8957	0.007384\\
4.8989	0.0072217\\
4.9022	0.0070625\\
4.9054	0.0069065\\
4.9086	0.0067536\\
4.9119	0.0066037\\
4.9151	0.0064568\\
4.9184	0.0063129\\
4.9216	0.0061718\\
4.9248	0.0060336\\
4.9281	0.0058982\\
4.9313	0.0057655\\
4.9345	0.0056356\\
4.9378	0.0055083\\
4.941	0.0053836\\
4.9443	0.0052614\\
4.9475	0.0051418\\
4.9507	0.0050246\\
4.954	0.0049099\\
4.9572	0.0047976\\
4.9604	0.0046876\\
4.9637	0.0045799\\
4.9669	0.0044744\\
4.9702	0.0043712\\
4.9734	0.0042701\\
4.9766	0.0041712\\
4.9799	0.0040744\\
4.9831	0.0039797\\
4.9864	0.0038869\\
4.9896	0.0037962\\
4.9928	0.0037074\\
4.9961	0.0036205\\
4.9993	0.0035354\\
5.0025	0.0034523\\
5.0058	0.0033709\\
5.009	0.0032913\\
5.0123	0.0032134\\
5.0155	0.0031372\\
5.0187	0.0030627\\
5.022	0.0029898\\
5.0252	0.0029185\\
5.0284	0.0028488\\
5.0317	0.0027807\\
5.0349	0.002714\\
5.0382	0.0026488\\
5.0414	0.0025851\\
5.0446	0.0025228\\
5.0479	0.0024619\\
5.0511	0.0024024\\
5.0543	0.0023442\\
5.0576	0.0022873\\
5.0608	0.0022317\\
5.0641	0.0021774\\
5.0673	0.0021243\\
5.0705	0.0020724\\
5.0738	0.0020217\\
5.077	0.0019721\\
5.0802	0.0019237\\
5.0835	0.0018764\\
5.0867	0.0018301\\
5.09	0.001785\\
5.0932	0.0017409\\
5.0964	0.0016978\\
5.0997	0.0016557\\
5.1029	0.0016145\\
5.1062	0.0015744\\
5.1094	0.0015352\\
5.1126	0.0014969\\
5.1159	0.0014594\\
5.1191	0.0014229\\
5.1223	0.0013872\\
5.1256	0.0013524\\
5.1288	0.0013184\\
5.1321	0.0012852\\
5.1353	0.0012528\\
5.1385	0.0012211\\
5.1418	0.0011902\\
5.145	0.0011601\\
5.1482	0.0011306\\
5.1515	0.0011019\\
5.1547	0.0010739\\
5.158	0.0010465\\
5.1612	0.0010198\\
5.1644	0.00099371\\
5.1677	0.00096828\\
5.1709	0.00094345\\
5.1741	0.00091923\\
5.1774	0.0008956\\
5.1806	0.00087254\\
5.1839	0.00085004\\
5.1871	0.0008281\\
5.1903	0.00080668\\
5.1936	0.0007858\\
5.1968	0.00076542\\
5.2	0.00074555\\
5.2033	0.00072617\\
5.2065	0.00070726\\
5.2098	0.00068882\\
5.213	0.00067084\\
5.2162	0.0006533\\
5.2195	0.0006362\\
5.2227	0.00061952\\
5.226	0.00060326\\
5.2292	0.0005874\\
5.2324	0.00057194\\
5.2357	0.00055687\\
5.2389	0.00054218\\
5.2421	0.00052785\\
5.2454	0.00051388\\
5.2486	0.00050027\\
5.2519	0.000487\\
5.2551	0.00047406\\
5.2583	0.00046145\\
5.2616	0.00044916\\
5.2648	0.00043718\\
5.268	0.00042551\\
5.2713	0.00041413\\
5.2745	0.00040305\\
5.2778	0.00039224\\
5.281	0.00038171\\
5.2842	0.00037146\\
5.2875	0.00036146\\
5.2907	0.00035172\\
5.2939	0.00034224\\
5.2972	0.00033299\\
5.3004	0.00032399\\
5.3037	0.00031521\\
5.3069	0.00030667\\
5.3101	0.00029834\\
5.3101	0.00021101\\
5.3069	0.00021732\\
5.3037	0.00022382\\
5.3004	0.0002305\\
5.2972	0.00023737\\
5.2939	0.00024444\\
5.2907	0.00025171\\
5.2875	0.00025918\\
5.2842	0.00026687\\
5.281	0.00027478\\
5.2778	0.00028291\\
5.2745	0.00029128\\
5.2713	0.00029988\\
5.268	0.00030872\\
5.2648	0.00031781\\
5.2616	0.00032716\\
5.2583	0.00033677\\
5.2551	0.00034665\\
5.2519	0.00035681\\
5.2486	0.00036725\\
5.2454	0.00037798\\
5.2421	0.00038902\\
5.2389	0.00040036\\
5.2357	0.00041201\\
5.2324	0.00042399\\
5.2292	0.00043631\\
5.226	0.00044896\\
5.2227	0.00046196\\
5.2195	0.00047532\\
5.2162	0.00048905\\
5.213	0.00050316\\
5.2098	0.00051766\\
5.2065	0.00053255\\
5.2033	0.00054785\\
5.2	0.00056357\\
5.1968	0.00057971\\
5.1936	0.0005963\\
5.1903	0.00061334\\
5.1871	0.00063084\\
5.1839	0.00064881\\
5.1806	0.00066727\\
5.1774	0.00068623\\
5.1741	0.0007057\\
5.1709	0.00072569\\
5.1677	0.00074622\\
5.1644	0.0007673\\
5.1612	0.00078895\\
5.158	0.00081117\\
5.1547	0.00083399\\
5.1515	0.00085741\\
5.1482	0.00088145\\
5.145	0.00090613\\
5.1418	0.00093146\\
5.1385	0.00095747\\
5.1353	0.00098415\\
5.1321	0.0010115\\
5.1288	0.0010396\\
5.1256	0.0010685\\
5.1223	0.0010981\\
5.1191	0.0011285\\
5.1159	0.0011596\\
5.1126	0.0011916\\
5.1094	0.0012244\\
5.1062	0.001258\\
5.1029	0.0012926\\
5.0997	0.001328\\
5.0964	0.0013643\\
5.0932	0.0014015\\
5.09	0.0014397\\
5.0867	0.0014789\\
5.0835	0.0015191\\
5.0802	0.0015603\\
5.077	0.0016025\\
5.0738	0.0016458\\
5.0705	0.0016902\\
5.0673	0.0017357\\
5.0641	0.0017824\\
5.0608	0.0018302\\
5.0576	0.0018793\\
5.0543	0.0019295\\
5.0511	0.001981\\
5.0479	0.0020338\\
5.0446	0.0020879\\
5.0414	0.0021433\\
5.0382	0.0022001\\
5.0349	0.0022583\\
5.0317	0.0023179\\
5.0284	0.002379\\
5.0252	0.0024416\\
5.022	0.0025057\\
5.0187	0.0025713\\
5.0155	0.0026386\\
5.0123	0.0027074\\
5.009	0.0027779\\
5.0058	0.0028501\\
5.0025	0.0029241\\
4.9993	0.0029997\\
4.9961	0.0030772\\
4.9928	0.0031566\\
4.9896	0.0032378\\
4.9864	0.0033209\\
4.9831	0.003406\\
4.9799	0.0034931\\
4.9766	0.0035822\\
4.9734	0.0036734\\
4.9702	0.0037667\\
4.9669	0.0038622\\
4.9637	0.0039599\\
4.9604	0.0040598\\
4.9572	0.004162\\
4.954	0.0042666\\
4.9507	0.0043735\\
4.9475	0.0044829\\
4.9443	0.0045947\\
4.941	0.0047091\\
4.9378	0.0048261\\
4.9345	0.0049456\\
4.9313	0.0050679\\
4.9281	0.0051929\\
4.9248	0.0053206\\
4.9216	0.0054512\\
4.9184	0.0055846\\
4.9151	0.005721\\
4.9119	0.0058603\\
4.9086	0.0060027\\
4.9054	0.0061482\\
4.9022	0.0062968\\
4.8989	0.0064486\\
4.8957	0.0066037\\
4.8925	0.0067621\\
4.8892	0.0069238\\
4.886	0.007089\\
4.8827	0.0072577\\
4.8795	0.0074299\\
4.8763	0.0076057\\
4.873	0.0077852\\
4.8698	0.0079685\\
4.8666	0.0081555\\
4.8633	0.0083463\\
4.8601	0.0085411\\
4.8568	0.0087398\\
4.8536	0.0089426\\
4.8504	0.0091495\\
4.8471	0.0093605\\
4.8439	0.0095758\\
4.8406	0.0097954\\
4.8374	0.010019\\
4.8342	0.010248\\
4.8309	0.01048\\
4.8277	0.010718\\
4.8245	0.01096\\
4.8212	0.011207\\
4.818	0.011458\\
4.8147	0.011714\\
4.8115	0.011975\\
4.8083	0.012241\\
4.805	0.012512\\
4.8018	0.012789\\
4.7986	0.01307\\
4.7953	0.013357\\
4.7921	0.013648\\
4.7888	0.013945\\
4.7856	0.014248\\
4.7824	0.014556\\
4.7791	0.01487\\
4.7759	0.015189\\
4.7727	0.015514\\
4.7694	0.015844\\
4.7662	0.016181\\
4.7629	0.016523\\
4.7597	0.016871\\
4.7565	0.017226\\
4.7532	0.017586\\
4.75	0.017952\\
4.7468	0.018325\\
4.7435	0.018704\\
4.7403	0.019089\\
4.737	0.01948\\
4.7338	0.019878\\
4.7306	0.020282\\
4.7273	0.020693\\
4.7241	0.021111\\
4.7209	0.021535\\
4.7176	0.021966\\
4.7144	0.022403\\
4.7111	0.022847\\
4.7079	0.023299\\
4.7047	0.023757\\
4.7014	0.024222\\
4.6982	0.024694\\
4.6949	0.025173\\
4.6917	0.025659\\
4.6885	0.026152\\
4.6852	0.026652\\
4.682	0.02716\\
4.6788	0.027675\\
4.6755	0.028197\\
4.6723	0.028726\\
4.669	0.029263\\
4.6658	0.029808\\
4.6626	0.030359\\
4.6593	0.030919\\
4.6561	0.031485\\
4.6529	0.03206\\
4.6496	0.032642\\
4.6464	0.033231\\
4.6431	0.033829\\
4.6399	0.034434\\
4.6367	0.035047\\
4.6334	0.035667\\
4.6302	0.036296\\
4.627	0.036932\\
4.6237	0.037576\\
4.6205	0.038229\\
4.6172	0.038889\\
4.614	0.039557\\
4.6108	0.040233\\
4.6075	0.040917\\
4.6043	0.04161\\
4.6011	0.04231\\
4.5978	0.043019\\
4.5946	0.043736\\
4.5913	0.044461\\
4.5881	0.045194\\
4.5849	0.045936\\
4.5816	0.046686\\
4.5784	0.047444\\
4.5751	0.04821\\
4.5719	0.048985\\
4.5687	0.049768\\
4.5654	0.05056\\
4.5622	0.05136\\
4.559	0.052168\\
4.5557	0.052985\\
4.5525	0.05381\\
4.5492	0.054644\\
4.546	0.055486\\
4.5428	0.056336\\
4.5395	0.057195\\
4.5363	0.058062\\
4.5331	0.058938\\
4.5298	0.059822\\
4.5266	0.060715\\
4.5233	0.061616\\
4.5201	0.062525\\
4.5169	0.063443\\
4.5136	0.064369\\
4.5104	0.065303\\
4.5072	0.066246\\
4.5039	0.067197\\
4.5007	0.068156\\
4.4974	0.069124\\
4.4942	0.070099\\
4.491	0.071083\\
4.4877	0.072075\\
4.4845	0.073076\\
4.4813	0.074084\\
4.478	0.0751\\
4.4748	0.076124\\
4.4715	0.077156\\
4.4683	0.078196\\
4.4651	0.079244\\
4.4618	0.0803\\
4.4586	0.081363\\
4.4553	0.082434\\
4.4521	0.083513\\
4.4489	0.084599\\
4.4456	0.085692\\
4.4424	0.086793\\
4.4392	0.087901\\
4.4359	0.089017\\
4.4327	0.090139\\
4.4294	0.091269\\
4.4262	0.092406\\
4.423	0.093549\\
4.4197	0.0947\\
4.4165	0.095857\\
4.4133	0.097021\\
4.41	0.098191\\
4.4068	0.099368\\
4.4035	0.10055\\
4.4003	0.10174\\
4.3971	0.10294\\
4.3938	0.10414\\
4.3906	0.10535\\
4.3874	0.10656\\
4.3841	0.10778\\
4.3809	0.109\\
4.3776	0.11023\\
4.3744	0.11147\\
4.3712	0.11271\\
4.3679	0.11395\\
4.3647	0.11521\\
4.3615	0.11646\\
4.3582	0.11772\\
4.355	0.11899\\
4.3517	0.12026\\
4.3485	0.12153\\
4.3453	0.12281\\
4.342	0.12409\\
4.3388	0.12537\\
4.3355	0.12666\\
4.3323	0.12796\\
4.3291	0.12925\\
4.3258	0.13055\\
4.3226	0.13186\\
4.3194	0.13316\\
4.3161	0.13447\\
4.3129	0.13578\\
4.3096	0.13709\\
4.3064	0.13841\\
4.3032	0.13973\\
4.2999	0.14105\\
4.2967	0.14237\\
4.2935	0.14369\\
4.2902	0.14502\\
4.287	0.14634\\
4.2837	0.14767\\
4.2805	0.14899\\
4.2773	0.15032\\
4.274	0.15165\\
4.2708	0.15298\\
4.2676	0.15431\\
4.2643	0.15564\\
4.2611	0.15696\\
4.2578	0.15829\\
4.2546	0.15962\\
4.2514	0.16095\\
4.2481	0.16227\\
4.2449	0.16359\\
4.2417	0.16492\\
4.2384	0.16624\\
4.2352	0.16755\\
4.2319	0.16887\\
4.2287	0.17018\\
4.2255	0.1715\\
4.2222	0.17281\\
4.219	0.17411\\
4.2157	0.17542\\
4.2125	0.17672\\
4.2093	0.17801\\
4.206	0.17931\\
4.2028	0.1806\\
4.1996	0.18188\\
4.1963	0.18317\\
4.1931	0.18445\\
4.1898	0.18572\\
4.1866	0.187\\
4.1834	0.18827\\
4.1801	0.18953\\
4.1769	0.19079\\
4.1737	0.19205\\
4.1704	0.1933\\
4.1672	0.19455\\
4.1639	0.19579\\
4.1607	0.19703\\
4.1575	0.19827\\
4.1542	0.1995\\
4.151	0.20073\\
4.1478	0.20196\\
4.1445	0.20318\\
4.1413	0.2044\\
4.138	0.20561\\
4.1348	0.20683\\
4.1316	0.20803\\
4.1283	0.20924\\
4.1251	0.21044\\
4.1219	0.21164\\
4.1186	0.21284\\
4.1154	0.21403\\
4.1121	0.21522\\
4.1089	0.21641\\
4.1057	0.2176\\
4.1024	0.21879\\
4.0992	0.21997\\
4.0959	0.22115\\
4.0927	0.22234\\
4.0895	0.22352\\
4.0862	0.2247\\
4.083	0.22588\\
4.0798	0.22706\\
4.0765	0.22823\\
4.0733	0.22941\\
4.07	0.23059\\
4.0668	0.23177\\
4.0636	0.23296\\
4.0603	0.23414\\
4.0571	0.23533\\
4.0539	0.23651\\
4.0506	0.2377\\
4.0474	0.2389\\
4.0441	0.24009\\
4.0409	0.24129\\
4.0377	0.24249\\
4.0344	0.2437\\
4.0312	0.24491\\
4.028	0.24613\\
4.0247	0.24735\\
4.0215	0.24858\\
4.0182	0.24982\\
4.015	0.25106\\
4.0118	0.25231\\
4.0085	0.25357\\
4.0053	0.25483\\
4.0021	0.25611\\
3.9988	0.25739\\
3.9956	0.25869\\
3.9923	0.25999\\
3.9891	0.2613\\
3.9859	0.26263\\
3.9826	0.26397\\
3.9794	0.26532\\
3.9762	0.26668\\
3.9729	0.26806\\
3.9697	0.26945\\
3.9664	0.27086\\
3.9632	0.27228\\
3.96	0.27372\\
3.9567	0.27518\\
3.9535	0.27665\\
3.9502	0.27814\\
3.947	0.27965\\
3.9438	0.28118\\
3.9405	0.28273\\
3.9373	0.2843\\
3.9341	0.28589\\
3.9308	0.2875\\
3.9276	0.28914\\
3.9243	0.2908\\
3.9211	0.29249\\
3.9179	0.2942\\
3.9146	0.29594\\
3.9114	0.2977\\
3.9082	0.29949\\
3.9049	0.30131\\
3.9017	0.30316\\
3.8984	0.30503\\
3.8952	0.30694\\
3.892	0.30888\\
3.8887	0.31085\\
3.8855	0.31286\\
3.8823	0.3149\\
3.879	0.31697\\
3.8758	0.31908\\
3.8725	0.32122\\
3.8693	0.3234\\
3.8661	0.32562\\
3.8628	0.32788\\
3.8596	0.33018\\
3.8564	0.33251\\
3.8531	0.33489\\
3.8499	0.33731\\
3.8466	0.33977\\
3.8434	0.34227\\
3.8402	0.34482\\
3.8369	0.34741\\
3.8337	0.35005\\
3.8304	0.35274\\
3.8272	0.35547\\
3.824	0.35824\\
3.8207	0.36107\\
3.8175	0.36394\\
3.8143	0.36686\\
3.811	0.36983\\
3.8078	0.37286\\
3.8045	0.37593\\
3.8013	0.37905\\
3.7981	0.38223\\
3.7948	0.38546\\
3.7916	0.38874\\
3.7884	0.39207\\
3.7851	0.39546\\
3.7819	0.3989\\
3.7786	0.4024\\
3.7754	0.40595\\
3.7722	0.40956\\
3.7689	0.41322\\
3.7657	0.41694\\
3.7625	0.42072\\
3.7592	0.42455\\
3.756	0.42845\\
3.7527	0.4324\\
3.7495	0.43641\\
3.7463	0.44048\\
3.743	0.44461\\
3.7398	0.4488\\
3.7366	0.45306\\
3.7333	0.45738\\
3.7301	0.46176\\
3.7268	0.4662\\
3.7236	0.47071\\
3.7204	0.47529\\
3.7171	0.47993\\
3.7139	0.48464\\
3.7106	0.48941\\
3.7074	0.49426\\
3.7042	0.49917\\
3.7009	0.50416\\
3.6977	0.50921\\
3.6945	0.51433\\
3.6912	0.51953\\
3.688	0.52479\\
3.6847	0.53013\\
3.6815	0.53554\\
3.6783	0.54102\\
3.675	0.54658\\
3.6718	0.5522\\
3.6686	0.5579\\
3.6653	0.56368\\
3.6621	0.56952\\
3.6588	0.57544\\
3.6556	0.58143\\
3.6524	0.5875\\
3.6491	0.59363\\
3.6459	0.59984\\
3.6427	0.60612\\
3.6394	0.61247\\
3.6362	0.6189\\
3.6329	0.62539\\
3.6297	0.63195\\
3.6265	0.63858\\
3.6232	0.64528\\
3.62	0.65205\\
3.6168	0.65889\\
3.6135	0.66579\\
3.6103	0.67276\\
3.607	0.67979\\
3.6038	0.68688\\
3.6006	0.69403\\
3.5973	0.70125\\
3.5941	0.70852\\
3.5908	0.71586\\
3.5876	0.72324\\
3.5844	0.73069\\
3.5811	0.73818\\
3.5779	0.74573\\
3.5747	0.75333\\
3.5714	0.76097\\
3.5682	0.76866\\
3.5649	0.7764\\
3.5617	0.78418\\
3.5585	0.79199\\
3.5552	0.79985\\
3.552	0.80774\\
3.5488	0.81567\\
3.5455	0.82362\\
3.5423	0.83161\\
3.539	0.83962\\
3.5358	0.84765\\
3.5326	0.8557\\
3.5293	0.86378\\
3.5261	0.87187\\
3.5229	0.87997\\
3.5196	0.88808\\
3.5164	0.89619\\
3.5131	0.90431\\
3.5099	0.91244\\
3.5067	0.92055\\
3.5034	0.92867\\
3.5002	0.93677\\
3.497	0.94486\\
3.4937	0.95294\\
3.4905	0.96099\\
3.4872	0.96903\\
3.484	0.97704\\
3.4808	0.98501\\
3.4775	0.99296\\
3.4743	1.0009\\
3.471	1.0087\\
3.4678	1.0166\\
3.4646	1.0243\\
3.4613	1.0321\\
3.4581	1.0397\\
3.4549	1.0473\\
3.4516	1.0549\\
3.4484	1.0623\\
3.4451	1.0697\\
3.4419	1.0771\\
3.4387	1.0843\\
3.4354	1.0915\\
3.4322	1.0986\\
3.429	1.1056\\
3.4257	1.1124\\
3.4225	1.1192\\
3.4192	1.1259\\
3.416	1.1325\\
3.4128	1.139\\
3.4095	1.1453\\
3.4063	1.1516\\
3.4031	1.1577\\
3.3998	1.1637\\
3.3966	1.1695\\
3.3933	1.1752\\
3.3901	1.1808\\
3.3869	1.1863\\
3.3836	1.1916\\
3.3804	1.1967\\
3.3772	1.2017\\
3.3739	1.2066\\
3.3707	1.2113\\
3.3674	1.2158\\
3.3642	1.2202\\
3.361	1.2244\\
3.3577	1.2284\\
3.3545	1.2323\\
3.3513	1.236\\
3.348	1.2395\\
3.3448	1.2429\\
3.3415	1.246\\
3.3383	1.249\\
3.3351	1.2518\\
3.3318	1.2545\\
3.3286	1.2569\\
3.3253	1.2592\\
3.3221	1.2612\\
3.3189	1.2631\\
3.3156	1.2648\\
3.3124	1.2663\\
3.3092	1.2676\\
3.3059	1.2687\\
3.3027	1.2696\\
3.2994	1.2703\\
3.2962	1.2708\\
3.293	1.2711\\
3.2897	1.2712\\
3.2865	1.2712\\
3.2833	1.2709\\
3.28	1.2704\\
3.2768	1.2697\\
3.2735	1.2689\\
3.2703	1.2678\\
3.2671	1.2665\\
3.2638	1.2651\\
3.2606	1.2634\\
3.2574	1.2615\\
3.2541	1.2595\\
3.2509	1.2572\\
3.2476	1.2548\\
3.2444	1.2521\\
3.2412	1.2493\\
3.2379	1.2462\\
3.2347	1.243\\
3.2315	1.2396\\
3.2282	1.236\\
3.225	1.2322\\
3.2217	1.2282\\
3.2185	1.2241\\
3.2153	1.2198\\
3.212	1.2152\\
3.2088	1.2105\\
3.2055	1.2057\\
3.2023	1.2006\\
3.1991	1.1954\\
3.1958	1.19\\
3.1926	1.1844\\
3.1894	1.1787\\
3.1861	1.1728\\
3.1829	1.1668\\
3.1796	1.1606\\
3.1764	1.1542\\
3.1732	1.1477\\
3.1699	1.141\\
3.1667	1.1342\\
3.1635	1.1272\\
3.1602	1.1201\\
3.157	1.1129\\
3.1537	1.1055\\
3.1505	1.098\\
3.1473	1.0904\\
3.144	1.0826\\
3.1408	1.0748\\
3.1376	1.0668\\
3.1343	1.0586\\
3.1311	1.0504\\
3.1278	1.0421\\
3.1246	1.0336\\
3.1214	1.0251\\
3.1181	1.0164\\
3.1149	1.0077\\
3.1117	0.99888\\
3.1084	0.98996\\
3.1052	0.98096\\
3.1019	0.97188\\
3.0987	0.96272\\
3.0955	0.95349\\
3.0922	0.94419\\
3.089	0.93482\\
3.0857	0.92538\\
3.0825	0.91589\\
3.0793	0.90634\\
3.076	0.89674\\
3.0728	0.88709\\
3.0696	0.8774\\
3.0663	0.86766\\
3.0631	0.85788\\
3.0598	0.84807\\
3.0566	0.83823\\
3.0534	0.82836\\
3.0501	0.81847\\
3.0469	0.80855\\
3.0437	0.79862\\
3.0404	0.78867\\
3.0372	0.77871\\
3.0339	0.76875\\
3.0307	0.75878\\
3.0275	0.74881\\
3.0242	0.73884\\
3.021	0.72888\\
3.0178	0.71893\\
3.0145	0.70899\\
3.0113	0.69907\\
3.008	0.68916\\
3.0048	0.67928\\
3.0016	0.66943\\
2.9983	0.65961\\
2.9951	0.64982\\
2.9919	0.64006\\
2.9886	0.63034\\
2.9854	0.62066\\
2.9821	0.61103\\
2.9789	0.60145\\
2.9757	0.59191\\
2.9724	0.58243\\
2.9692	0.57301\\
2.9659	0.56364\\
2.9627	0.55433\\
2.9595	0.54508\\
2.9562	0.5359\\
2.953	0.52679\\
2.9498	0.51775\\
2.9465	0.50878\\
2.9433	0.49988\\
2.94	0.49106\\
2.9368	0.48231\\
2.9336	0.47365\\
2.9303	0.46506\\
2.9271	0.45656\\
2.9239	0.44814\\
2.9206	0.4398\\
2.9174	0.43156\\
2.9141	0.42339\\
2.9109	0.41532\\
2.9077	0.40734\\
2.9044	0.39945\\
2.9012	0.39165\\
2.898	0.38395\\
2.8947	0.37633\\
2.8915	0.36881\\
2.8882	0.36139\\
2.885	0.35406\\
2.8818	0.34683\\
2.8785	0.33969\\
2.8753	0.33265\\
2.8721	0.32571\\
2.8688	0.31886\\
2.8656	0.31212\\
2.8623	0.30546\\
2.8591	0.29891\\
2.8559	0.29246\\
2.8526	0.2861\\
2.8494	0.27984\\
2.8461	0.27368\\
2.8429	0.26761\\
2.8397	0.26164\\
2.8364	0.25577\\
2.8332	0.24999\\
2.83	0.24431\\
2.8267	0.23873\\
2.8235	0.23324\\
2.8202	0.22784\\
2.817	0.22254\\
2.8138	0.21733\\
2.8105	0.21222\\
2.8073	0.2072\\
2.8041	0.20227\\
2.8008	0.19743\\
2.7976	0.19268\\
2.7943	0.18801\\
2.7911	0.18344\\
2.7879	0.17896\\
2.7846	0.17456\\
2.7814	0.17025\\
2.7782	0.16602\\
2.7749	0.16188\\
2.7717	0.15782\\
2.7684	0.15384\\
2.7652	0.14995\\
2.762	0.14613\\
2.7587	0.1424\\
2.7555	0.13874\\
2.7523	0.13516\\
2.749	0.13166\\
2.7458	0.12823\\
2.7425	0.12487\\
2.7393	0.12159\\
2.7361	0.11839\\
2.7328	0.11525\\
2.7296	0.11218\\
2.7263	0.10918\\
2.7231	0.10625\\
2.7199	0.10339\\
2.7166	0.10059\\
2.7134	0.097863\\
2.7102	0.095195\\
2.7069	0.092589\\
2.7037	0.090045\\
2.7004	0.087562\\
2.6972	0.085138\\
2.694	0.082772\\
2.6907	0.080464\\
2.6875	0.078212\\
2.6843	0.076016\\
2.681	0.073874\\
2.6778	0.071785\\
2.6745	0.069748\\
2.6713	0.067762\\
2.6681	0.065827\\
2.6648	0.06394\\
2.6616	0.062102\\
2.6584	0.060311\\
2.6551	0.058567\\
2.6519	0.056867\\
2.6486	0.055212\\
2.6454	0.053601\\
2.6422	0.052031\\
2.6389	0.050504\\
2.6357	0.049017\\
2.6325	0.04757\\
2.6292	0.046161\\
2.626	0.044791\\
2.6227	0.043458\\
2.6195	0.042161\\
2.6163	0.0409\\
2.613	0.039673\\
2.6098	0.03848\\
2.6066	0.03732\\
2.6033	0.036193\\
2.6001	0.035097\\
2.5968	0.034031\\
2.5936	0.032996\\
2.5904	0.03199\\
2.5871	0.031013\\
2.5839	0.030063\\
2.5806	0.02914\\
2.5774	0.028244\\
2.5742	0.027374\\
2.5709	0.026529\\
2.5677	0.025709\\
2.5645	0.024912\\
2.5612	0.024138\\
2.558	0.023388\\
2.5547	0.022659\\
2.5515	0.021952\\
2.5483	0.021265\\
2.545	0.0206\\
2.5418	0.019953\\
2.5386	0.019327\\
2.5353	0.018719\\
2.5321	0.018129\\
2.5288	0.017557\\
2.5256	0.017002\\
2.5224	0.016464\\
2.5191	0.015943\\
2.5159	0.015437\\
2.5127	0.014947\\
2.5094	0.014471\\
2.5062	0.014011\\
2.5029	0.013564\\
2.4997	0.013131\\
2.4965	0.012712\\
2.4932	0.012306\\
2.49	0.011912\\
2.4868	0.011531\\
2.4835	0.011161\\
2.4803	0.010803\\
2.477	0.010456\\
2.4738	0.01012\\
2.4706	0.009795\\
2.4673	0.00948\\
2.4641	0.0091749\\
2.4608	0.0088795\\
2.4576	0.0085934\\
2.4544	0.0083165\\
2.4511	0.0080483\\
2.4479	0.0077888\\
2.4447	0.0075375\\
2.4414	0.0072943\\
2.4382	0.0070588\\
2.4349	0.006831\\
2.4317	0.0066104\\
2.4285	0.006397\\
2.4252	0.0061905\\
2.422	0.0059906\\
2.4188	0.0057972\\
2.4155	0.0056101\\
2.4123	0.005429\\
2.409	0.0052538\\
2.4058	0.0050843\\
2.4026	0.0049204\\
2.3993	0.0047618\\
2.3961	0.0046083\\
2.3929	0.0044599\\
2.3896	0.0043163\\
2.3864	0.0041774\\
2.3831	0.0040431\\
2.3799	0.0039132\\
2.3767	0.0037875\\
2.3734	0.0036659\\
2.3702	0.0035484\\
2.367	0.0034347\\
2.3637	0.0033247\\
2.3605	0.0032184\\
2.3572	0.0031156\\
2.354	0.0030161\\
2.3508	0.0029199\\
2.3475	0.0028269\\
2.3443	0.002737\\
2.341	0.00265\\
2.3378	0.0025658\\
2.3346	0.0024845\\
2.3313	0.0024058\\
2.3281	0.0023297\\
2.3249	0.0022561\\
2.3216	0.002185\\
2.3184	0.0021162\\
2.3151	0.0020496\\
2.3119	0.0019853\\
2.3087	0.001923\\
2.3054	0.0018628\\
2.3022	0.0018046\\
2.299	0.0017483\\
2.2957	0.0016938\\
2.2925	0.0016411\\
2.2892	0.0015901\\
2.286	0.0015408\\
2.2828	0.0014932\\
2.2795	0.001447\\
2.2763	0.0014024\\
2.2731	0.0013592\\
2.2698	0.0013175\\
2.2666	0.0012771\\
2.2633	0.001238\\
2.2601	0.0012001\\
2.2569	0.0011635\\
2.2536	0.0011281\\
2.2504	0.0010938\\
2.2472	0.0010607\\
2.2439	0.0010286\\
2.2407	0.00099748\\
2.2374	0.0009674\\
2.2342	0.00093829\\
2.231	0.00091011\\
2.2277	0.00088282\\
2.2245	0.00085641\\
2.2212	0.00083084\\
2.218	0.00080608\\
2.2148	0.0007821\\
2.2115	0.00075888\\
2.2083	0.0007364\\
2.2051	0.00071462\\
2.2018	0.00069353\\
2.1986	0.0006731\\
2.1953	0.0006533\\
2.1921	0.00063413\\
2.1889	0.00061555\\
2.1856	0.00059756\\
2.1824	0.00058012\\
2.1792	0.00056322\\
2.1759	0.00054684\\
2.1727	0.00053097\\
2.1694	0.00051558\\
2.1662	0.00050067\\
2.163	0.00048621\\
2.1597	0.0004722\\
2.1565	0.00045861\\
2.1533	0.00044544\\
2.15	0.00043266\\
2.1468	0.00042027\\
2.1435	0.00040826\\
2.1403	0.0003966\\
2.1371	0.0003853\\
2.1338	0.00037433\\
2.1306	0.0003637\\
2.1274	0.00035338\\
2.1241	0.00034337\\
2.1209	0.00033365\\
2.1176	0.00032422\\
2.1144	0.00031507\\
2.1112	0.00030619\\
2.1079	0.00029757\\
2.1047	0.00028921\\
2.1014	0.00028109\\
2.0982	0.0002732\\
2.095	0.00026555\\
2.0917	0.00025812\\
2.0885	0.0002509\\
2.0853	0.00024389\\
2.082	0.00023709\\
2.0788	0.00023048\\
2.0755	0.00022406\\
}--cycle;
\addplot[ybar interval, fill=black, fill opacity=0.35, area legend, draw=none] table[row sep=crcr, x=Lower, y=Count] {%
Lower	Upper	Count\\
2.0755	2.3244	0.075022\\
2.3244	2.5732	0.13933\\
2.5732	2.822	0.35368\\
2.822	3.0708	0.76094\\
3.0708	3.3196	0.94314\\
3.3196	3.5684	0.75022\\
3.5684	3.8172	0.52516\\
3.8172	4.0661	0.2465\\
4.0661	4.3149	0.10717\\
4.3149	4.5637	0.053587\\
4.5637	4.8125	0.04287\\
4.8125	5.0613	0\\
5.0613	5.3101	0.021435\\
5.3101	5.3101	0.021435\\
};
\addlegendentry{Istogramma reale}

\addplot [color=black]
  table[row sep=crcr]{%
2.0755	0.016864\\
2.1144	0.021976\\
2.15	0.027791\\
2.1824	0.034175\\
2.2115	0.040944\\
2.2407	0.048804\\
2.2666	0.056806\\
2.2925	0.065853\\
2.3184	0.076034\\
2.3443	0.087435\\
2.3702	0.10014\\
2.3929	0.11239\\
2.4155	0.12576\\
2.4382	0.14028\\
2.4608	0.15599\\
2.4835	0.17292\\
2.5062	0.19111\\
2.5288	0.21056\\
2.5515	0.23127\\
2.5774	0.25648\\
2.6033	0.28329\\
2.6292	0.31165\\
2.6584	0.34529\\
2.6875	0.38062\\
2.7199	0.42161\\
2.7555	0.4684\\
2.8041	0.5341\\
2.8947	0.65712\\
2.9303	0.70335\\
2.9595	0.73941\\
2.9854	0.76973\\
3.008	0.79468\\
3.0307	0.81793\\
3.0501	0.83636\\
3.0696	0.8533\\
3.089	0.86863\\
3.1052	0.88011\\
3.1214	0.89036\\
3.1376	0.89934\\
3.1537	0.907\\
3.1699	0.91332\\
3.1861	0.91826\\
3.1991	0.92122\\
3.212	0.92327\\
3.225	0.92442\\
3.2379	0.92466\\
3.2509	0.92399\\
3.2638	0.92242\\
3.2768	0.91996\\
3.2897	0.91661\\
3.3059	0.91119\\
3.3221	0.90444\\
3.3383	0.89637\\
3.3545	0.88705\\
3.3707	0.8765\\
3.3869	0.86479\\
3.4063	0.84927\\
3.4257	0.83225\\
3.4451	0.81384\\
3.4678	0.79078\\
3.4905	0.76619\\
3.5164	0.73648\\
3.5455	0.70139\\
3.5811	0.65673\\
3.6362	0.58564\\
3.7042	0.49801\\
3.743	0.44963\\
3.7754	0.4109\\
3.8045	0.37756\\
3.8337	0.34584\\
3.8596	0.31913\\
3.8855	0.29389\\
3.9114	0.27016\\
3.9373	0.24797\\
3.9632	0.22732\\
3.9891	0.20819\\
4.015	0.19056\\
4.0409	0.17437\\
4.0668	0.15957\\
4.0927	0.14609\\
4.1186	0.13385\\
4.1445	0.12277\\
4.1704	0.11277\\
4.1996	0.10269\\
4.2287	0.093758\\
4.2611	0.085015\\
4.2935	0.077372\\
4.3291	0.070062\\
4.3679	0.063196\\
4.41	0.056832\\
4.4586	0.050596\\
4.5136	0.044623\\
4.5784	0.03868\\
4.6593	0.032391\\
4.7565	0.025978\\
4.8666	0.019813\\
4.9799	0.014524\\
5.0964	0.010153\\
5.2162	0.006727\\
5.3101	0.0047195\\
};
\addlegendentry{Adatt. BLN reale}

            % \input{../../../.tex/20/KS/SizeDistributionFittingsfig3.tex}
        \end{groupplot}
    
        \begin{groupplot}[
            group style={
                group name=Row3,
                plotGroup,
            },
            plotRow,
            row2Keys,
            plotLegendNE,
            %
            xmin=3.4,xmax=5.4,
            ymin=0.0053191,ymax=1,
        ]
            \nextgroupplot[%
                at={(Row2 c1r1.south east)},yshift=-\yPlotSep em,anchor=north east,
                % xmin=3.4,xmax=5.4,
                % ymin=0.0053191,ymax=1,
            ] % This file was created by matlab2tikz.
%
\definecolor{mycolor1}{rgb}{0.00000,0.44706,0.69804}%
\definecolor{mycolor2}{rgb}{0.00000,0.44700,0.74100}%
\definecolor{mycolor3}{rgb}{0.49400,0.18400,0.55600}%
%
\addplot[only marks, mark=*, mark size=1.1180pt, color=mycolor1, fill=mycolor1, opacity=0.60, draw=none, mark options={draw=none,line width=0pt}] table[row sep=crcr]{%
x	y\\
3.595	0.99468\\
3.5991	0.98404\\
3.603	0.9734\\
3.6072	0.96277\\
3.6112	0.95213\\
3.616	0.94149\\
3.6206	0.93085\\
3.6244	0.92021\\
3.628	0.90957\\
3.632	0.89894\\
3.6361	0.8883\\
3.6414	0.87766\\
3.6463	0.86702\\
3.6506	0.85638\\
3.6554	0.84574\\
3.6603	0.83511\\
3.6654	0.82447\\
3.6705	0.81383\\
3.6755	0.80319\\
3.6801	0.79255\\
3.6846	0.78191\\
3.689	0.77128\\
3.6953	0.76064\\
3.7007	0.75\\
3.7071	0.73936\\
3.7137	0.72872\\
3.719	0.71809\\
3.7252	0.70745\\
3.7319	0.69681\\
3.7385	0.68617\\
3.7446	0.67553\\
3.7513	0.66489\\
3.7579	0.65426\\
3.7637	0.64362\\
3.7697	0.63298\\
3.7771	0.62234\\
3.7851	0.6117\\
3.792	0.60106\\
3.7972	0.59043\\
3.8042	0.57979\\
3.8108	0.56915\\
3.8187	0.55851\\
3.8253	0.54787\\
3.8332	0.53723\\
3.842	0.5266\\
3.8503	0.51596\\
3.8578	0.50532\\
3.867	0.49468\\
3.8745	0.48404\\
3.8835	0.4734\\
3.8932	0.46277\\
3.9021	0.45213\\
3.9128	0.44149\\
3.9218	0.43085\\
3.9333	0.42021\\
3.9438	0.40957\\
3.9556	0.39894\\
3.9664	0.3883\\
3.9784	0.37766\\
3.9905	0.36702\\
4.0044	0.35638\\
4.0158	0.34574\\
4.0305	0.33511\\
4.0434	0.32447\\
4.0548	0.31383\\
4.0668	0.30319\\
4.0796	0.29255\\
4.0945	0.28191\\
4.1105	0.27128\\
4.1247	0.26064\\
4.1395	0.25\\
4.1526	0.23936\\
4.1668	0.22872\\
4.1835	0.21809\\
4.2006	0.20745\\
4.2187	0.19681\\
4.234	0.18617\\
4.25	0.17553\\
4.2681	0.16489\\
4.2855	0.15426\\
4.3055	0.14362\\
4.3291	0.13298\\
4.3464	0.12234\\
4.3699	0.1117\\
4.3915	0.10106\\
4.4131	0.090426\\
4.4393	0.079787\\
4.469	0.069149\\
4.5028	0.058511\\
4.5315	0.047872\\
4.5633	0.037234\\
4.6027	0.026596\\
4.6624	0.015957\\
4.7723	0.0053191\\
};
\addlegendentry{FRC empirica esatta}

\addplot [color=mycolor1, only marks, every error bar/.append style={opacity=0.30}, mark=*, mark size=0pt, draw=none, forget plot]
 plot [error bars/.cd, x dir=both, x explicit, error bar style={line width=1pt, color=mycolor1}, error mark options={mark=none,mark size=0pt}]
 table[row sep=crcr, x error plus index=2, x error minus index=3]{%
3.595	0.99468	0.016766	0.016766\\
3.5991	0.98404	0.016782	0.016782\\
3.603	0.9734	0.016717	0.016717\\
3.6072	0.96277	0.016839	0.016839\\
3.6112	0.95213	0.016912	0.016912\\
3.616	0.94149	0.016953	0.016953\\
3.6206	0.93085	0.01706	0.01706\\
3.6244	0.92021	0.017084	0.017084\\
3.628	0.90957	0.017009	0.017009\\
3.632	0.89894	0.017084	0.017084\\
3.6361	0.8883	0.017114	0.017114\\
3.6414	0.87766	0.017155	0.017155\\
3.6463	0.86702	0.017092	0.017092\\
3.6506	0.85638	0.01694	0.01694\\
3.6554	0.84574	0.016867	0.016867\\
3.6603	0.83511	0.016956	0.016956\\
3.6654	0.82447	0.016933	0.016933\\
3.6705	0.81383	0.017048	0.017048\\
3.6755	0.80319	0.017097	0.017097\\
3.6801	0.79255	0.016913	0.016913\\
3.6846	0.78191	0.01693	0.01693\\
3.689	0.77128	0.017115	0.017115\\
3.6953	0.76064	0.017195	0.017195\\
3.7007	0.75	0.017192	0.017192\\
3.7071	0.73936	0.017346	0.017346\\
3.7137	0.72872	0.017581	0.017581\\
3.719	0.71809	0.01758	0.01758\\
3.7252	0.70745	0.017585	0.017585\\
3.7319	0.69681	0.017497	0.017497\\
3.7385	0.68617	0.017381	0.017381\\
3.7446	0.67553	0.017427	0.017427\\
3.7513	0.66489	0.017893	0.017893\\
3.7579	0.65426	0.017745	0.017745\\
3.7637	0.64362	0.017909	0.017909\\
3.7697	0.63298	0.018096	0.018096\\
3.7771	0.62234	0.017864	0.017864\\
3.7851	0.6117	0.01782	0.01782\\
3.792	0.60106	0.017856	0.017856\\
3.7972	0.59043	0.017793	0.017793\\
3.8042	0.57979	0.017565	0.017565\\
3.8108	0.56915	0.01747	0.01747\\
3.8187	0.55851	0.017521	0.017521\\
3.8253	0.54787	0.017571	0.017571\\
3.8332	0.53723	0.01761	0.01761\\
3.842	0.5266	0.017466	0.017466\\
3.8503	0.51596	0.017322	0.017322\\
3.8578	0.50532	0.017264	0.017264\\
3.867	0.49468	0.017161	0.017161\\
3.8745	0.48404	0.017213	0.017213\\
3.8835	0.4734	0.01698	0.01698\\
3.8932	0.46277	0.01723	0.01723\\
3.9021	0.45213	0.017293	0.017293\\
3.9128	0.44149	0.017509	0.017509\\
3.9218	0.43085	0.017639	0.017639\\
3.9333	0.42021	0.017595	0.017595\\
3.9438	0.40957	0.01781	0.01781\\
3.9556	0.39894	0.017385	0.017385\\
3.9664	0.3883	0.01757	0.01757\\
3.9784	0.37766	0.017472	0.017472\\
3.9905	0.36702	0.017342	0.017342\\
4.0044	0.35638	0.017587	0.017587\\
4.0158	0.34574	0.017576	0.017576\\
4.0305	0.33511	0.017734	0.017734\\
4.0434	0.32447	0.01792	0.01792\\
4.0548	0.31383	0.018094	0.018094\\
4.0668	0.30319	0.018204	0.018204\\
4.0796	0.29255	0.017912	0.017912\\
4.0945	0.28191	0.017799	0.017799\\
4.1105	0.27128	0.017972	0.017972\\
4.1247	0.26064	0.017952	0.017952\\
4.1395	0.25	0.018059	0.018059\\
4.1526	0.23936	0.01802	0.01802\\
4.1668	0.22872	0.017883	0.017883\\
4.1835	0.21809	0.0181	0.0181\\
4.2006	0.20745	0.018599	0.018599\\
4.2187	0.19681	0.018791	0.018791\\
4.234	0.18617	0.019126	0.019126\\
4.25	0.17553	0.019253	0.019253\\
4.2681	0.16489	0.019124	0.019124\\
4.2855	0.15426	0.018721	0.018721\\
4.3055	0.14362	0.018333	0.018333\\
4.3291	0.13298	0.01891	0.01891\\
4.3464	0.12234	0.019257	0.019257\\
4.3699	0.1117	0.01886	0.01886\\
4.3915	0.10106	0.018397	0.018397\\
4.4131	0.090426	0.019204	0.019204\\
4.4393	0.079787	0.019281	0.019281\\
4.469	0.069149	0.020489	0.020489\\
4.5028	0.058511	0.020691	0.020691\\
4.5315	0.047872	0.021265	0.021265\\
4.5633	0.037234	0.021561	0.021561\\
4.6027	0.026596	0.022743	0.022743\\
4.6624	0.015957	0.024013	0.024013\\
4.7723	0.0053191	0.028039	0.028039\\
};
\addplot [color=mycolor1]
  table[row sep=crcr]{%
3.6074	0.89722\\
3.6246	0.87254\\
3.6418	0.84594\\
3.659	0.81797\\
3.6762	0.78757\\
3.6934	0.75589\\
3.7278	0.69272\\
3.7622	0.63288\\
3.7966	0.576\\
3.8826	0.45231\\
4.347	0.12087\\
4.8286	0.030889\\
5.3101	0.0079275\\
};
\addlegendentry{Adatt. Pareto esatto}


\addplot[area legend, draw=none, fill=mycolor1, fill opacity=0.15, forget plot]
table[row sep=crcr] {%
x	y\\
3.6074	0.92101\\
3.6246	0.89854\\
3.6418	0.87384\\
3.659	0.84743\\
3.6762	0.81789\\
3.6934	0.78655\\
3.7106	0.75465\\
3.7278	0.72316\\
3.745	0.69254\\
3.7622	0.66265\\
3.7794	0.63331\\
3.7966	0.60416\\
3.8138	0.5761\\
3.831	0.54941\\
3.8482	0.52403\\
3.8654	0.49924\\
3.8826	0.47541\\
3.8998	0.45272\\
3.917	0.43112\\
3.9342	0.41055\\
3.9514	0.39097\\
3.9686	0.37232\\
3.9858	0.35457\\
4.003	0.33767\\
4.0202	0.32157\\
4.0374	0.30625\\
4.0546	0.29166\\
4.0718	0.27777\\
4.089	0.26454\\
4.1062	0.25194\\
4.1234	0.23995\\
4.1406	0.22853\\
4.1578	0.21765\\
4.175	0.2073\\
4.1922	0.19744\\
4.2094	0.18805\\
4.2266	0.1791\\
4.2438	0.17059\\
4.261	0.16248\\
4.2782	0.15476\\
4.2954	0.1474\\
4.3126	0.1404\\
4.3298	0.13373\\
4.347	0.12738\\
4.3642	0.12133\\
4.3814	0.11557\\
4.3986	0.11009\\
4.4158	0.10486\\
4.433	0.099888\\
4.4502	0.095149\\
4.4674	0.090636\\
4.4846	0.086338\\
4.5018	0.082244\\
4.519	0.078345\\
4.5362	0.074632\\
4.5534	0.071095\\
4.5706	0.067726\\
4.5878	0.064518\\
4.605	0.061462\\
4.6222	0.058551\\
4.6394	0.055779\\
4.6566	0.053138\\
4.6738	0.050623\\
4.691	0.048227\\
4.7082	0.045945\\
4.7254	0.043771\\
4.7426	0.041701\\
4.7598	0.039728\\
4.777	0.03785\\
4.7942	0.03606\\
4.8114	0.034355\\
4.8286	0.032732\\
4.8458	0.031185\\
4.863	0.029711\\
4.8802	0.028308\\
4.8974	0.02697\\
4.9146	0.025697\\
4.9318	0.024483\\
4.949	0.023327\\
4.9662	0.022226\\
4.9834	0.021177\\
5.0006	0.020177\\
5.0177	0.019225\\
5.0349	0.018318\\
5.0521	0.017454\\
5.0693	0.016631\\
5.0865	0.015847\\
5.1037	0.015099\\
5.1209	0.014388\\
5.1381	0.013709\\
5.1553	0.013063\\
5.1725	0.012448\\
5.1897	0.011861\\
5.2069	0.011302\\
5.2241	0.01077\\
5.2413	0.010263\\
5.2585	0.0097794\\
5.2757	0.0093189\\
5.2929	0.0088802\\
5.3101	0.0084622\\
5.3101	0.0073927\\
5.2929	0.0077626\\
5.2757	0.0081511\\
5.2585	0.0085591\\
5.2413	0.0089874\\
5.2241	0.0094373\\
5.2069	0.0099097\\
5.1897	0.010406\\
5.1725	0.010927\\
5.1553	0.011474\\
5.1381	0.012048\\
5.1209	0.012652\\
5.1037	0.013285\\
5.0865	0.013951\\
5.0693	0.014649\\
5.0521	0.015383\\
5.0349	0.016154\\
5.0177	0.016963\\
5.0006	0.017812\\
4.9834	0.018705\\
4.9662	0.019642\\
4.949	0.020626\\
4.9318	0.02166\\
4.9146	0.022745\\
4.8974	0.023885\\
4.8802	0.025082\\
4.863	0.026339\\
4.8458	0.027659\\
4.8286	0.029046\\
4.8114	0.030502\\
4.7942	0.032031\\
4.777	0.033637\\
4.7598	0.035323\\
4.7426	0.037094\\
4.7254	0.038954\\
4.7082	0.040908\\
4.691	0.042959\\
4.6738	0.045113\\
4.6566	0.047376\\
4.6394	0.049752\\
4.6222	0.052247\\
4.605	0.054868\\
4.5878	0.05762\\
4.5706	0.06051\\
4.5534	0.063546\\
4.5362	0.066734\\
4.519	0.070082\\
4.5018	0.073598\\
4.4846	0.07729\\
4.4674	0.081169\\
4.4502	0.085241\\
4.433	0.089519\\
4.4158	0.094011\\
4.3986	0.098729\\
4.3814	0.10368\\
4.3642	0.10889\\
4.347	0.11435\\
4.3298	0.12009\\
4.3126	0.12612\\
4.2954	0.13245\\
4.2782	0.1391\\
4.261	0.14608\\
4.2438	0.15341\\
4.2266	0.16111\\
4.2094	0.1692\\
4.1922	0.17769\\
4.175	0.18661\\
4.1578	0.19598\\
4.1406	0.20582\\
4.1234	0.21616\\
4.1062	0.22701\\
4.089	0.23841\\
4.0718	0.25038\\
4.0546	0.26295\\
4.0374	0.27615\\
4.0202	0.29002\\
4.003	0.30458\\
3.9858	0.31988\\
3.9686	0.33594\\
3.9514	0.35281\\
3.9342	0.37053\\
3.917	0.38913\\
3.8998	0.40867\\
3.8826	0.4292\\
3.8654	0.45075\\
3.8482	0.47336\\
3.831	0.49701\\
3.8138	0.52182\\
3.7966	0.54784\\
3.7794	0.57503\\
3.7622	0.60312\\
3.745	0.6322\\
3.7278	0.66227\\
3.7106	0.69338\\
3.6934	0.72524\\
3.6762	0.75724\\
3.659	0.78851\\
3.6418	0.81805\\
3.6246	0.84654\\
3.6074	0.87343\\
}--cycle;
\addplot[only marks, mark=*, mark size=1.1180pt, color=black, fill=black, opacity=0.60, draw=none, mark options={draw=none,line width=0pt}] table[row sep=crcr]{%
x	y\\
3.5623	0.99468\\
3.5723	0.98404\\
3.573	0.9734\\
3.5753	0.96277\\
3.5766	0.95213\\
3.5781	0.94149\\
3.5826	0.93085\\
3.5851	0.92021\\
3.5911	0.90957\\
3.5942	0.89894\\
3.596	0.8883\\
3.6016	0.87766\\
3.603	0.86702\\
3.6047	0.85638\\
3.6086	0.84574\\
3.6086	0.83511\\
3.6097	0.82447\\
3.6142	0.81383\\
3.6261	0.80319\\
3.6383	0.79255\\
3.6464	0.78191\\
3.6535	0.77128\\
3.657	0.76064\\
3.6658	0.75\\
3.6678	0.73936\\
3.6687	0.72872\\
3.6702	0.71809\\
3.6723	0.70745\\
3.6793	0.69681\\
3.6867	0.68617\\
3.6949	0.67553\\
3.7015	0.66489\\
3.7031	0.65426\\
3.719	0.64362\\
3.7213	0.63298\\
3.7265	0.62234\\
3.7307	0.6117\\
3.737	0.60106\\
3.7374	0.59043\\
3.7423	0.57979\\
3.7487	0.56915\\
3.7677	0.55851\\
3.7685	0.54787\\
3.7768	0.53723\\
3.8005	0.5266\\
3.8015	0.51596\\
3.8022	0.50532\\
3.8032	0.49468\\
3.8039	0.48404\\
3.8108	0.4734\\
3.8214	0.46277\\
3.8405	0.45213\\
3.8605	0.44149\\
3.863	0.43085\\
3.8645	0.42021\\
3.8662	0.40957\\
3.8806	0.39894\\
3.8878	0.3883\\
3.8964	0.37766\\
3.905	0.36702\\
3.9294	0.35638\\
3.9303	0.34574\\
3.954	0.33511\\
3.9604	0.32447\\
3.9669	0.31383\\
3.9747	0.30319\\
3.9822	0.29255\\
3.9929	0.28191\\
4.0051	0.27128\\
4.0144	0.26064\\
4.0161	0.25\\
4.0433	0.23936\\
4.0578	0.22872\\
4.073	0.21809\\
4.0857	0.20745\\
4.0904	0.19681\\
4.1168	0.18617\\
4.1265	0.17553\\
4.127	0.16489\\
4.143	0.15426\\
4.1756	0.14362\\
4.2156	0.13298\\
4.3116	0.12234\\
4.3276	0.1117\\
4.3662	0.10106\\
4.4791	0.090426\\
4.4912	0.079787\\
4.517	0.069149\\
4.5743	0.058511\\
4.5914	0.047872\\
4.6138	0.037234\\
4.7898	0.026596\\
5.0876	0.015957\\
5.3101	0.0053191\\
};
\addlegendentry{FRC empirica reale}

\addplot [color=black]
  table[row sep=crcr]{%
3.5608	1\\
5.3101	0.005939\\
};
\addlegendentry{Adatt. Pareto reale}

            % \input{../../../.tex/20/KS/SizeDistributionFittingsfig8.tex}
    
            \nextgroupplot[%
                % xmin=3.4,xmax=5.4,
                % ymin=0.0053191,ymax=1,
            ] % This file was created by matlab2tikz.
%
\definecolor{mycolor1}{rgb}{0.00000,0.44706,0.69804}%
\definecolor{mycolor2}{rgb}{0.00000,0.44700,0.74100}%
\definecolor{mycolor3}{rgb}{0.49400,0.18400,0.55600}%
%
\addplot[only marks, mark=*, mark size=1.1180pt, color=mycolor1, fill=mycolor1, opacity=0.60, draw=none, mark options={draw=none,line width=0pt}] table[row sep=crcr]{%
x	y\\
3.6059	0.99468\\
3.6099	0.98404\\
3.6137	0.9734\\
3.6172	0.96277\\
3.6211	0.95213\\
3.6248	0.94149\\
3.6293	0.93085\\
3.6333	0.92021\\
3.6373	0.90957\\
3.6417	0.89894\\
3.6459	0.8883\\
3.6501	0.87766\\
3.6549	0.86702\\
3.6599	0.85638\\
3.6645	0.84574\\
3.6688	0.83511\\
3.6732	0.82447\\
3.6783	0.81383\\
3.6835	0.80319\\
3.6888	0.79255\\
3.6945	0.78191\\
3.6991	0.77128\\
3.7049	0.76064\\
3.7098	0.75\\
3.7158	0.73936\\
3.7213	0.72872\\
3.7272	0.71809\\
3.7335	0.70745\\
3.7397	0.69681\\
3.7454	0.68617\\
3.7522	0.67553\\
3.7594	0.66489\\
3.7665	0.65426\\
3.7732	0.64362\\
3.7794	0.63298\\
3.7864	0.62234\\
3.7933	0.6117\\
3.7997	0.60106\\
3.8082	0.59043\\
3.8166	0.57979\\
3.8247	0.56915\\
3.8323	0.55851\\
3.8389	0.54787\\
3.8468	0.53723\\
3.8561	0.5266\\
3.8655	0.51596\\
3.8746	0.50532\\
3.8839	0.49468\\
3.8931	0.48404\\
3.903	0.4734\\
3.9128	0.46277\\
3.9236	0.45213\\
3.9347	0.44149\\
3.9445	0.43085\\
3.9551	0.42021\\
3.9678	0.40957\\
3.9788	0.39894\\
3.9925	0.3883\\
4.0046	0.37766\\
4.0168	0.36702\\
4.0286	0.35638\\
4.0404	0.34574\\
4.0558	0.33511\\
4.0684	0.32447\\
4.0836	0.31383\\
4.099	0.30319\\
4.1131	0.29255\\
4.1277	0.28191\\
4.1421	0.27128\\
4.1543	0.26064\\
4.1683	0.25\\
4.1839	0.23936\\
4.1981	0.22872\\
4.2149	0.21809\\
4.231	0.20745\\
4.2506	0.19681\\
4.2659	0.18617\\
4.2858	0.17553\\
4.3036	0.16489\\
4.3185	0.15426\\
4.334	0.14362\\
4.3523	0.13298\\
4.3721	0.12234\\
4.3902	0.1117\\
4.4077	0.10106\\
4.4258	0.090426\\
4.4437	0.079787\\
4.4617	0.069149\\
4.4855	0.058511\\
4.5094	0.047872\\
4.5418	0.037234\\
4.5716	0.026596\\
4.6123	0.015957\\
4.6623	0.0053191\\
};
\addlegendentry{FRC empirica esatta}

\addplot [color=mycolor1, only marks, every error bar/.append style={opacity=0.30}, mark=*, mark size=0pt, draw=none, forget plot]
 plot [error bars/.cd, x dir=both, x explicit, error bar style={line width=1pt, color=mycolor1}, error mark options={mark=none,mark size=0pt}]
 table[row sep=crcr, x error plus index=2, x error minus index=3]{%
3.6059	0.99468	0.003094	0.003094\\
3.6099	0.98404	0.0031692	0.0031692\\
3.6137	0.9734	0.0030718	0.0030718\\
3.6172	0.96277	0.0031693	0.0031693\\
3.6211	0.95213	0.0032353	0.0032353\\
3.6248	0.94149	0.0032585	0.0032585\\
3.6293	0.93085	0.003201	0.003201\\
3.6333	0.92021	0.0032879	0.0032879\\
3.6373	0.90957	0.0032566	0.0032566\\
3.6417	0.89894	0.0032233	0.0032233\\
3.6459	0.8883	0.0033792	0.0033792\\
3.6501	0.87766	0.0033426	0.0033426\\
3.6549	0.86702	0.0033991	0.0033991\\
3.6599	0.85638	0.0033465	0.0033465\\
3.6645	0.84574	0.0033965	0.0033965\\
3.6688	0.83511	0.0034419	0.0034419\\
3.6732	0.82447	0.0034482	0.0034482\\
3.6783	0.81383	0.0034786	0.0034786\\
3.6835	0.80319	0.0036611	0.0036611\\
3.6888	0.79255	0.0035673	0.0035673\\
3.6945	0.78191	0.0036178	0.0036178\\
3.6991	0.77128	0.0036915	0.0036915\\
3.7049	0.76064	0.0039275	0.0039275\\
3.7098	0.75	0.0040974	0.0040974\\
3.7158	0.73936	0.0042161	0.0042161\\
3.7213	0.72872	0.0040783	0.0040783\\
3.7272	0.71809	0.0041318	0.0041318\\
3.7335	0.70745	0.0042297	0.0042297\\
3.7397	0.69681	0.004108	0.004108\\
3.7454	0.68617	0.0041197	0.0041197\\
3.7522	0.67553	0.0043036	0.0043036\\
3.7594	0.66489	0.0042925	0.0042925\\
3.7665	0.65426	0.0044916	0.0044916\\
3.7732	0.64362	0.0044742	0.0044742\\
3.7794	0.63298	0.0044369	0.0044369\\
3.7864	0.62234	0.004425	0.004425\\
3.7933	0.6117	0.0042513	0.0042513\\
3.7997	0.60106	0.0043225	0.0043225\\
3.8082	0.59043	0.0043351	0.0043351\\
3.8166	0.57979	0.0043132	0.0043132\\
3.8247	0.56915	0.0042772	0.0042772\\
3.8323	0.55851	0.0039948	0.0039948\\
3.8389	0.54787	0.0041372	0.0041372\\
3.8468	0.53723	0.0043442	0.0043442\\
3.8561	0.5266	0.0045954	0.0045954\\
3.8655	0.51596	0.0045341	0.0045341\\
3.8746	0.50532	0.0047735	0.0047735\\
3.8839	0.49468	0.0046397	0.0046397\\
3.8931	0.48404	0.0046701	0.0046701\\
3.903	0.4734	0.0046291	0.0046291\\
3.9128	0.46277	0.0047467	0.0047467\\
3.9236	0.45213	0.0048018	0.0048018\\
3.9347	0.44149	0.0049219	0.0049219\\
3.9445	0.43085	0.0046167	0.0046167\\
3.9551	0.42021	0.0045622	0.0045622\\
3.9678	0.40957	0.0045148	0.0045148\\
3.9788	0.39894	0.004715	0.004715\\
3.9925	0.3883	0.0051674	0.0051674\\
4.0046	0.37766	0.0056007	0.0056007\\
4.0168	0.36702	0.0055549	0.0055549\\
4.0286	0.35638	0.0055284	0.0055284\\
4.0404	0.34574	0.0053456	0.0053456\\
4.0558	0.33511	0.0055047	0.0055047\\
4.0684	0.32447	0.0056848	0.0056848\\
4.0836	0.31383	0.0053806	0.0053806\\
4.099	0.30319	0.0050695	0.0050695\\
4.1131	0.29255	0.0049377	0.0049377\\
4.1277	0.28191	0.0054952	0.0054952\\
4.1421	0.27128	0.0055687	0.0055687\\
4.1543	0.26064	0.0056016	0.0056016\\
4.1683	0.25	0.0058282	0.0058282\\
4.1839	0.23936	0.0056262	0.0056262\\
4.1981	0.22872	0.0053676	0.0053676\\
4.2149	0.21809	0.0053149	0.0053149\\
4.231	0.20745	0.0053388	0.0053388\\
4.2506	0.19681	0.0052303	0.0052303\\
4.2659	0.18617	0.0058319	0.0058319\\
4.2858	0.17553	0.0060002	0.0060002\\
4.3036	0.16489	0.006501	0.006501\\
4.3185	0.15426	0.0064755	0.0064755\\
4.334	0.14362	0.0060002	0.0060002\\
4.3523	0.13298	0.0058033	0.0058033\\
4.3721	0.12234	0.005256	0.005256\\
4.3902	0.1117	0.0047916	0.0047916\\
4.4077	0.10106	0.0050181	0.0050181\\
4.4258	0.090426	0.0048308	0.0048308\\
4.4437	0.079787	0.0050423	0.0050423\\
4.4617	0.069149	0.0052417	0.0052417\\
4.4855	0.058511	0.0057599	0.0057599\\
4.5094	0.047872	0.0052978	0.0052978\\
4.5418	0.037234	0.0062925	0.0062925\\
4.5716	0.026596	0.0063138	0.0063138\\
4.6123	0.015957	0.0068066	0.0068066\\
4.6623	0.0053191	0.0082205	0.0082205\\
};
\addplot [color=mycolor1]
  table[row sep=crcr]{%
3.6041	0.98281\\
3.6214	0.95075\\
3.6386	0.90801\\
3.7592	0.64579\\
4.3107	0.13612\\
4.8621	0.028809\\
5.3101	0.0081826\\
};
\addlegendentry{Adatt. Pareto esatto}


\addplot[area legend, draw=none, fill=mycolor1, fill opacity=0.15, forget plot]
table[row sep=crcr] {%
x	y\\
3.6041	0.98731\\
3.6214	0.95761\\
3.6386	0.91511\\
3.6558	0.87172\\
3.6731	0.82996\\
3.6903	0.7902\\
3.7075	0.75236\\
3.7248	0.71634\\
3.742	0.68204\\
3.7592	0.6494\\
3.7765	0.61833\\
3.7937	0.58876\\
3.8109	0.56062\\
3.8282	0.53383\\
3.8454	0.50833\\
3.8626	0.48407\\
3.8799	0.46098\\
3.8971	0.439\\
3.9143	0.4181\\
3.9315	0.39821\\
3.9488	0.37928\\
3.966	0.36127\\
3.9832	0.34413\\
4.0005	0.32782\\
4.0177	0.3123\\
4.0349	0.29753\\
4.0522	0.28346\\
4.0694	0.27006\\
4.0866	0.25731\\
4.1039	0.24516\\
4.1211	0.23359\\
4.1383	0.22257\\
4.1556	0.21208\\
4.1728	0.20208\\
4.19	0.19255\\
4.2073	0.18348\\
4.2245	0.17483\\
4.2417	0.1666\\
4.259	0.15875\\
4.2762	0.15127\\
4.2934	0.14415\\
4.3107	0.13736\\
4.3279	0.1309\\
4.3451	0.12473\\
4.3624	0.11886\\
4.3796	0.11327\\
4.3968	0.10794\\
4.4141	0.10286\\
4.4313	0.098022\\
4.4485	0.093412\\
4.4658	0.089018\\
4.483	0.084831\\
4.5002	0.080842\\
4.5174	0.077041\\
4.5347	0.073419\\
4.5519	0.069967\\
4.5691	0.066678\\
4.5864	0.063544\\
4.6036	0.060557\\
4.6208	0.057711\\
4.6381	0.054999\\
4.6553	0.052415\\
4.6725	0.049952\\
4.6898	0.047605\\
4.707	0.045369\\
4.7242	0.043238\\
4.7415	0.041207\\
4.7587	0.039272\\
4.7759	0.037428\\
4.7932	0.03567\\
4.8104	0.033995\\
4.8276	0.032399\\
4.8449	0.030878\\
4.8621	0.029429\\
4.8793	0.028048\\
4.8966	0.026731\\
4.9138	0.025477\\
4.931	0.024281\\
4.9483	0.023142\\
4.9655	0.022056\\
4.9827	0.021021\\
5	0.020035\\
5.0172	0.019095\\
5.0344	0.0182\\
5.0516	0.017346\\
5.0689	0.016533\\
5.0861	0.015757\\
5.1033	0.015018\\
5.1206	0.014314\\
5.1378	0.013643\\
5.155	0.013004\\
5.1723	0.012394\\
5.1895	0.011813\\
5.2067	0.01126\\
5.224	0.010732\\
5.2412	0.010229\\
5.2584	0.0097497\\
5.2757	0.0092929\\
5.2929	0.0088575\\
5.3101	0.0084426\\
5.3101	0.0079225\\
5.2929	0.0083185\\
5.2757	0.0087342\\
5.2584	0.0091708\\
5.2412	0.0096293\\
5.224	0.010111\\
5.2067	0.010616\\
5.1895	0.011147\\
5.1723	0.011704\\
5.155	0.01229\\
5.1378	0.012904\\
5.1206	0.01355\\
5.1033	0.014228\\
5.0861	0.014939\\
5.0689	0.015687\\
5.0516	0.016472\\
5.0344	0.017296\\
5.0172	0.018162\\
5	0.01907\\
4.9827	0.020025\\
4.9655	0.021027\\
4.9483	0.02208\\
4.931	0.023185\\
4.9138	0.024346\\
4.8966	0.025565\\
4.8793	0.026845\\
4.8621	0.028189\\
4.8449	0.029601\\
4.8276	0.031083\\
4.8104	0.03264\\
4.7932	0.034275\\
4.7759	0.035992\\
4.7587	0.037796\\
4.7415	0.039689\\
4.7242	0.041678\\
4.707	0.043766\\
4.6898	0.04596\\
4.6725	0.048263\\
4.6553	0.050682\\
4.6381	0.053223\\
4.6208	0.055891\\
4.6036	0.058693\\
4.5864	0.061635\\
4.5691	0.064726\\
4.5519	0.067971\\
4.5347	0.07138\\
4.5174	0.074959\\
4.5002	0.078719\\
4.483	0.082667\\
4.4658	0.086813\\
4.4485	0.091168\\
4.4313	0.095742\\
4.4141	0.10054\\
4.3968	0.10559\\
4.3796	0.11089\\
4.3624	0.11645\\
4.3451	0.12229\\
4.3279	0.12843\\
4.3107	0.13488\\
4.2934	0.14164\\
4.2762	0.14875\\
4.259	0.15622\\
4.2417	0.16406\\
4.2245	0.17229\\
4.2073	0.18094\\
4.19	0.19002\\
4.1728	0.19955\\
4.1556	0.20956\\
4.1383	0.22008\\
4.1211	0.23112\\
4.1039	0.24271\\
4.0866	0.25488\\
4.0694	0.26765\\
4.0522	0.28106\\
4.0349	0.29514\\
4.0177	0.30991\\
4.0005	0.32542\\
3.9832	0.34168\\
3.966	0.35875\\
3.9488	0.37665\\
3.9315	0.39543\\
3.9143	0.41513\\
3.8971	0.4358\\
3.8799	0.45747\\
3.8626	0.48021\\
3.8454	0.50407\\
3.8282	0.52909\\
3.8109	0.55535\\
3.7937	0.58291\\
3.7765	0.61183\\
3.7592	0.64217\\
3.742	0.67401\\
3.7248	0.70743\\
3.7075	0.74251\\
3.6903	0.77932\\
3.6731	0.81795\\
3.6558	0.8585\\
3.6386	0.90092\\
3.6214	0.9439\\
3.6041	0.97831\\
}--cycle;
\addplot[only marks, mark=*, mark size=1.1180pt, color=black, fill=black, opacity=0.60, draw=none, mark options={draw=none,line width=0pt}] table[row sep=crcr]{%
x	y\\
3.5623	0.99468\\
3.5723	0.98404\\
3.573	0.9734\\
3.5753	0.96277\\
3.5766	0.95213\\
3.5781	0.94149\\
3.5826	0.93085\\
3.5851	0.92021\\
3.5911	0.90957\\
3.5942	0.89894\\
3.596	0.8883\\
3.6016	0.87766\\
3.603	0.86702\\
3.6047	0.85638\\
3.6086	0.84574\\
3.6086	0.83511\\
3.6097	0.82447\\
3.6142	0.81383\\
3.6261	0.80319\\
3.6383	0.79255\\
3.6464	0.78191\\
3.6535	0.77128\\
3.657	0.76064\\
3.6658	0.75\\
3.6678	0.73936\\
3.6687	0.72872\\
3.6702	0.71809\\
3.6723	0.70745\\
3.6793	0.69681\\
3.6867	0.68617\\
3.6949	0.67553\\
3.7015	0.66489\\
3.7031	0.65426\\
3.719	0.64362\\
3.7213	0.63298\\
3.7265	0.62234\\
3.7307	0.6117\\
3.737	0.60106\\
3.7374	0.59043\\
3.7423	0.57979\\
3.7487	0.56915\\
3.7677	0.55851\\
3.7685	0.54787\\
3.7768	0.53723\\
3.8005	0.5266\\
3.8015	0.51596\\
3.8022	0.50532\\
3.8032	0.49468\\
3.8039	0.48404\\
3.8108	0.4734\\
3.8214	0.46277\\
3.8405	0.45213\\
3.8605	0.44149\\
3.863	0.43085\\
3.8645	0.42021\\
3.8662	0.40957\\
3.8806	0.39894\\
3.8878	0.3883\\
3.8964	0.37766\\
3.905	0.36702\\
3.9294	0.35638\\
3.9303	0.34574\\
3.954	0.33511\\
3.9604	0.32447\\
3.9669	0.31383\\
3.9747	0.30319\\
3.9822	0.29255\\
3.9929	0.28191\\
4.0051	0.27128\\
4.0144	0.26064\\
4.0161	0.25\\
4.0433	0.23936\\
4.0578	0.22872\\
4.073	0.21809\\
4.0857	0.20745\\
4.0904	0.19681\\
4.1168	0.18617\\
4.1265	0.17553\\
4.127	0.16489\\
4.143	0.15426\\
4.1756	0.14362\\
4.2156	0.13298\\
4.3116	0.12234\\
4.3276	0.1117\\
4.3662	0.10106\\
4.4791	0.090426\\
4.4912	0.079787\\
4.517	0.069149\\
4.5743	0.058511\\
4.5914	0.047872\\
4.6138	0.037234\\
4.7898	0.026596\\
5.0876	0.015957\\
5.3101	0.0053191\\
};
\addlegendentry{FRC empirica reale}

\addplot [color=black]
  table[row sep=crcr]{%
3.5608	1\\
5.3101	0.005939\\
};
\addlegendentry{Adatt. Pareto reale}

            % \input{../../../.tex/20/KS/SizeDistributionFittingsfig8.tex}
        \end{groupplot}
    
        \begin{groupplot}[
            group style={
                group name=Row4,
                plotLastGroup,
            },
            plotRow,
            row3Keys,
            plotLegendNE,
            %
            xmin=3.4,xmax=5.4,
            ymin=0.0053191,ymax=1,
        ]
            \nextgroupplot[%
                at={($(Row3 c1r1.south west)-(0,.5*\yPlotSep em)-(0,\plotHeight cm)$)},
                % xmin=3.4,xmax=5.4,
                % ymin=0.0053191,ymax=1,
            ] % This file was created by matlab2tikz.
%
\definecolor{mycolor1}{rgb}{0.83529,0.36863,0.00000}%
\definecolor{mycolor2}{rgb}{0.00000,0.44700,0.74100}%
\definecolor{mycolor3}{rgb}{0.49400,0.18400,0.55600}%
%
\addplot[only marks, mark=*, mark size=1.1180pt, color=mycolor1, fill=mycolor1, opacity=0.60, draw=none, mark options={draw=none,line width=0pt}] table[row sep=crcr]{%
x	y\\
3.6194	0.99468\\
3.6232	0.98404\\
3.6278	0.9734\\
3.6314	0.96277\\
3.6357	0.95213\\
3.6392	0.94149\\
3.6437	0.93085\\
3.6471	0.92021\\
3.6516	0.90957\\
3.6566	0.89894\\
3.6607	0.8883\\
3.6653	0.87766\\
3.6698	0.86702\\
3.674	0.85638\\
3.6787	0.84574\\
3.6822	0.83511\\
3.6863	0.82447\\
3.6915	0.81383\\
3.6972	0.80319\\
3.7027	0.79255\\
3.7086	0.78191\\
3.7141	0.77128\\
3.7198	0.76064\\
3.7264	0.75\\
3.7327	0.73936\\
3.7381	0.72872\\
3.7433	0.71809\\
3.7488	0.70745\\
3.7555	0.69681\\
3.7601	0.68617\\
3.766	0.67553\\
3.7711	0.66489\\
3.7773	0.65426\\
3.7848	0.64362\\
3.7921	0.63298\\
3.7981	0.62234\\
3.8059	0.6117\\
3.813	0.60106\\
3.8201	0.59043\\
3.8259	0.57979\\
3.8325	0.56915\\
3.8406	0.55851\\
3.848	0.54787\\
3.8549	0.53723\\
3.8639	0.5266\\
3.8714	0.51596\\
3.8788	0.50532\\
3.8859	0.49468\\
3.8938	0.48404\\
3.9011	0.4734\\
3.9113	0.46277\\
3.9198	0.45213\\
3.9301	0.44149\\
3.9386	0.43085\\
3.9481	0.42021\\
3.9586	0.40957\\
3.968	0.39894\\
3.9784	0.3883\\
3.987	0.37766\\
3.9966	0.36702\\
4.0062	0.35638\\
4.0169	0.34574\\
4.0263	0.33511\\
4.0359	0.32447\\
4.0484	0.31383\\
4.0607	0.30319\\
4.0709	0.29255\\
4.085	0.28191\\
4.098	0.27128\\
4.1118	0.26064\\
4.126	0.25\\
4.1402	0.23936\\
4.1546	0.22872\\
4.1704	0.21809\\
4.1874	0.20745\\
4.2028	0.19681\\
4.2165	0.18617\\
4.236	0.17553\\
4.2531	0.16489\\
4.27	0.15426\\
4.2881	0.14362\\
4.304	0.13298\\
4.3264	0.12234\\
4.3455	0.1117\\
4.3679	0.10106\\
4.395	0.090426\\
4.4211	0.079787\\
4.4498	0.069149\\
4.4855	0.058511\\
4.527	0.047872\\
4.5704	0.037234\\
4.6275	0.026596\\
4.6955	0.015957\\
4.8084	0.0053191\\
};
\addlegendentry{FRC empirica appr.}

\addplot [color=mycolor1, only marks, every error bar/.append style={opacity=0.30}, mark=*, mark size=0pt, draw=none, forget plot]
 plot [error bars/.cd, x dir=both, x explicit, error bar style={line width=1pt, color=mycolor1}, error mark options={mark=none,mark size=0pt}]
 table[row sep=crcr, x error plus index=2, x error minus index=3]{%
3.6194	0.99468	0.017772	0.017772\\
3.6232	0.98404	0.017681	0.017681\\
3.6278	0.9734	0.017691	0.017691\\
3.6314	0.96277	0.017694	0.017694\\
3.6357	0.95213	0.017687	0.017687\\
3.6392	0.94149	0.017664	0.017664\\
3.6437	0.93085	0.017454	0.017454\\
3.6471	0.92021	0.017419	0.017419\\
3.6516	0.90957	0.01749	0.01749\\
3.6566	0.89894	0.017525	0.017525\\
3.6607	0.8883	0.017461	0.017461\\
3.6653	0.87766	0.017654	0.017654\\
3.6698	0.86702	0.017589	0.017589\\
3.674	0.85638	0.017581	0.017581\\
3.6787	0.84574	0.017479	0.017479\\
3.6822	0.83511	0.017412	0.017412\\
3.6863	0.82447	0.017412	0.017412\\
3.6915	0.81383	0.017531	0.017531\\
3.6972	0.80319	0.017589	0.017589\\
3.7027	0.79255	0.017679	0.017679\\
3.7086	0.78191	0.017648	0.017648\\
3.7141	0.77128	0.017578	0.017578\\
3.7198	0.76064	0.017546	0.017546\\
3.7264	0.75	0.017439	0.017439\\
3.7327	0.73936	0.017293	0.017293\\
3.7381	0.72872	0.017183	0.017183\\
3.7433	0.71809	0.017131	0.017131\\
3.7488	0.70745	0.017006	0.017006\\
3.7555	0.69681	0.016872	0.016872\\
3.7601	0.68617	0.016837	0.016837\\
3.766	0.67553	0.016584	0.016584\\
3.7711	0.66489	0.016694	0.016694\\
3.7773	0.65426	0.016691	0.016691\\
3.7848	0.64362	0.016725	0.016725\\
3.7921	0.63298	0.016643	0.016643\\
3.7981	0.62234	0.016667	0.016667\\
3.8059	0.6117	0.016673	0.016673\\
3.813	0.60106	0.016683	0.016683\\
3.8201	0.59043	0.016721	0.016721\\
3.8259	0.57979	0.016512	0.016512\\
3.8325	0.56915	0.016443	0.016443\\
3.8406	0.55851	0.016483	0.016483\\
3.848	0.54787	0.016418	0.016418\\
3.8549	0.53723	0.016679	0.016679\\
3.8639	0.5266	0.016817	0.016817\\
3.8714	0.51596	0.016747	0.016747\\
3.8788	0.50532	0.016814	0.016814\\
3.8859	0.49468	0.016787	0.016787\\
3.8938	0.48404	0.016676	0.016676\\
3.9011	0.4734	0.016795	0.016795\\
3.9113	0.46277	0.016812	0.016812\\
3.9198	0.45213	0.016652	0.016652\\
3.9301	0.44149	0.016576	0.016576\\
3.9386	0.43085	0.016963	0.016963\\
3.9481	0.42021	0.017223	0.017223\\
3.9586	0.40957	0.01699	0.01699\\
3.968	0.39894	0.016786	0.016786\\
3.9784	0.3883	0.01685	0.01685\\
3.987	0.37766	0.016722	0.016722\\
3.9966	0.36702	0.01674	0.01674\\
4.0062	0.35638	0.01698	0.01698\\
4.0169	0.34574	0.016764	0.016764\\
4.0263	0.33511	0.016704	0.016704\\
4.0359	0.32447	0.016868	0.016868\\
4.0484	0.31383	0.017537	0.017537\\
4.0607	0.30319	0.017639	0.017639\\
4.0709	0.29255	0.01799	0.01799\\
4.085	0.28191	0.018081	0.018081\\
4.098	0.27128	0.018429	0.018429\\
4.1118	0.26064	0.018371	0.018371\\
4.126	0.25	0.018434	0.018434\\
4.1402	0.23936	0.018205	0.018205\\
4.1546	0.22872	0.018244	0.018244\\
4.1704	0.21809	0.017889	0.017889\\
4.1874	0.20745	0.018223	0.018223\\
4.2028	0.19681	0.018367	0.018367\\
4.2165	0.18617	0.017727	0.017727\\
4.236	0.17553	0.017623	0.017623\\
4.2531	0.16489	0.017565	0.017565\\
4.27	0.15426	0.017665	0.017665\\
4.2881	0.14362	0.018058	0.018058\\
4.304	0.13298	0.017936	0.017936\\
4.3264	0.12234	0.017809	0.017809\\
4.3455	0.1117	0.018126	0.018126\\
4.3679	0.10106	0.018748	0.018748\\
4.395	0.090426	0.019538	0.019538\\
4.4211	0.079787	0.019795	0.019795\\
4.4498	0.069149	0.020804	0.020804\\
4.4855	0.058511	0.021834	0.021834\\
4.527	0.047872	0.020927	0.020927\\
4.5704	0.037234	0.021241	0.021241\\
4.6275	0.026596	0.021802	0.021802\\
4.6955	0.015957	0.024085	0.024085\\
4.8084	0.0053191	0.02809	0.02809\\
};
\addplot [color=mycolor1]
  table[row sep=crcr]{%
3.6288	0.8898\\
3.6458	0.86612\\
3.6628	0.84072\\
3.6798	0.81335\\
3.6968	0.78546\\
3.7137	0.75637\\
3.7307	0.72596\\
3.7647	0.66501\\
3.7987	0.60537\\
3.8326	0.54862\\
4.2402	0.16363\\
4.6478	0.049018\\
5.0554	0.014749\\
5.3101	0.0069776\\
};
\addlegendentry{Adatt. Pareto appr.}


\addplot[area legend, draw=none, fill=mycolor1, fill opacity=0.15, forget plot]
table[row sep=crcr] {%
x	y\\
3.6288	0.91426\\
3.6458	0.89303\\
3.6628	0.86973\\
3.6798	0.84403\\
3.6968	0.81757\\
3.7137	0.78945\\
3.7307	0.7594\\
3.7477	0.72865\\
3.7647	0.69822\\
3.7817	0.66777\\
3.7987	0.63708\\
3.8156	0.60753\\
3.8326	0.57785\\
3.8496	0.54932\\
3.8666	0.5222\\
3.8836	0.49643\\
3.9006	0.47194\\
3.9175	0.44866\\
3.9345	0.42653\\
3.9515	0.4055\\
3.9685	0.38552\\
3.9855	0.36652\\
4.0025	0.34846\\
4.0194	0.3313\\
4.0364	0.31498\\
4.0534	0.29948\\
4.0704	0.28474\\
4.0874	0.27073\\
4.1043	0.25741\\
4.1213	0.24475\\
4.1383	0.23271\\
4.1553	0.22127\\
4.1723	0.2104\\
4.1893	0.20006\\
4.2062	0.19023\\
4.2232	0.18089\\
4.2402	0.17201\\
4.2572	0.16357\\
4.2742	0.15554\\
4.2912	0.14791\\
4.3081	0.14065\\
4.3251	0.13376\\
4.3421	0.1272\\
4.3591	0.12097\\
4.3761	0.11504\\
4.3931	0.1094\\
4.41	0.10404\\
4.427	0.09895\\
4.444	0.094106\\
4.461	0.0895\\
4.478	0.085121\\
4.495	0.080958\\
4.5119	0.076998\\
4.5289	0.073234\\
4.5459	0.069654\\
4.5629	0.06625\\
4.5799	0.063013\\
4.5969	0.059936\\
4.6138	0.057009\\
4.6308	0.054226\\
4.6478	0.051579\\
4.6648	0.049062\\
4.6818	0.046668\\
4.6988	0.044392\\
4.7157	0.042228\\
4.7327	0.040169\\
4.7497	0.038211\\
4.7667	0.036349\\
4.7837	0.034578\\
4.8006	0.032894\\
4.8176	0.031293\\
4.8346	0.029769\\
4.8516	0.02832\\
4.8686	0.026942\\
4.8856	0.025631\\
4.9025	0.024385\\
4.9195	0.023199\\
4.9365	0.022071\\
4.9535	0.020999\\
4.9705	0.019978\\
4.9875	0.019008\\
5.0044	0.018084\\
5.0214	0.017206\\
5.0384	0.016371\\
5.0554	0.015576\\
5.0724	0.01482\\
5.0894	0.014101\\
5.1063	0.013417\\
5.1233	0.012767\\
5.1403	0.012148\\
5.1573	0.011559\\
5.1743	0.010999\\
5.1913	0.010466\\
5.2082	0.0099586\\
5.2252	0.0094763\\
5.2422	0.0090174\\
5.2592	0.0085808\\
5.2762	0.0081655\\
5.2932	0.0077703\\
5.3101	0.0073944\\
5.3101	0.0065609\\
5.2932	0.0068982\\
5.2762	0.0072528\\
5.2592	0.0076257\\
5.2422	0.0080178\\
5.2252	0.0084301\\
5.2082	0.0088636\\
5.1913	0.0093195\\
5.1743	0.0097988\\
5.1573	0.010303\\
5.1403	0.010833\\
5.1233	0.01139\\
5.1063	0.011976\\
5.0894	0.012592\\
5.0724	0.01324\\
5.0554	0.013921\\
5.0384	0.014638\\
5.0214	0.015391\\
5.0044	0.016183\\
4.9875	0.017016\\
4.9705	0.017892\\
4.9535	0.018813\\
4.9365	0.019782\\
4.9195	0.0208\\
4.9025	0.021871\\
4.8856	0.022997\\
4.8686	0.024181\\
4.8516	0.025426\\
4.8346	0.026736\\
4.8176	0.028113\\
4.8006	0.029561\\
4.7837	0.031083\\
4.7667	0.032684\\
4.7497	0.034368\\
4.7327	0.036138\\
4.7157	0.038\\
4.6988	0.039958\\
4.6818	0.042016\\
4.6648	0.044181\\
4.6478	0.046458\\
4.6308	0.048852\\
4.6138	0.051369\\
4.5969	0.054016\\
4.5799	0.0568\\
4.5629	0.059727\\
4.5459	0.062806\\
4.5289	0.066043\\
4.5119	0.069447\\
4.495	0.073027\\
4.478	0.076791\\
4.461	0.08075\\
4.444	0.084913\\
4.427	0.08929\\
4.41	0.093894\\
4.3931	0.098735\\
4.3761	0.10383\\
4.3591	0.10918\\
4.3421	0.11481\\
4.3251	0.12073\\
4.3081	0.12696\\
4.2912	0.1335\\
4.2742	0.14039\\
4.2572	0.14763\\
4.2402	0.15524\\
4.2232	0.16325\\
4.2062	0.17167\\
4.1893	0.18053\\
4.1723	0.18984\\
4.1553	0.19964\\
4.1383	0.20993\\
4.1213	0.22077\\
4.1043	0.23216\\
4.0874	0.24414\\
4.0704	0.25673\\
4.0534	0.26998\\
4.0364	0.28391\\
4.0194	0.29856\\
4.0025	0.31397\\
3.9855	0.33018\\
3.9685	0.34722\\
3.9515	0.36514\\
3.9345	0.38399\\
3.9175	0.40381\\
3.9006	0.42466\\
3.8836	0.44658\\
3.8666	0.46964\\
3.8496	0.49388\\
3.8326	0.51939\\
3.8156	0.54611\\
3.7987	0.57367\\
3.7817	0.60243\\
3.7647	0.6318\\
3.7477	0.66181\\
3.7307	0.69251\\
3.7137	0.72329\\
3.6968	0.75335\\
3.6798	0.78267\\
3.6628	0.8117\\
3.6458	0.83922\\
3.6288	0.86533\\
}--cycle;
\addplot[only marks, mark=*, mark size=1.1180pt, color=black, fill=black, opacity=0.60, draw=none, mark options={draw=none,line width=0pt}] table[row sep=crcr]{%
x	y\\
3.5623	0.99468\\
3.5723	0.98404\\
3.573	0.9734\\
3.5753	0.96277\\
3.5766	0.95213\\
3.5781	0.94149\\
3.5826	0.93085\\
3.5851	0.92021\\
3.5911	0.90957\\
3.5942	0.89894\\
3.596	0.8883\\
3.6016	0.87766\\
3.603	0.86702\\
3.6047	0.85638\\
3.6086	0.84574\\
3.6086	0.83511\\
3.6097	0.82447\\
3.6142	0.81383\\
3.6261	0.80319\\
3.6383	0.79255\\
3.6464	0.78191\\
3.6535	0.77128\\
3.657	0.76064\\
3.6658	0.75\\
3.6678	0.73936\\
3.6687	0.72872\\
3.6702	0.71809\\
3.6723	0.70745\\
3.6793	0.69681\\
3.6867	0.68617\\
3.6949	0.67553\\
3.7015	0.66489\\
3.7031	0.65426\\
3.719	0.64362\\
3.7213	0.63298\\
3.7265	0.62234\\
3.7307	0.6117\\
3.737	0.60106\\
3.7374	0.59043\\
3.7423	0.57979\\
3.7487	0.56915\\
3.7677	0.55851\\
3.7685	0.54787\\
3.7768	0.53723\\
3.8005	0.5266\\
3.8015	0.51596\\
3.8022	0.50532\\
3.8032	0.49468\\
3.8039	0.48404\\
3.8108	0.4734\\
3.8214	0.46277\\
3.8405	0.45213\\
3.8605	0.44149\\
3.863	0.43085\\
3.8645	0.42021\\
3.8662	0.40957\\
3.8806	0.39894\\
3.8878	0.3883\\
3.8964	0.37766\\
3.905	0.36702\\
3.9294	0.35638\\
3.9303	0.34574\\
3.954	0.33511\\
3.9604	0.32447\\
3.9669	0.31383\\
3.9747	0.30319\\
3.9822	0.29255\\
3.9929	0.28191\\
4.0051	0.27128\\
4.0144	0.26064\\
4.0161	0.25\\
4.0433	0.23936\\
4.0578	0.22872\\
4.073	0.21809\\
4.0857	0.20745\\
4.0904	0.19681\\
4.1168	0.18617\\
4.1265	0.17553\\
4.127	0.16489\\
4.143	0.15426\\
4.1756	0.14362\\
4.2156	0.13298\\
4.3116	0.12234\\
4.3276	0.1117\\
4.3662	0.10106\\
4.4791	0.090426\\
4.4912	0.079787\\
4.517	0.069149\\
4.5743	0.058511\\
4.5914	0.047872\\
4.6138	0.037234\\
4.7898	0.026596\\
5.0876	0.015957\\
5.3101	0.0053191\\
};
\addlegendentry{FRC empirica reale}

\addplot [color=black]
  table[row sep=crcr]{%
3.5608	1\\
5.3101	0.005939\\
};
\addlegendentry{Adatt. Pareto reale}

            % \input{../../../.tex/20/KS/SizeDistributionFittingsfig9.tex}
    
            \nextgroupplot[%
                % xmin=3.4,xmax=5.4,
                % ymin=0.0053191,ymax=1,
            ] % This file was created by matlab2tikz.
%
\definecolor{mycolor1}{rgb}{0.83529,0.36863,0.00000}%
\definecolor{mycolor2}{rgb}{0.00000,0.44700,0.74100}%
\definecolor{mycolor3}{rgb}{0.49400,0.18400,0.55600}%
%
\addplot[only marks, mark=*, mark size=1.1180pt, color=mycolor1, fill=mycolor1, opacity=0.60, draw=none, mark options={draw=none,line width=0pt}] table[row sep=crcr]{%
x	y\\
3.6228	0.99468\\
3.6263	0.98404\\
3.6303	0.9734\\
3.6345	0.96277\\
3.6384	0.95213\\
3.6422	0.94149\\
3.6459	0.93085\\
3.6502	0.92021\\
3.6545	0.90957\\
3.6592	0.89894\\
3.6635	0.8883\\
3.6679	0.87766\\
3.6725	0.86702\\
3.677	0.85638\\
3.6818	0.84574\\
3.686	0.83511\\
3.6907	0.82447\\
3.695	0.81383\\
3.6993	0.80319\\
3.7045	0.79255\\
3.7103	0.78191\\
3.7159	0.77128\\
3.7213	0.76064\\
3.7259	0.75\\
3.7315	0.73936\\
3.7376	0.72872\\
3.7436	0.71809\\
3.749	0.70745\\
3.7549	0.69681\\
3.7613	0.68617\\
3.769	0.67553\\
3.775	0.66489\\
3.7822	0.65426\\
3.789	0.64362\\
3.7959	0.63298\\
3.8022	0.62234\\
3.8105	0.6117\\
3.8171	0.60106\\
3.8247	0.59043\\
3.8324	0.57979\\
3.8391	0.56915\\
3.8469	0.55851\\
3.8533	0.54787\\
3.8606	0.53723\\
3.8689	0.5266\\
3.8766	0.51596\\
3.8862	0.50532\\
3.8938	0.49468\\
3.9018	0.48404\\
3.9107	0.4734\\
3.9196	0.46277\\
3.9288	0.45213\\
3.936	0.44149\\
3.9448	0.43085\\
3.9549	0.42021\\
3.9638	0.40957\\
3.9737	0.39894\\
3.9837	0.3883\\
3.9943	0.37766\\
4.003	0.36702\\
4.0139	0.35638\\
4.0231	0.34574\\
4.036	0.33511\\
4.0469	0.32447\\
4.0587	0.31383\\
4.073	0.30319\\
4.0862	0.29255\\
4.0998	0.28191\\
4.1134	0.27128\\
4.1263	0.26064\\
4.1389	0.25\\
4.1536	0.23936\\
4.1667	0.22872\\
4.1809	0.21809\\
4.1942	0.20745\\
4.2098	0.19681\\
4.2282	0.18617\\
4.2451	0.17553\\
4.2634	0.16489\\
4.2839	0.15426\\
4.3018	0.14362\\
4.319	0.13298\\
4.339	0.12234\\
4.3591	0.1117\\
4.3807	0.10106\\
4.3999	0.090426\\
4.4271	0.079787\\
4.454	0.069149\\
4.4834	0.058511\\
4.5198	0.047872\\
4.5651	0.037234\\
4.6107	0.026596\\
4.6816	0.015957\\
4.7722	0.0053191\\
};
\addlegendentry{FRC empirica appr.}

\addplot [color=mycolor1, only marks, every error bar/.append style={opacity=0.30}, mark=*, mark size=0pt, draw=none, forget plot]
 plot [error bars/.cd, x dir=both, x explicit, error bar style={line width=1pt, color=mycolor1}, error mark options={mark=none,mark size=0pt}]
 table[row sep=crcr, x error plus index=2, x error minus index=3]{%
3.6228	0.99468	0.0030118	0.0030118\\
3.6263	0.98404	0.0030221	0.0030221\\
3.6303	0.9734	0.003048	0.003048\\
3.6345	0.96277	0.0029435	0.0029435\\
3.6384	0.95213	0.0029654	0.0029654\\
3.6422	0.94149	0.0029683	0.0029683\\
3.6459	0.93085	0.0031145	0.0031145\\
3.6502	0.92021	0.0030723	0.0030723\\
3.6545	0.90957	0.0031095	0.0031095\\
3.6592	0.89894	0.0032137	0.0032137\\
3.6635	0.8883	0.0033616	0.0033616\\
3.6679	0.87766	0.0033587	0.0033587\\
3.6725	0.86702	0.003436	0.003436\\
3.677	0.85638	0.0033124	0.0033124\\
3.6818	0.84574	0.0034363	0.0034363\\
3.686	0.83511	0.0035117	0.0035117\\
3.6907	0.82447	0.0037475	0.0037475\\
3.695	0.81383	0.0038168	0.0038168\\
3.6993	0.80319	0.0038794	0.0038794\\
3.7045	0.79255	0.0039683	0.0039683\\
3.7103	0.78191	0.0039696	0.0039696\\
3.7159	0.77128	0.0040093	0.0040093\\
3.7213	0.76064	0.0040001	0.0040001\\
3.7259	0.75	0.0040127	0.0040127\\
3.7315	0.73936	0.0040247	0.0040247\\
3.7376	0.72872	0.0039412	0.0039412\\
3.7436	0.71809	0.0039168	0.0039168\\
3.749	0.70745	0.004012	0.004012\\
3.7549	0.69681	0.0040165	0.0040165\\
3.7613	0.68617	0.0040394	0.0040394\\
3.769	0.67553	0.0042063	0.0042063\\
3.775	0.66489	0.0041416	0.0041416\\
3.7822	0.65426	0.004108	0.004108\\
3.789	0.64362	0.0042736	0.0042736\\
3.7959	0.63298	0.0043112	0.0043112\\
3.8022	0.62234	0.0043687	0.0043687\\
3.8105	0.6117	0.0043986	0.0043986\\
3.8171	0.60106	0.0044699	0.0044699\\
3.8247	0.59043	0.0045414	0.0045414\\
3.8324	0.57979	0.0045054	0.0045054\\
3.8391	0.56915	0.0045316	0.0045316\\
3.8469	0.55851	0.004509	0.004509\\
3.8533	0.54787	0.0045624	0.0045624\\
3.8606	0.53723	0.0044381	0.0044381\\
3.8689	0.5266	0.0045637	0.0045637\\
3.8766	0.51596	0.0045172	0.0045172\\
3.8862	0.50532	0.0046976	0.0046976\\
3.8938	0.49468	0.0044914	0.0044914\\
3.9018	0.48404	0.0045069	0.0045069\\
3.9107	0.4734	0.0045815	0.0045815\\
3.9196	0.46277	0.0045831	0.0045831\\
3.9288	0.45213	0.0046655	0.0046655\\
3.936	0.44149	0.0048751	0.0048751\\
3.9448	0.43085	0.0050417	0.0050417\\
3.9549	0.42021	0.0050819	0.0050819\\
3.9638	0.40957	0.0052092	0.0052092\\
3.9737	0.39894	0.0051505	0.0051505\\
3.9837	0.3883	0.0051306	0.0051306\\
3.9943	0.37766	0.0051166	0.0051166\\
4.003	0.36702	0.0051951	0.0051951\\
4.0139	0.35638	0.0050619	0.0050619\\
4.0231	0.34574	0.0051686	0.0051686\\
4.036	0.33511	0.0049584	0.0049584\\
4.0469	0.32447	0.0049068	0.0049068\\
4.0587	0.31383	0.0050858	0.0050858\\
4.073	0.30319	0.0051435	0.0051435\\
4.0862	0.29255	0.0050315	0.0050315\\
4.0998	0.28191	0.0049629	0.0049629\\
4.1134	0.27128	0.0051363	0.0051363\\
4.1263	0.26064	0.0051984	0.0051984\\
4.1389	0.25	0.0051243	0.0051243\\
4.1536	0.23936	0.005389	0.005389\\
4.1667	0.22872	0.0052368	0.0052368\\
4.1809	0.21809	0.0053127	0.0053127\\
4.1942	0.20745	0.0056331	0.0056331\\
4.2098	0.19681	0.0055258	0.0055258\\
4.2282	0.18617	0.0052646	0.0052646\\
4.2451	0.17553	0.0053001	0.0053001\\
4.2634	0.16489	0.0055714	0.0055714\\
4.2839	0.15426	0.0058965	0.0058965\\
4.3018	0.14362	0.0059847	0.0059847\\
4.319	0.13298	0.0056011	0.0056011\\
4.339	0.12234	0.0059027	0.0059027\\
4.3591	0.1117	0.0056182	0.0056182\\
4.3807	0.10106	0.0060889	0.0060889\\
4.3999	0.090426	0.0057706	0.0057706\\
4.4271	0.079787	0.0065299	0.0065299\\
4.454	0.069149	0.006593	0.006593\\
4.4834	0.058511	0.0065854	0.0065854\\
4.5198	0.047872	0.0071316	0.0071316\\
4.5651	0.037234	0.007653	0.007653\\
4.6107	0.026596	0.0087925	0.0087925\\
4.6816	0.015957	0.0096303	0.0096303\\
4.7722	0.0053191	0.0093005	0.0093005\\
};
\addplot [color=mycolor1]
  table[row sep=crcr]{%
3.6208	0.98323\\
3.6379	0.9491\\
3.6549	0.90581\\
4.3375	0.12129\\
4.8835	0.024401\\
5.3101	0.0069926\\
};
\addlegendentry{Adatt. Pareto appr.}


\addplot[area legend, draw=none, fill=mycolor1, fill opacity=0.15, forget plot]
table[row sep=crcr] {%
x	y\\
3.6208	0.9875\\
3.6379	0.95578\\
3.6549	0.91301\\
3.672	0.86791\\
3.6891	0.82499\\
3.7061	0.78421\\
3.7232	0.74545\\
3.7402	0.70862\\
3.7573	0.67362\\
3.7744	0.64035\\
3.7914	0.60875\\
3.8085	0.57871\\
3.8256	0.55017\\
3.8426	0.52306\\
3.8597	0.49729\\
3.8768	0.47282\\
3.8938	0.44956\\
3.9109	0.42747\\
3.928	0.40649\\
3.945	0.38655\\
3.9621	0.36761\\
3.9791	0.34962\\
3.9962	0.33252\\
4.0133	0.31627\\
4.0303	0.30083\\
4.0474	0.28615\\
4.0645	0.27219\\
4.0815	0.25892\\
4.0986	0.24631\\
4.1157	0.23431\\
4.1327	0.2229\\
4.1498	0.21205\\
4.1668	0.20173\\
4.1839	0.19191\\
4.201	0.18258\\
4.218	0.1737\\
4.2351	0.16525\\
4.2522	0.15722\\
4.2692	0.14958\\
4.2863	0.14231\\
4.3034	0.13539\\
4.3204	0.12881\\
4.3375	0.12255\\
4.3546	0.1166\\
4.3716	0.11094\\
4.3887	0.10555\\
4.4057	0.10042\\
4.4228	0.095547\\
4.4399	0.090908\\
4.4569	0.086494\\
4.474	0.082296\\
4.4911	0.078301\\
4.5081	0.074501\\
4.5252	0.070885\\
4.5423	0.067445\\
4.5593	0.064173\\
4.5764	0.061059\\
4.5934	0.058097\\
4.6105	0.055279\\
4.6276	0.052598\\
4.6446	0.050047\\
4.6617	0.047619\\
4.6788	0.04531\\
4.6958	0.043113\\
4.7129	0.041023\\
4.73	0.039034\\
4.747	0.037142\\
4.7641	0.035341\\
4.7812	0.033628\\
4.7982	0.031999\\
4.8153	0.030448\\
4.8323	0.028972\\
4.8494	0.027569\\
4.8665	0.026233\\
4.8835	0.024962\\
4.9006	0.023753\\
4.9177	0.022602\\
4.9347	0.021507\\
4.9518	0.020466\\
4.9689	0.019475\\
4.9859	0.018532\\
5.003	0.017634\\
5.02	0.01678\\
5.0371	0.015968\\
5.0542	0.015195\\
5.0712	0.014459\\
5.0883	0.01376\\
5.1054	0.013094\\
5.1224	0.01246\\
5.1395	0.011857\\
5.1566	0.011283\\
5.1736	0.010737\\
5.1907	0.010218\\
5.2078	0.0097235\\
5.2248	0.0092532\\
5.2419	0.0088056\\
5.2589	0.0083797\\
5.276	0.0079744\\
5.2931	0.0075888\\
5.3101	0.0072218\\
5.3101	0.0067633\\
5.2931	0.0071125\\
5.276	0.0074798\\
5.2589	0.0078661\\
5.2419	0.0082723\\
5.2248	0.0086996\\
5.2078	0.009149\\
5.1907	0.0096216\\
5.1736	0.010119\\
5.1566	0.010642\\
5.1395	0.011191\\
5.1224	0.01177\\
5.1054	0.012378\\
5.0883	0.013018\\
5.0712	0.013691\\
5.0542	0.014399\\
5.0371	0.015143\\
5.02	0.015926\\
5.003	0.016749\\
4.9859	0.017616\\
4.9689	0.018527\\
4.9518	0.019485\\
4.9347	0.020493\\
4.9177	0.021553\\
4.9006	0.022668\\
4.8835	0.023841\\
4.8665	0.025074\\
4.8494	0.026372\\
4.8323	0.027737\\
4.8153	0.029172\\
4.7982	0.030682\\
4.7812	0.03227\\
4.7641	0.033941\\
4.747	0.035698\\
4.73	0.037546\\
4.7129	0.039491\\
4.6958	0.041536\\
4.6788	0.043687\\
4.6617	0.04595\\
4.6446	0.04833\\
4.6276	0.050834\\
4.6105	0.053467\\
4.5934	0.056237\\
4.5764	0.059151\\
4.5593	0.062216\\
4.5423	0.065441\\
4.5252	0.068832\\
4.5081	0.0724\\
4.4911	0.076153\\
4.474	0.0801\\
4.4569	0.084253\\
4.4399	0.088621\\
4.4228	0.093216\\
4.4057	0.09805\\
4.3887	0.10313\\
4.3716	0.10848\\
4.3546	0.11411\\
4.3375	0.12003\\
4.3204	0.12625\\
4.3034	0.1328\\
4.2863	0.13969\\
4.2692	0.14694\\
4.2522	0.15456\\
4.2351	0.16257\\
4.218	0.17101\\
4.201	0.17988\\
4.1839	0.18921\\
4.1668	0.19902\\
4.1498	0.20935\\
4.1327	0.2202\\
4.1157	0.23162\\
4.0986	0.24363\\
4.0815	0.25626\\
4.0645	0.26954\\
4.0474	0.2835\\
4.0303	0.29818\\
4.0133	0.31361\\
3.9962	0.32983\\
3.9791	0.34688\\
3.9621	0.36479\\
3.945	0.38362\\
3.928	0.4034\\
3.9109	0.42418\\
3.8938	0.44602\\
3.8768	0.46896\\
3.8597	0.49306\\
3.8426	0.5184\\
3.8256	0.54501\\
3.8085	0.57299\\
3.7914	0.60239\\
3.7744	0.63328\\
3.7573	0.66576\\
3.7402	0.69989\\
3.7232	0.73577\\
3.7061	0.77349\\
3.6891	0.81313\\
3.672	0.85481\\
3.6549	0.8986\\
3.6379	0.94243\\
3.6208	0.97897\\
}--cycle;
\addplot[only marks, mark=*, mark size=1.1180pt, color=black, fill=black, opacity=0.60, draw=none, mark options={draw=none,line width=0pt}] table[row sep=crcr]{%
x	y\\
3.5623	0.99468\\
3.5723	0.98404\\
3.573	0.9734\\
3.5753	0.96277\\
3.5766	0.95213\\
3.5781	0.94149\\
3.5826	0.93085\\
3.5851	0.92021\\
3.5911	0.90957\\
3.5942	0.89894\\
3.596	0.8883\\
3.6016	0.87766\\
3.603	0.86702\\
3.6047	0.85638\\
3.6086	0.84574\\
3.6086	0.83511\\
3.6097	0.82447\\
3.6142	0.81383\\
3.6261	0.80319\\
3.6383	0.79255\\
3.6464	0.78191\\
3.6535	0.77128\\
3.657	0.76064\\
3.6658	0.75\\
3.6678	0.73936\\
3.6687	0.72872\\
3.6702	0.71809\\
3.6723	0.70745\\
3.6793	0.69681\\
3.6867	0.68617\\
3.6949	0.67553\\
3.7015	0.66489\\
3.7031	0.65426\\
3.719	0.64362\\
3.7213	0.63298\\
3.7265	0.62234\\
3.7307	0.6117\\
3.737	0.60106\\
3.7374	0.59043\\
3.7423	0.57979\\
3.7487	0.56915\\
3.7677	0.55851\\
3.7685	0.54787\\
3.7768	0.53723\\
3.8005	0.5266\\
3.8015	0.51596\\
3.8022	0.50532\\
3.8032	0.49468\\
3.8039	0.48404\\
3.8108	0.4734\\
3.8214	0.46277\\
3.8405	0.45213\\
3.8605	0.44149\\
3.863	0.43085\\
3.8645	0.42021\\
3.8662	0.40957\\
3.8806	0.39894\\
3.8878	0.3883\\
3.8964	0.37766\\
3.905	0.36702\\
3.9294	0.35638\\
3.9303	0.34574\\
3.954	0.33511\\
3.9604	0.32447\\
3.9669	0.31383\\
3.9747	0.30319\\
3.9822	0.29255\\
3.9929	0.28191\\
4.0051	0.27128\\
4.0144	0.26064\\
4.0161	0.25\\
4.0433	0.23936\\
4.0578	0.22872\\
4.073	0.21809\\
4.0857	0.20745\\
4.0904	0.19681\\
4.1168	0.18617\\
4.1265	0.17553\\
4.127	0.16489\\
4.143	0.15426\\
4.1756	0.14362\\
4.2156	0.13298\\
4.3116	0.12234\\
4.3276	0.1117\\
4.3662	0.10106\\
4.4791	0.090426\\
4.4912	0.079787\\
4.517	0.069149\\
4.5743	0.058511\\
4.5914	0.047872\\
4.6138	0.037234\\
4.7898	0.026596\\
5.0876	0.015957\\
5.3101	0.0053191\\
};
\addlegendentry{FRC empirica reale}

\addplot [color=black]
  table[row sep=crcr]{%
3.5608	1\\
5.3101	0.005939\\
};
\addlegendentry{Adatt. Pareto reale}

            % \input{../../../.tex/20/KS/SizeDistributionFittingsfig9.tex}
        \end{groupplot}
    
        \begin{groupplot}[
            group style={
                group name=Row5,
                plotLastGroup,
            },
            plotRow,
            row5KeysLin,
            plotLegendNW,
            %
            xmin=0,xmax=10000,
            ymin=4150,ymax=4650,
        ]
            \nextgroupplot[%
                at={(Row4 c1r1.south east)},yshift=-\yGroupSep em,anchor=north east,
                % xmin=0,xmax=10000,
                % ymin=4150,ymax=4650,
            ] % This file was created by matlab2tikz.
%
\definecolor{mycolor1}{rgb}{0.00000,0.44706,0.69804}%
\definecolor{mycolor2}{rgb}{0.83529,0.36863,0.00000}%
%
\addplot [color=mycolor1]
  table[row sep=crcr]{%
0	4395.3\\
200	4396.6\\
600	4383\\
800	4370.2\\
1000	4369.3\\
1200	4369.7\\
1400	4365.3\\
1600	4354.6\\
1800	4354.6\\
2000	4356.4\\
2200	4351.3\\
2400	4350.3\\
2600	4358.9\\
2800	4359.4\\
3000	4355.1\\
3200	4330.5\\
3400	4318.3\\
3600	4315\\
3800	4322\\
4000	4315.8\\
4200	4327\\
4400	4323.6\\
4600	4333.7\\
4800	4345.2\\
5000	4350.5\\
5200	4370.2\\
5400	4390.3\\
5600	4399.9\\
5800	4394.4\\
6000	4384.5\\
6200	4374.4\\
6400	4363.7\\
6600	4337.3\\
6800	4325.5\\
7000	4322.2\\
7200	4340.8\\
7400	4345.4\\
7600	4349.4\\
7800	4357.9\\
8000	4359.8\\
8200	4376.8\\
8400	4376.9\\
8600	4374.3\\
8800	4380.5\\
9000	4392.4\\
9200	4385.5\\
9400	4385.6\\
9600	4382.2\\
9800	4387.4\\
10000	4392.1\\
};
\addlegendentry{Taglia media esatta}


\addplot[area legend, draw=none, fill=mycolor1, fill opacity=0.15, forget plot]
table[row sep=crcr] {%
x	y\\
0	4395.3\\
200	4401.7\\
400	4403.1\\
600	4401.7\\
800	4395.9\\
1000	4400.9\\
1200	4411.5\\
1400	4413.2\\
1600	4406.5\\
1800	4410.9\\
2000	4417.4\\
2200	4415.5\\
2400	4416.4\\
2600	4429.2\\
2800	4434.3\\
3000	4435.4\\
3200	4415\\
3400	4405.7\\
3600	4407\\
3800	4419.8\\
4000	4412.9\\
4200	4426.4\\
4400	4425.7\\
4600	4439.7\\
4800	4457.2\\
5000	4462.7\\
5200	4485.4\\
5400	4512.2\\
5600	4526.2\\
5800	4522.4\\
6000	4516.9\\
6200	4508.8\\
6400	4500.8\\
6600	4473.1\\
6800	4461.1\\
7000	4461.5\\
7200	4483.9\\
7400	4496.3\\
7600	4506.9\\
7800	4520\\
8000	4524\\
8200	4544.1\\
8400	4543.6\\
8600	4544.6\\
8800	4554.7\\
9000	4566.4\\
9200	4561.9\\
9400	4564.7\\
9600	4563.5\\
9800	4570.9\\
10000	4575.6\\
10000	4208.5\\
9800	4203.9\\
9600	4200.9\\
9400	4206.6\\
9200	4209.1\\
9000	4218.3\\
8800	4206.2\\
8600	4204.1\\
8400	4210.3\\
8200	4209.4\\
8000	4195.7\\
7800	4195.7\\
7600	4191.9\\
7400	4194.6\\
7200	4197.8\\
7000	4182.8\\
6800	4189.9\\
6600	4201.4\\
6400	4226.5\\
6200	4240.1\\
6000	4252.2\\
5800	4266.4\\
5600	4273.6\\
5400	4268.3\\
5200	4254.9\\
5000	4238.2\\
4800	4233.2\\
4600	4227.6\\
4400	4221.5\\
4200	4227.6\\
4000	4218.7\\
3800	4224.2\\
3600	4223\\
3400	4230.9\\
3200	4246\\
3000	4274.8\\
2800	4284.6\\
2600	4288.7\\
2400	4284.2\\
2200	4287\\
2000	4295.4\\
1800	4298.3\\
1600	4302.6\\
1400	4317.4\\
1200	4327.9\\
1000	4337.7\\
800	4344.5\\
600	4364.4\\
400	4376.5\\
200	4391.4\\
0	4395.3\\
}--cycle;
\addplot [color=mycolor2]
  table[row sep=crcr]{%
0	4395.3\\
200	4398.4\\
400	4406.4\\
600	4399.9\\
800	4399.6\\
1000	4409.6\\
1200	4397.2\\
1400	4408.1\\
1600	4422\\
1800	4428.3\\
2000	4411.7\\
2200	4409.6\\
2400	4395.4\\
2600	4389.8\\
2800	4376.4\\
3000	4376\\
3200	4374.3\\
3400	4376.8\\
3600	4385.1\\
3800	4403.8\\
4000	4393.7\\
4200	4391.1\\
4400	4396.3\\
4600	4406.2\\
4800	4417.3\\
5000	4408.2\\
5200	4404.8\\
5400	4404.6\\
5600	4384.3\\
5800	4396.1\\
6000	4403\\
6200	4408\\
6400	4402.7\\
6600	4390.7\\
6800	4388.8\\
7000	4383.1\\
7200	4380.8\\
7400	4381.5\\
7600	4388.7\\
7800	4397.9\\
8000	4398.4\\
8200	4414.3\\
8400	4431.2\\
8600	4420.2\\
8800	4420.6\\
9000	4443.9\\
9200	4457.7\\
9400	4459.4\\
9600	4452.7\\
9800	4462\\
10000	4463.2\\
};
\addlegendentry{Taglia media appr.}


\addplot[area legend, draw=none, fill=mycolor2, fill opacity=0.15, forget plot]
table[row sep=crcr] {%
x	y\\
0	4395.3\\
200	4403.9\\
400	4419.6\\
600	4419.7\\
800	4429.3\\
1000	4447.1\\
1200	4440.3\\
1400	4460.4\\
1600	4482.2\\
1800	4492.7\\
2000	4483.2\\
2200	4486\\
2400	4474.6\\
2600	4471.5\\
2800	4462.4\\
3000	4465.6\\
3200	4464.2\\
3400	4470.9\\
3600	4486.1\\
3800	4513.8\\
4000	4504.8\\
4200	4503.8\\
4400	4511.8\\
4600	4525.6\\
4800	4538.6\\
5000	4533.8\\
5200	4532.1\\
5400	4529.6\\
5600	4509.6\\
5800	4525\\
6000	4533.7\\
6200	4539.6\\
6400	4535.8\\
6600	4524\\
6800	4529.1\\
7000	4523.3\\
7200	4523.3\\
7400	4526.2\\
7600	4536.1\\
7800	4549\\
8000	4549.5\\
8200	4570.6\\
8400	4590.3\\
8600	4579.1\\
8800	4577.7\\
9000	4602.7\\
9200	4620.3\\
9400	4625.4\\
9600	4622.1\\
9800	4634.5\\
10000	4637.9\\
10000	4288.5\\
9800	4289.6\\
9600	4283.4\\
9400	4293.4\\
9200	4295\\
9000	4285.1\\
8800	4263.5\\
8600	4261.2\\
8400	4272.1\\
8200	4258\\
8000	4247.2\\
7800	4246.9\\
7600	4241.2\\
7400	4236.8\\
7200	4238.3\\
7000	4242.9\\
6800	4248.4\\
6600	4257.3\\
6400	4269.6\\
6200	4276.3\\
6000	4272.2\\
5800	4267.1\\
5600	4259.1\\
5400	4279.7\\
5200	4277.6\\
5000	4282.6\\
4800	4296\\
4600	4286.9\\
4400	4280.7\\
4200	4278.4\\
4000	4282.7\\
3800	4293.7\\
3600	4284.1\\
3400	4282.7\\
3200	4284.4\\
3000	4286.5\\
2800	4290.4\\
2600	4308.2\\
2400	4316.2\\
2200	4333.3\\
2000	4340.3\\
1800	4363.9\\
1600	4361.9\\
1400	4355.7\\
1200	4354.2\\
1000	4372\\
800	4369.8\\
600	4380\\
400	4393.2\\
200	4392.9\\
0	4395.3\\
}--cycle;
\addplot [color=black]
  table[row sep=crcr]{%
0	4395.3\\
10000	4395.3\\
};
\addlegendentry{Taglia media reale}

            % \input{../../../.tex/20/KS/AverageSize.tex}
    
            \nextgroupplot[%
                % xmin=0,xmax=10000,
                % ymin=4395.3,ymax=4395.3,
            ] % This file was created by matlab2tikz.
%
\definecolor{mycolor1}{rgb}{0.00000,0.44706,0.69804}%
\definecolor{mycolor2}{rgb}{0.83529,0.36863,0.00000}%
%
\addplot [color=mycolor1]
  table[row sep=crcr]{%
0	4395.3\\
200	4395.3\\
800	4395.3\\
1000	4395.3\\
1200	4395.3\\
1400	4395.3\\
1800	4395.3\\
2000	4395.3\\
2200	4395.3\\
2400	4395.3\\
5400	4395.3\\
5600	4395.3\\
5800	4395.3\\
6200	4395.3\\
6400	4395.3\\
6600	4395.3\\
6800	4395.3\\
7800	4395.3\\
8000	4395.3\\
8400	4395.3\\
8600	4395.3\\
8800	4395.3\\
9000	4395.3\\
9400	4395.3\\
9600	4395.3\\
9800	4395.3\\
10000	4395.3\\
};
\addlegendentry{Taglia media esatta}


\addplot[area legend, draw=none, fill=mycolor1, fill opacity=0.15, forget plot]
table[row sep=crcr] {%
x	y\\
0	4395.3\\
200	4395.3\\
400	4395.3\\
600	4395.3\\
800	4395.3\\
1000	4395.3\\
1200	4395.3\\
1400	4395.3\\
1600	4395.3\\
1800	4395.3\\
2000	4395.3\\
2200	4395.3\\
2400	4395.3\\
2600	4395.3\\
2800	4395.3\\
3000	4395.3\\
3200	4395.3\\
3400	4395.3\\
3600	4395.3\\
3800	4395.3\\
4000	4395.3\\
4200	4395.3\\
4400	4395.3\\
4600	4395.3\\
4800	4395.3\\
5000	4395.3\\
5200	4395.3\\
5400	4395.3\\
5600	4395.3\\
5800	4395.3\\
6000	4395.3\\
6200	4395.3\\
6400	4395.3\\
6600	4395.3\\
6800	4395.3\\
7000	4395.3\\
7200	4395.3\\
7400	4395.3\\
7600	4395.3\\
7800	4395.3\\
8000	4395.3\\
8200	4395.3\\
8400	4395.3\\
8600	4395.3\\
8800	4395.3\\
9000	4395.3\\
9200	4395.3\\
9400	4395.3\\
9600	4395.3\\
9800	4395.3\\
10000	4395.3\\
10000	4395.3\\
9800	4395.3\\
9600	4395.3\\
9400	4395.3\\
9200	4395.3\\
9000	4395.3\\
8800	4395.3\\
8600	4395.3\\
8400	4395.3\\
8200	4395.3\\
8000	4395.3\\
7800	4395.3\\
7600	4395.3\\
7400	4395.3\\
7200	4395.3\\
7000	4395.3\\
6800	4395.3\\
6600	4395.3\\
6400	4395.3\\
6200	4395.3\\
6000	4395.3\\
5800	4395.3\\
5600	4395.3\\
5400	4395.3\\
5200	4395.3\\
5000	4395.3\\
4800	4395.3\\
4600	4395.3\\
4400	4395.3\\
4200	4395.3\\
4000	4395.3\\
3800	4395.3\\
3600	4395.3\\
3400	4395.3\\
3200	4395.3\\
3000	4395.3\\
2800	4395.3\\
2600	4395.3\\
2400	4395.3\\
2200	4395.3\\
2000	4395.3\\
1800	4395.3\\
1600	4395.3\\
1400	4395.3\\
1200	4395.3\\
1000	4395.3\\
800	4395.3\\
600	4395.3\\
400	4395.3\\
200	4395.3\\
0	4395.3\\
}--cycle;
\addplot [color=mycolor2]
  table[row sep=crcr]{%
0	4395.3\\
200	4395.3\\
800	4395.3\\
1000	4395.3\\
1200	4395.3\\
1400	4395.3\\
4600	4395.3\\
4800	4395.3\\
5000	4395.3\\
5200	4395.3\\
5400	4395.3\\
5600	4395.3\\
6400	4395.3\\
6600	4395.3\\
7200	4395.3\\
7400	4395.3\\
7600	4395.3\\
7800	4395.3\\
8200	4395.3\\
8400	4395.3\\
8600	4395.3\\
8800	4395.3\\
9200	4395.3\\
9800	4395.3\\
10000	4395.3\\
};
\addlegendentry{Taglia media appr.}


\addplot[area legend, draw=none, fill=mycolor2, fill opacity=0.15, forget plot]
table[row sep=crcr] {%
x	y\\
0	4395.3\\
200	4395.3\\
400	4395.3\\
600	4395.3\\
800	4395.3\\
1000	4395.3\\
1200	4395.3\\
1400	4395.3\\
1600	4395.3\\
1800	4395.3\\
2000	4395.3\\
2200	4395.3\\
2400	4395.3\\
2600	4395.3\\
2800	4395.3\\
3000	4395.3\\
3200	4395.3\\
3400	4395.3\\
3600	4395.3\\
3800	4395.3\\
4000	4395.3\\
4200	4395.3\\
4400	4395.3\\
4600	4395.3\\
4800	4395.3\\
5000	4395.3\\
5200	4395.3\\
5400	4395.3\\
5600	4395.3\\
5800	4395.3\\
6000	4395.3\\
6200	4395.3\\
6400	4395.3\\
6600	4395.3\\
6800	4395.3\\
7000	4395.3\\
7200	4395.3\\
7400	4395.3\\
7600	4395.3\\
7800	4395.3\\
8000	4395.3\\
8200	4395.3\\
8400	4395.3\\
8600	4395.3\\
8800	4395.3\\
9000	4395.3\\
9200	4395.3\\
9400	4395.3\\
9600	4395.3\\
9800	4395.3\\
10000	4395.3\\
10000	4395.3\\
9800	4395.3\\
9600	4395.3\\
9400	4395.3\\
9200	4395.3\\
9000	4395.3\\
8800	4395.3\\
8600	4395.3\\
8400	4395.3\\
8200	4395.3\\
8000	4395.3\\
7800	4395.3\\
7600	4395.3\\
7400	4395.3\\
7200	4395.3\\
7000	4395.3\\
6800	4395.3\\
6600	4395.3\\
6400	4395.3\\
6200	4395.3\\
6000	4395.3\\
5800	4395.3\\
5600	4395.3\\
5400	4395.3\\
5200	4395.3\\
5000	4395.3\\
4800	4395.3\\
4600	4395.3\\
4400	4395.3\\
4200	4395.3\\
4000	4395.3\\
3800	4395.3\\
3600	4395.3\\
3400	4395.3\\
3200	4395.3\\
3000	4395.3\\
2800	4395.3\\
2600	4395.3\\
2400	4395.3\\
2200	4395.3\\
2000	4395.3\\
1800	4395.3\\
1600	4395.3\\
1400	4395.3\\
1200	4395.3\\
1000	4395.3\\
800	4395.3\\
600	4395.3\\
400	4395.3\\
200	4395.3\\
0	4395.3\\
}--cycle;
\addplot [color=black]
  table[row sep=crcr]{%
0	4395.3\\
10000	4395.3\\
};
\addlegendentry{Taglia media reale}

            % \input{../../../.tex/20/KS/AverageSize.tex}
        \end{groupplot}
    %\endgroup

    %\begingroup Titoli e didascalie
        \node[title] at (Row1 c1r1.north) {Con fluttuazioni \(\gamma\)};
        \node[title] at (Row1 c2r1.north) {Senza fluttuazioni \(\gamma\)};

        \foreach \r in {1,...,5}
            \foreach \c in {1,2}{
                \ifcase\r\relax% Caso 0
                \or\writeSubCaption{Row\r\space c\c r1}{north west}{north west}{0em}{0em}  % Caso 1
                \or\writeSubCaption{Row\r\space c\c r1}{north west}{north west}{0em}{0em}  % Caso 2
                \or\writeSubCaption{Row\r\space c\c r1}{south west}{south west}{0em}{0em}  % Caso 3
                \or\writeSubCaption{Row\r\space c\c r1}{south west}{south west}{0em}{0em}  % Caso 4
                \else\ifnum\c<3
                    \writeSubCaption{Row\r\space c\c r1}{south west}{south west}{0em}{0em} % Caso 5
                \else\relax\fi\fi
                
            } % https://tex.stackexchange.com/questions/632045/get-subcaption-for-every-plot-using-groupplot-pgf-tikz-after-2020-update
    %\endgroup

    %\begingroup Parametri
        %\begingroup beta
            \newcommand\betaExactList[1]{
                \ifnum#1=1\num{1.238+-0.012} % Caso con fluttuazioni
                \else\num{1.227+-0.010}\fi   % Caso senza fluttuazioni
            }
            \newcommand\betaAprxList[1]{
                \ifnum#1=1\num{1.287+-0.015} % Caso con fluttuazioni
                \else\num{1.280+-0.011}\fi   % Caso senza fluttuazioni
            }
        %\endgroup

        %\begingroup xi
        \newcommand\xiExactList[1]{
            \ifnum#1=1\num{0.766+-0.024} % Caso con fluttuazioni
            \else\num{0.786+-0.021}\fi   % Caso senza fluttuazioni
        }
        \newcommand\xiAprxList[1]{
            \ifnum#1=1\num{0.707+-0.025} % Caso con fluttuazioni
            \else\num{0.723+-0.022}\fi   % Caso senza fluttuazioni
        }
        %\endgroup

        %\begingroup mI
        \newcommand\mIExactList[1]{
            \ifnum#1=1\num{3.281+-0.018} % Caso con fluttuazioni
            \else\num{3.302+-0.005}\fi   % Caso senza fluttuazioni
        }
        \newcommand\mIAprxList[1]{
            \ifnum#1=1\num{3.268+-0.017} % Caso con fluttuazioni
            \else\num{3.276+-0.006}\fi   % Caso senza fluttuazioni
        }
        %\endgroup

        %\begingroup sI
        \newcommand\sIExactList[1]{
            \ifnum#1=1\num{0.256+-0.005} % Caso con fluttuazioni
            \else\num{0.241+-0.005}\fi   % Caso senza fluttuazioni
        }
        \newcommand\sIAprxList[1]{
            \ifnum#1=1\num{0.253+-0.006} % Caso con fluttuazioni
            \else\num{0.245+-0.005}\fi   % Caso senza fluttuazioni
        }
        %\endgroup

        %\begingroup mII
        \newcommand\mIIExactList[1]{
            \ifnum#1=1\num{3.926+-0.048} % Caso con fluttuazioni
            \else\num{3.990+-0.042}\fi   % Caso senza fluttuazioni
        }
        \newcommand\mIIAprxList[1]{
            \ifnum#1=1\num{3.836+-0.042} % Caso con fluttuazioni
            \else\num{3.870+-0.037}\fi   % Caso senza fluttuazioni
        }
        %\endgroup

        %\begingroup sII
        \newcommand\sIIExactList[1]{
            \ifnum#1=1\num{0.368+-0.019} % Caso con fluttuazioni
            \else\num{0.331+-0.019}\fi   % Caso senza fluttuazioni
        }
        \newcommand\sIIAprxList[1]{
            \ifnum#1=1\num{0.392+-0.014} % Caso con fluttuazioni
            \else\num{0.377+-0.013}\fi   % Caso senza fluttuazioni
        }
        %\endgroup

        \foreach \c/\sl/\a/\p/\xS/\yS/\yLS in {
            1/2/south west/south west/2em/0em/-.5em,
            2/2/south west/south west/2em/0em/-.5em
        }{
            \def\betaExact{\betaExactList\c}
            \def\betaAprx{\betaAprxList\c}

            \foreach \r/\fl in {3/0,4/1}{
                \writeParetoParameters[\a]\fl\sl{Row\r\space c\c r1.\p}\xS\yS\yLS
            }
        } % «\fl=first line», «\sl=second line», «\a=anchor»
          % «\xS=\xShift», «\yS=\yShift» e «\yLS=\yLineShift»

        \foreach \c/\a/\p/\xS/\yS/\yLS in {
            1/south west/south west/0em/2em/-.5em,
            2/south west/south west/0em/2em/-.5em
        }{
            \def\xiExact{\xiExactList\c}
            \def\xiAprx{\xiAprxList\c}

            \def\mIExact{\mIExactList\c}
            \def\mIAprx{\mIAprxList\c}

            \def\sIExact{\sIExactList\c}
            \def\sIAprx{\sIAprxList\c}

            \def\mIIExact{\mIIExactList\c}
            \def\mIIAprx{\mIIAprxList\c}

            \def\sIIExact{\sIIExactList\c}
            \def\sIIAprx{\sIIAprxList\c}

            \foreach \r/\t in {3/0,4/1}{
                \writeBiLognormalParameters[\a]\t{Row\r\space c\c r1.\p}\xS\yS\yLS
            } % «\fl=first line», «\sl=second line», «\a=anchor»
              % «\xS=\xShift», «\yS=\yShift» e «\yLS=\yLineShift»
        }
    %\endgroup

    \redefineTikZbounds{Row1 c1r1}{Row5 c2r1}
\end{tikzpicture}%

% \input{../../Figure/.fnl/Capitolo4/ConfrontoConESenzaFluttuazioni/Parametri.txt}
